% =====================================================================
% ODD 2010 model description — MalariaSentinel vector ABM (mal_abm_fast)
%
% Structure follows the official ODD 2010 template (Grimm et al. 2010,
% Ecological Modelling 221:2760-2768), seven numbered elements.
% The shaded "Question" boxes of the template have been removed, as the
% official instructions require for a submitted description.
%
% Source of truth for this description: the repository source tree read
% on 2026-10-28 (HEAD 5e28644 plus working-tree modifications present on
% that date). Parameter values: mal-core/config/model_constants.yaml,
% schema version 2. State contract 4.0; env-NC contract 3.0.
% =====================================================================
\documentclass[11pt,a4paper]{article}

\usepackage[utf8]{inputenc}
\usepackage{lmodern}
% Latin Modern: a scalable Type 1 font. With T1 and the default
% Computer Modern, TeX asks for the EC fonts, and when the Type 1 EC
% ones (cm-super) are missing it falls back to bitmaps -- and
% microtype's font expansion cannot work with a bitmap font, so the
% build fails on a TeX Live without cm-super.
\usepackage[T1]{fontenc}
\usepackage[english]{babel}
\usepackage{geometry}
\geometry{margin=2.3cm}
\usepackage{xcolor}
\usepackage{booktabs}
\usepackage{longtable}
\usepackage{array}
\usepackage{amsmath,amssymb}
\usepackage{listings}
\usepackage{fancyhdr}
\usepackage{microtype}
\usepackage{etoolbox}

% The generated parameter table carries a superscript minus in a unit
% annotation (psu^-1), which inputenc does not define by itself.
\DeclareUnicodeCharacter{207B}{\ensuremath{-}}

% -- generated values --------------------------------------------------------
% Every model constant quoted in this document is emitted by the generator
% from mal-core/config/model_constants.yaml. A key that no longer exists
% makes \mskey fail the build.
% Generated by tools/gen_concepts_values.py -- do not edit.
% Values come from mal-core/config/model_constants.yaml via the
% validated loader; a key that no longer exists makes \mskey fail
% the build with an error, in any document that references it.

\makeatletter
\newcommand{\mskey}[1]{%
  \expandafter\ifx\csname ms@k@#1\endcsname\relax
    \errmessage{Concept document: unknown model constant key `#1'}%
  \else
    \csname ms@k@#1\endcsname%
  \fi}
\makeatother

\expandafter\def\csname ms@k@adults.daily_mort_basal\endcsname{0.87}
\expandafter\def\csname ms@k@adults.mort_cap\endcsname{0.93}
\expandafter\def\csname ms@k@adults.mort_floor\endcsname{0.6}
\expandafter\def\csname ms@k@adults.opt_c\endcsname{25.0}
\expandafter\def\csname ms@k@adults.sigma_c\endcsname{15.0}
\expandafter\def\csname ms@k@aquatic_stages.egg_daily_mort\endcsname{0.08}
\expandafter\def\csname ms@k@aquatic_stages.larva_bh_alpha\endcsname{0.05}
\expandafter\def\csname ms@k@aquatic_stages.larva_bh_s0\endcsname{0.95}
\expandafter\def\csname ms@k@aquatic_stages.last_larval_instar\endcsname{4}
\expandafter\def\csname ms@k@aquatic_stages.pupa_daily_mort\endcsname{0.05}
\expandafter\def\csname ms@k@aquatic_stages.salinity_missing_psu\endcsname{0.0}
\expandafter\def\csname ms@k@aquatic_stress.desicc_base_daily\endcsname{0.1}
\expandafter\def\csname ms@k@aquatic_stress.desicc_grace_days\endcsname{5.0}
\expandafter\def\csname ms@k@aquatic_stress.desicc_ramp_decay\endcsname{0.3}
\expandafter\def\csname ms@k@aquatic_stress.vulnerable_instar_max\endcsname{2}
\expandafter\def\csname ms@k@aquatic_stress.wash_fraction_max\endcsname{0.6}
\expandafter\def\csname ms@k@aquatic_stress.wash_threshold_mm\endcsname{40.0}
\expandafter\def\csname ms@k@capacity.persistent_availability\endcsname{1.0}
\expandafter\def\csname ms@k@contracts.cogs_state\endcsname{"4.0"}
\expandafter\def\csname ms@k@contracts.env_nc\endcsname{"3.0"}
\expandafter\def\csname ms@k@contracts.habitat\endcsname{"4.0"}
\expandafter\def\csname ms@k@dispersal.flight_hour\endcsname{20}
\expandafter\def\csname ms@k@dispersal.night_end_hour\endcsname{6}
\expandafter\def\csname ms@k@dispersal.night_start_hour\endcsname{18}
\expandafter\def\csname ms@k@dispersal.snap_radius_m\endcsname{500.0}
\expandafter\def\csname ms@k@dispersal.wind_flight_hours\endcsname{4.0}
\expandafter\def\csname ms@k@dispersal.wind_migration_prob\endcsname{0.05}
\expandafter\def\csname ms@k@dispersal.wind_self_speed_ms\endcsname{1.0}
\expandafter\def\csname ms@k@dispersal.wind_survival\endcsname{0.85}
\expandafter\def\csname ms@k@eip.base_c\endcsname{16.0}
\expandafter\def\csname ms@k@eip.threshold_gd\endcsname{110.0}
\expandafter\def\csname ms@k@geodesy.earth_radius_m\endcsname{6371008.8}
\expandafter\def\csname ms@k@geodesy.metres_per_deg_lat\endcsname{111195.08023353292}
\expandafter\def\csname ms@k@gonotrophic.cycle_duration_days\endcsname{2.0}
\expandafter\def\csname ms@k@gonotrophic.egg_batch_max\endcsname{170}
\expandafter\def\csname ms@k@gonotrophic.egg_batch_mean\endcsname{120}
\expandafter\def\csname ms@k@gonotrophic.egg_batch_min\endcsname{30}
\expandafter\def\csname ms@k@gonotrophic.egg_batch_n\endcsname{240}
\expandafter\def\csname ms@k@gonotrophic.feeding_success_rate\endcsname{0.825}
\expandafter\def\csname ms@k@gonotrophic.first_cycle_days\endcsname{4.0}
\expandafter\def\csname ms@k@gonotrophic.resting_duration_days\endcsname{0.5}
\expandafter\def\csname ms@k@habitat.activation_min_engine\endcsname{0.01}
\expandafter\def\csname ms@k@habitat.candidate_min_engine\endcsname{0.01}
\expandafter\def\csname ms@k@habitat_types.dynamics.crop_pluvial\endcsname{"pluvial"}
\expandafter\def\csname ms@k@habitat_types.dynamics.open_pluvial\endcsname{"pluvial"}
\expandafter\def\csname ms@k@habitat_types.dynamics.surface_persistent\endcsname{"persistent"}
\expandafter\def\csname ms@k@habitat_types.dynamics.urban_persistent\endcsname{"persistent"}
\expandafter\def\csname ms@k@habitat_types.dynamics.wetland_persistent\endcsname{"persistent"}
\expandafter\def\csname ms@k@habitat_types.dynamics.wetland_pluvial\endcsname{"pluvial"}
\expandafter\def\csname ms@k@habitat_types.types\endcsname{"surface\_persistent, wetland\_persistent, urban\_persistent, wetland\_pluvial, open\_pluvial, crop\_pluvial"}
\expandafter\def\csname ms@k@host_seeking.indoor_attraction_add\endcsname{0.72}
\expandafter\def\csname ms@k@host_seeking.indoor_fraction_rural\endcsname{0.3}
\expandafter\def\csname ms@k@host_seeking.indoor_fraction_urban\endcsname{0.72}
\expandafter\def\csname ms@k@host_seeking.starvation_days\endcsname{7}
\expandafter\def\csname ms@k@host_seeking.urbanicity_add\endcsname{0.2}
\expandafter\def\csname ms@k@ingest.candidate_min_ingest\endcsname{1e-06}
\expandafter\def\csname ms@k@ingest.cell_size_hint_m\endcsname{1000.0}
\expandafter\def\csname ms@k@ingest.coastline_buffer_m\endcsname{2000.0}
\expandafter\def\csname ms@k@ingest.coastline_resolution\endcsname{"f"}
\expandafter\def\csname ms@k@ingest.fill_frac_valid_max\endcsname{0.5}
\expandafter\def\csname ms@k@ingest.grid_degrees_step_max\endcsname{1}
\expandafter\def\csname ms@k@ingest.land_mask_threshold\endcsname{0.5}
\expandafter\def\csname ms@k@ingest.min_dilation_radius_cells\endcsname{1}
\expandafter\def\csname ms@k@ingest.nc_compression_level\endcsname{4}
\expandafter\def\csname ms@k@ingest.nodata_sentinel\endcsname{-9999.0}
\expandafter\def\csname ms@k@ingest.nodata_sentinel_int\endcsname{-9999}
\expandafter\def\csname ms@k@ingest.ocean_mask_threshold\endcsname{0.5}
\expandafter\def\csname ms@k@ingest.residential_from_indoor_frac\endcsname{0.8}
\expandafter\def\csname ms@k@ingest.sea_water_floor\endcsname{0.05}
\expandafter\def\csname ms@k@ingest.valid_count_epsilon\endcsname{1e-09}
\expandafter\def\csname ms@k@integration.directed_flight_budget_h\endcsname{0.4856}
\expandafter\def\csname ms@k@integration.dt_olfactory_s\endcsname{0.25}
\expandafter\def\csname ms@k@integration.dt_terminal_s\endcsname{0.04}
\expandafter\def\csname ms@k@integration.landing_speed_hold_s\endcsname{3.0}
\expandafter\def\csname ms@k@integration.landing_speed_mm_s\endcsname{50.0}
\expandafter\def\csname ms@k@integration.phase_duration_hours\endcsname{"12.0, 4.0, 6.0, 2.0"}
\expandafter\def\csname ms@k@integration.plume_memory_s\endcsname{0.7}
\expandafter\def\csname ms@k@mating.female_attendance_rate\endcsname{0.1033}
\expandafter\def\csname ms@k@mating.local_swarm_male_offer_factor\endcsname{0.0}
\expandafter\def\csname ms@k@mating.male_females_per_night_max\endcsname{3}
\expandafter\def\csname ms@k@mating.male_move_capacity_m\endcsname{50.0}
\expandafter\def\csname ms@k@mating.male_swarm_reference_size\endcsname{100.0}
\expandafter\def\csname ms@k@mating.p_mate_male_fraction\endcsname{0.072}
\expandafter\def\csname ms@k@mating.p_mate_pairs_per_individual\endcsname{0.167}
\expandafter\def\csname ms@k@mating.p_mate_size_exponent\endcsname{0.0}
\expandafter\def\csname ms@k@mating.rain_suppresses_swarming\endcsname{true}
\expandafter\def\csname ms@k@mating.swarm_duration_min\endcsname{28.3}
\expandafter\def\csname ms@k@mating.swarm_height_m\endcsname{2.87}
\expandafter\def\csname ms@k@mating.swarm_markers_per_km2\endcsname{1.0}
\expandafter\def\csname ms@k@mating.swarm_markers_source\endcsname{"engineering prior"}
\expandafter\def\csname ms@k@mating.swarm_start_min_after_sunset\endcsname{5.0}
\expandafter\def\csname ms@k@oviposition.detect_range_m\endcsname{2000.0}
\expandafter\def\csname ms@k@oviposition.min_attraction\endcsname{0.3}
\expandafter\def\csname ms@k@oviposition.scale_m\endcsname{50.0}
\expandafter\def\csname ms@k@oviposition.step_m\endcsname{50.0}
\expandafter\def\csname ms@k@oviposition.timeout_days\endcsname{7}
\expandafter\def\csname ms@k@pluvial.cold_start_z\endcsname{0.0}
\expandafter\def\csname ms@k@pluvial.delta_days\endcsname{1.0}
\expandafter\def\csname ms@k@pluvial.dry_memory_days\endcsname{5.0}
\expandafter\def\csname ms@k@pluvial.rain_scale_mm\endcsname{2.5}
\expandafter\def\csname ms@k@rng.aquatic_seed_larva_mortality\endcsname{123}
\expandafter\def\csname ms@k@rng.aquatic_seed_stage_mortality\endcsname{42}
\expandafter\def\csname ms@k@rng.aquatic_seed_washout\endcsname{789}
\expandafter\def\csname ms@k@species.arabiensis.activity_weights\endcsname{"0.02, 0.1, 0.8, 0.08"}
\expandafter\def\csname ms@k@species.arabiensis.aestivation_supported\endcsname{true}
\expandafter\def\csname ms@k@species.arabiensis.anemotaxis_deviation_deg\endcsname{30.0}
\expandafter\def\csname ms@k@species.arabiensis.anemotaxis_rule\endcsname{"upwind"}
\expandafter\def\csname ms@k@species.arabiensis.arrested_daily_survival\endcsname{0.986}
\expandafter\def\csname ms@k@species.arabiensis.blood_meal_refractory_days\endcsname{0.5}
\expandafter\def\csname ms@k@species.arabiensis.blood_meal_size_ul\endcsname{0.3}
\expandafter\def\csname ms@k@species.arabiensis.blood_meal_target_ul\endcsname{0.5}
\expandafter\def\csname ms@k@species.arabiensis.casting_angle_deg\endcsname{90.0}
\expandafter\def\csname ms@k@species.arabiensis.casting_probability\endcsname{0.1}
\expandafter\def\csname ms@k@species.arabiensis.co2_activation_radius_m\endcsname{100.0}
\expandafter\def\csname ms@k@species.arabiensis.daily_move_capacity_m\endcsname{590.0}
\expandafter\def\csname ms@k@species.arabiensis.directed_flight_budget_h\endcsname{0.4856}
\expandafter\def\csname ms@k@species.arabiensis.eggs_per_ul\endcsname{223.0}
\expandafter\def\csname ms@k@species.arabiensis.exploration_range_m\endcsname{100.0}
\expandafter\def\csname ms@k@species.arabiensis.exploration_step_correlation\endcsname{0.5}
\expandafter\def\csname ms@k@species.arabiensis.exploration_step_m\endcsname{50.0}
\expandafter\def\csname ms@k@species.arabiensis.first_cycle_target_ul\endcsname{0.8}
\expandafter\def\csname ms@k@species.arabiensis.flight_speed_ms\endcsname{0.3375}
\expandafter\def\csname ms@k@species.arabiensis.gonotrophic_arrest_prob_dry\endcsname{0.4454}
\expandafter\def\csname ms@k@species.arabiensis.male_daily_survival\endcsname{0.773}
\expandafter\def\csname ms@k@species.arabiensis.male_ratio\endcsname{0.5}
\expandafter\def\csname ms@k@species.arabiensis.meal_size_min_ul\endcsname{0.1}
\expandafter\def\csname ms@k@species.arabiensis.olfactory_radius_m\endcsname{50.0}
\expandafter\def\csname ms@k@species.arabiensis.plume_memory_s\endcsname{0.7}
\expandafter\def\csname ms@k@species.arabiensis.pref.cattle\endcsname{0.2}
\expandafter\def\csname ms@k@species.arabiensis.pref.goat\endcsname{0.15}
\expandafter\def\csname ms@k@species.arabiensis.pref.human\endcsname{0.5}
\expandafter\def\csname ms@k@species.arabiensis.pref.sheep\endcsname{0.1}
\expandafter\def\csname ms@k@species.arabiensis.pref.wildlife\endcsname{0.05}
\expandafter\def\csname ms@k@species.arabiensis.probing_success_rate\endcsname{0.7275}
\expandafter\def\csname ms@k@species.arabiensis.salinity_hi_tol_ppt\endcsname{15.0}
\expandafter\def\csname ms@k@species.arabiensis.salinity_opt_ppt\endcsname{0.0}
\expandafter\def\csname ms@k@species.arabiensis.salinity_slope\endcsname{0.15}
\expandafter\def\csname ms@k@species.arabiensis.terminal_radius_m\endcsname{0.15}
\expandafter\def\csname ms@k@species.arabiensis.visual_radius_m\endcsname{0.485}
\expandafter\def\csname ms@k@species.arabiensis.weights.crop_pluvial\endcsname{0.1}
\expandafter\def\csname ms@k@species.arabiensis.weights.open_pluvial\endcsname{0.55}
\expandafter\def\csname ms@k@species.arabiensis.weights.surface_persistent\endcsname{0.1}
\expandafter\def\csname ms@k@species.arabiensis.weights.urban_persistent\endcsname{0.15}
\expandafter\def\csname ms@k@species.arabiensis.weights.wetland_persistent\endcsname{0.1}
\expandafter\def\csname ms@k@species.arabiensis.weights.wetland_pluvial\endcsname{0.0}
\expandafter\def\csname ms@k@species.arabiensis.wind_block_threshold_ms\endcsname{0.3375}
\expandafter\def\csname ms@k@species.arabiensis.wind_reference\endcsname{"u100/v100"}
\expandafter\def\csname ms@k@species.coluzzii.activity_weights\endcsname{"0.0, 0.255, 0.48, 0.265"}
\expandafter\def\csname ms@k@species.coluzzii.aestivation_supported\endcsname{true}
\expandafter\def\csname ms@k@species.coluzzii.anemotaxis_deviation_deg\endcsname{30.0}
\expandafter\def\csname ms@k@species.coluzzii.anemotaxis_rule\endcsname{"upwind"}
\expandafter\def\csname ms@k@species.coluzzii.arrested_daily_survival\endcsname{0.986}
\expandafter\def\csname ms@k@species.coluzzii.blood_meal_refractory_days\endcsname{0.5}
\expandafter\def\csname ms@k@species.coluzzii.blood_meal_size_ul\endcsname{0.3}
\expandafter\def\csname ms@k@species.coluzzii.blood_meal_target_ul\endcsname{0.5}
\expandafter\def\csname ms@k@species.coluzzii.casting_angle_deg\endcsname{90.0}
\expandafter\def\csname ms@k@species.coluzzii.casting_probability\endcsname{0.1}
\expandafter\def\csname ms@k@species.coluzzii.co2_activation_radius_m\endcsname{100.0}
\expandafter\def\csname ms@k@species.coluzzii.daily_move_capacity_m\endcsname{590.0}
\expandafter\def\csname ms@k@species.coluzzii.directed_flight_budget_h\endcsname{0.4856}
\expandafter\def\csname ms@k@species.coluzzii.eggs_per_ul\endcsname{223.0}
\expandafter\def\csname ms@k@species.coluzzii.exploration_range_m\endcsname{100.0}
\expandafter\def\csname ms@k@species.coluzzii.exploration_step_correlation\endcsname{0.5}
\expandafter\def\csname ms@k@species.coluzzii.exploration_step_m\endcsname{50.0}
\expandafter\def\csname ms@k@species.coluzzii.first_cycle_target_ul\endcsname{0.8}
\expandafter\def\csname ms@k@species.coluzzii.flight_speed_ms\endcsname{0.3375}
\expandafter\def\csname ms@k@species.coluzzii.gonotrophic_arrest_prob_dry\endcsname{0.2704}
\expandafter\def\csname ms@k@species.coluzzii.male_daily_survival\endcsname{0.773}
\expandafter\def\csname ms@k@species.coluzzii.male_ratio\endcsname{0.5}
\expandafter\def\csname ms@k@species.coluzzii.meal_size_min_ul\endcsname{0.1}
\expandafter\def\csname ms@k@species.coluzzii.olfactory_radius_m\endcsname{50.0}
\expandafter\def\csname ms@k@species.coluzzii.plume_memory_s\endcsname{0.7}
\expandafter\def\csname ms@k@species.coluzzii.pref.cattle\endcsname{0.005}
\expandafter\def\csname ms@k@species.coluzzii.pref.goat\endcsname{0.003}
\expandafter\def\csname ms@k@species.coluzzii.pref.human\endcsname{0.99}
\expandafter\def\csname ms@k@species.coluzzii.pref.sheep\endcsname{0.001}
\expandafter\def\csname ms@k@species.coluzzii.pref.wildlife\endcsname{0.001}
\expandafter\def\csname ms@k@species.coluzzii.probing_success_rate\endcsname{0.7275}
\expandafter\def\csname ms@k@species.coluzzii.salinity_hi_tol_ppt\endcsname{30.0}
\expandafter\def\csname ms@k@species.coluzzii.salinity_opt_ppt\endcsname{0.0}
\expandafter\def\csname ms@k@species.coluzzii.salinity_slope\endcsname{0.08}
\expandafter\def\csname ms@k@species.coluzzii.terminal_radius_m\endcsname{0.15}
\expandafter\def\csname ms@k@species.coluzzii.visual_radius_m\endcsname{0.485}
\expandafter\def\csname ms@k@species.coluzzii.weights.crop_pluvial\endcsname{0.05}
\expandafter\def\csname ms@k@species.coluzzii.weights.open_pluvial\endcsname{0.25}
\expandafter\def\csname ms@k@species.coluzzii.weights.surface_persistent\endcsname{0.4}
\expandafter\def\csname ms@k@species.coluzzii.weights.urban_persistent\endcsname{0.2}
\expandafter\def\csname ms@k@species.coluzzii.weights.wetland_persistent\endcsname{0.1}
\expandafter\def\csname ms@k@species.coluzzii.weights.wetland_pluvial\endcsname{0.0}
\expandafter\def\csname ms@k@species.coluzzii.wind_block_threshold_ms\endcsname{0.3375}
\expandafter\def\csname ms@k@species.coluzzii.wind_reference\endcsname{"u100/v100"}
\expandafter\def\csname ms@k@species.default_species\endcsname{"coluzzii"}
\expandafter\def\csname ms@k@species.funestus.activity_weights\endcsname{"0.02, 0.1, 0.8, 0.08"}
\expandafter\def\csname ms@k@species.funestus.aestivation_supported\endcsname{false}
\expandafter\def\csname ms@k@species.funestus.anemotaxis_deviation_deg\endcsname{30.0}
\expandafter\def\csname ms@k@species.funestus.anemotaxis_rule\endcsname{"upwind"}
\expandafter\def\csname ms@k@species.funestus.arrested_daily_survival\endcsname{0.986}
\expandafter\def\csname ms@k@species.funestus.blood_meal_refractory_days\endcsname{0.5}
\expandafter\def\csname ms@k@species.funestus.blood_meal_size_ul\endcsname{0.3}
\expandafter\def\csname ms@k@species.funestus.blood_meal_target_ul\endcsname{0.5}
\expandafter\def\csname ms@k@species.funestus.casting_angle_deg\endcsname{90.0}
\expandafter\def\csname ms@k@species.funestus.casting_probability\endcsname{0.1}
\expandafter\def\csname ms@k@species.funestus.co2_activation_radius_m\endcsname{100.0}
\expandafter\def\csname ms@k@species.funestus.daily_move_capacity_m\endcsname{590.0}
\expandafter\def\csname ms@k@species.funestus.directed_flight_budget_h\endcsname{0.4856}
\expandafter\def\csname ms@k@species.funestus.eggs_per_ul\endcsname{223.0}
\expandafter\def\csname ms@k@species.funestus.exploration_range_m\endcsname{100.0}
\expandafter\def\csname ms@k@species.funestus.exploration_step_correlation\endcsname{0.5}
\expandafter\def\csname ms@k@species.funestus.exploration_step_m\endcsname{50.0}
\expandafter\def\csname ms@k@species.funestus.first_cycle_target_ul\endcsname{0.8}
\expandafter\def\csname ms@k@species.funestus.flight_speed_ms\endcsname{0.3375}
\expandafter\def\csname ms@k@species.funestus.gonotrophic_arrest_prob_dry\endcsname{0.0}
\expandafter\def\csname ms@k@species.funestus.male_daily_survival\endcsname{0.773}
\expandafter\def\csname ms@k@species.funestus.male_ratio\endcsname{0.5}
\expandafter\def\csname ms@k@species.funestus.meal_size_min_ul\endcsname{0.1}
\expandafter\def\csname ms@k@species.funestus.olfactory_radius_m\endcsname{50.0}
\expandafter\def\csname ms@k@species.funestus.plume_memory_s\endcsname{0.7}
\expandafter\def\csname ms@k@species.funestus.pref.cattle\endcsname{0.005}
\expandafter\def\csname ms@k@species.funestus.pref.goat\endcsname{0.003}
\expandafter\def\csname ms@k@species.funestus.pref.human\endcsname{0.99}
\expandafter\def\csname ms@k@species.funestus.pref.sheep\endcsname{0.001}
\expandafter\def\csname ms@k@species.funestus.pref.wildlife\endcsname{0.001}
\expandafter\def\csname ms@k@species.funestus.probing_success_rate\endcsname{0.7275}
\expandafter\def\csname ms@k@species.funestus.salinity_hi_tol_ppt\endcsname{15.0}
\expandafter\def\csname ms@k@species.funestus.salinity_opt_ppt\endcsname{0.0}
\expandafter\def\csname ms@k@species.funestus.salinity_slope\endcsname{0.15}
\expandafter\def\csname ms@k@species.funestus.terminal_radius_m\endcsname{0.15}
\expandafter\def\csname ms@k@species.funestus.visual_radius_m\endcsname{0.485}
\expandafter\def\csname ms@k@species.funestus.weights.crop_pluvial\endcsname{0.05}
\expandafter\def\csname ms@k@species.funestus.weights.open_pluvial\endcsname{0.1}
\expandafter\def\csname ms@k@species.funestus.weights.surface_persistent\endcsname{0.45}
\expandafter\def\csname ms@k@species.funestus.weights.urban_persistent\endcsname{0.1}
\expandafter\def\csname ms@k@species.funestus.weights.wetland_persistent\endcsname{0.3}
\expandafter\def\csname ms@k@species.funestus.weights.wetland_pluvial\endcsname{0.0}
\expandafter\def\csname ms@k@species.funestus.wind_block_threshold_ms\endcsname{0.3375}
\expandafter\def\csname ms@k@species.funestus.wind_reference\endcsname{"u100/v100"}
\expandafter\def\csname ms@k@species.gambiae_ss.activity_weights\endcsname{"0.02, 0.1, 0.8, 0.08"}
\expandafter\def\csname ms@k@species.gambiae_ss.aestivation_supported\endcsname{true}
\expandafter\def\csname ms@k@species.gambiae_ss.anemotaxis_deviation_deg\endcsname{30.0}
\expandafter\def\csname ms@k@species.gambiae_ss.anemotaxis_rule\endcsname{"upwind"}
\expandafter\def\csname ms@k@species.gambiae_ss.arrested_daily_survival\endcsname{0.986}
\expandafter\def\csname ms@k@species.gambiae_ss.blood_meal_refractory_days\endcsname{0.5}
\expandafter\def\csname ms@k@species.gambiae_ss.blood_meal_size_ul\endcsname{0.3}
\expandafter\def\csname ms@k@species.gambiae_ss.blood_meal_target_ul\endcsname{0.5}
\expandafter\def\csname ms@k@species.gambiae_ss.casting_angle_deg\endcsname{90.0}
\expandafter\def\csname ms@k@species.gambiae_ss.casting_probability\endcsname{0.1}
\expandafter\def\csname ms@k@species.gambiae_ss.co2_activation_radius_m\endcsname{100.0}
\expandafter\def\csname ms@k@species.gambiae_ss.daily_move_capacity_m\endcsname{590.0}
\expandafter\def\csname ms@k@species.gambiae_ss.directed_flight_budget_h\endcsname{0.4856}
\expandafter\def\csname ms@k@species.gambiae_ss.eggs_per_ul\endcsname{223.0}
\expandafter\def\csname ms@k@species.gambiae_ss.exploration_range_m\endcsname{100.0}
\expandafter\def\csname ms@k@species.gambiae_ss.exploration_step_correlation\endcsname{0.5}
\expandafter\def\csname ms@k@species.gambiae_ss.exploration_step_m\endcsname{50.0}
\expandafter\def\csname ms@k@species.gambiae_ss.first_cycle_target_ul\endcsname{0.8}
\expandafter\def\csname ms@k@species.gambiae_ss.flight_speed_ms\endcsname{0.3375}
\expandafter\def\csname ms@k@species.gambiae_ss.gonotrophic_arrest_prob_dry\endcsname{0.6838}
\expandafter\def\csname ms@k@species.gambiae_ss.male_daily_survival\endcsname{0.773}
\expandafter\def\csname ms@k@species.gambiae_ss.male_ratio\endcsname{0.5}
\expandafter\def\csname ms@k@species.gambiae_ss.meal_size_min_ul\endcsname{0.1}
\expandafter\def\csname ms@k@species.gambiae_ss.olfactory_radius_m\endcsname{50.0}
\expandafter\def\csname ms@k@species.gambiae_ss.plume_memory_s\endcsname{0.7}
\expandafter\def\csname ms@k@species.gambiae_ss.pref.cattle\endcsname{0.005}
\expandafter\def\csname ms@k@species.gambiae_ss.pref.goat\endcsname{0.003}
\expandafter\def\csname ms@k@species.gambiae_ss.pref.human\endcsname{0.99}
\expandafter\def\csname ms@k@species.gambiae_ss.pref.sheep\endcsname{0.001}
\expandafter\def\csname ms@k@species.gambiae_ss.pref.wildlife\endcsname{0.001}
\expandafter\def\csname ms@k@species.gambiae_ss.probing_success_rate\endcsname{0.7275}
\expandafter\def\csname ms@k@species.gambiae_ss.salinity_hi_tol_ppt\endcsname{20.0}
\expandafter\def\csname ms@k@species.gambiae_ss.salinity_opt_ppt\endcsname{0.0}
\expandafter\def\csname ms@k@species.gambiae_ss.salinity_slope\endcsname{0.12}
\expandafter\def\csname ms@k@species.gambiae_ss.terminal_radius_m\endcsname{0.15}
\expandafter\def\csname ms@k@species.gambiae_ss.visual_radius_m\endcsname{0.485}
\expandafter\def\csname ms@k@species.gambiae_ss.weights.crop_pluvial\endcsname{0.1}
\expandafter\def\csname ms@k@species.gambiae_ss.weights.open_pluvial\endcsname{0.65}
\expandafter\def\csname ms@k@species.gambiae_ss.weights.surface_persistent\endcsname{0.05}
\expandafter\def\csname ms@k@species.gambiae_ss.weights.urban_persistent\endcsname{0.15}
\expandafter\def\csname ms@k@species.gambiae_ss.weights.wetland_persistent\endcsname{0.05}
\expandafter\def\csname ms@k@species.gambiae_ss.weights.wetland_pluvial\endcsname{0.0}
\expandafter\def\csname ms@k@species.gambiae_ss.wind_block_threshold_ms\endcsname{0.3375}
\expandafter\def\csname ms@k@species.gambiae_ss.wind_reference\endcsname{"u100/v100"}
\expandafter\def\csname ms@k@species.k_ref\endcsname{50.0}
\expandafter\def\csname ms@k@species.melas.activity_weights\endcsname{"0.02, 0.1, 0.8, 0.08"}
\expandafter\def\csname ms@k@species.melas.aestivation_supported\endcsname{false}
\expandafter\def\csname ms@k@species.melas.anemotaxis_deviation_deg\endcsname{30.0}
\expandafter\def\csname ms@k@species.melas.anemotaxis_rule\endcsname{"upwind"}
\expandafter\def\csname ms@k@species.melas.arrested_daily_survival\endcsname{0.986}
\expandafter\def\csname ms@k@species.melas.blood_meal_refractory_days\endcsname{0.5}
\expandafter\def\csname ms@k@species.melas.blood_meal_size_ul\endcsname{0.3}
\expandafter\def\csname ms@k@species.melas.blood_meal_target_ul\endcsname{0.5}
\expandafter\def\csname ms@k@species.melas.casting_angle_deg\endcsname{90.0}
\expandafter\def\csname ms@k@species.melas.casting_probability\endcsname{0.1}
\expandafter\def\csname ms@k@species.melas.co2_activation_radius_m\endcsname{100.0}
\expandafter\def\csname ms@k@species.melas.daily_move_capacity_m\endcsname{590.0}
\expandafter\def\csname ms@k@species.melas.directed_flight_budget_h\endcsname{0.4856}
\expandafter\def\csname ms@k@species.melas.eggs_per_ul\endcsname{223.0}
\expandafter\def\csname ms@k@species.melas.exploration_range_m\endcsname{100.0}
\expandafter\def\csname ms@k@species.melas.exploration_step_correlation\endcsname{0.5}
\expandafter\def\csname ms@k@species.melas.exploration_step_m\endcsname{50.0}
\expandafter\def\csname ms@k@species.melas.first_cycle_target_ul\endcsname{0.8}
\expandafter\def\csname ms@k@species.melas.flight_speed_ms\endcsname{0.3375}
\expandafter\def\csname ms@k@species.melas.gonotrophic_arrest_prob_dry\endcsname{0.0}
\expandafter\def\csname ms@k@species.melas.male_daily_survival\endcsname{0.773}
\expandafter\def\csname ms@k@species.melas.male_ratio\endcsname{0.5}
\expandafter\def\csname ms@k@species.melas.meal_size_min_ul\endcsname{0.1}
\expandafter\def\csname ms@k@species.melas.olfactory_radius_m\endcsname{50.0}
\expandafter\def\csname ms@k@species.melas.plume_memory_s\endcsname{0.7}
\expandafter\def\csname ms@k@species.melas.pref.cattle\endcsname{0.005}
\expandafter\def\csname ms@k@species.melas.pref.goat\endcsname{0.003}
\expandafter\def\csname ms@k@species.melas.pref.human\endcsname{0.99}
\expandafter\def\csname ms@k@species.melas.pref.sheep\endcsname{0.001}
\expandafter\def\csname ms@k@species.melas.pref.wildlife\endcsname{0.001}
\expandafter\def\csname ms@k@species.melas.probing_success_rate\endcsname{0.7275}
\expandafter\def\csname ms@k@species.melas.salinity_hi_tol_ppt\endcsname{38.0}
\expandafter\def\csname ms@k@species.melas.salinity_opt_ppt\endcsname{0.0}
\expandafter\def\csname ms@k@species.melas.salinity_slope\endcsname{0.02}
\expandafter\def\csname ms@k@species.melas.terminal_radius_m\endcsname{0.15}
\expandafter\def\csname ms@k@species.melas.visual_radius_m\endcsname{0.485}
\expandafter\def\csname ms@k@species.melas.weights.crop_pluvial\endcsname{0.05}
\expandafter\def\csname ms@k@species.melas.weights.open_pluvial\endcsname{0.4}
\expandafter\def\csname ms@k@species.melas.weights.surface_persistent\endcsname{0.25}
\expandafter\def\csname ms@k@species.melas.weights.urban_persistent\endcsname{0.15}
\expandafter\def\csname ms@k@species.melas.weights.wetland_persistent\endcsname{0.15}
\expandafter\def\csname ms@k@species.melas.weights.wetland_pluvial\endcsname{0.0}
\expandafter\def\csname ms@k@species.melas.wind_block_threshold_ms\endcsname{0.3375}
\expandafter\def\csname ms@k@species.melas.wind_reference\endcsname{"u100/v100"}
\expandafter\def\csname ms@k@species.neutral.weights.crop_pluvial\endcsname{0.05}
\expandafter\def\csname ms@k@species.neutral.weights.open_pluvial\endcsname{0.4}
\expandafter\def\csname ms@k@species.neutral.weights.surface_persistent\endcsname{0.25}
\expandafter\def\csname ms@k@species.neutral.weights.urban_persistent\endcsname{0.15}
\expandafter\def\csname ms@k@species.neutral.weights.wetland_persistent\endcsname{0.15}
\expandafter\def\csname ms@k@species.neutral.weights.wetland_pluvial\endcsname{0.0}
\expandafter\def\csname ms@k@species.stephensi.activity_weights\endcsname{"0.02, 0.1, 0.8, 0.08"}
\expandafter\def\csname ms@k@species.stephensi.aestivation_supported\endcsname{false}
\expandafter\def\csname ms@k@species.stephensi.anemotaxis_deviation_deg\endcsname{30.0}
\expandafter\def\csname ms@k@species.stephensi.anemotaxis_rule\endcsname{"upwind"}
\expandafter\def\csname ms@k@species.stephensi.arrested_daily_survival\endcsname{0.986}
\expandafter\def\csname ms@k@species.stephensi.blood_meal_refractory_days\endcsname{0.5}
\expandafter\def\csname ms@k@species.stephensi.blood_meal_size_ul\endcsname{0.3}
\expandafter\def\csname ms@k@species.stephensi.blood_meal_target_ul\endcsname{0.5}
\expandafter\def\csname ms@k@species.stephensi.casting_angle_deg\endcsname{90.0}
\expandafter\def\csname ms@k@species.stephensi.casting_probability\endcsname{0.1}
\expandafter\def\csname ms@k@species.stephensi.co2_activation_radius_m\endcsname{100.0}
\expandafter\def\csname ms@k@species.stephensi.daily_move_capacity_m\endcsname{590.0}
\expandafter\def\csname ms@k@species.stephensi.directed_flight_budget_h\endcsname{0.4856}
\expandafter\def\csname ms@k@species.stephensi.eggs_per_ul\endcsname{223.0}
\expandafter\def\csname ms@k@species.stephensi.exploration_range_m\endcsname{100.0}
\expandafter\def\csname ms@k@species.stephensi.exploration_step_correlation\endcsname{0.5}
\expandafter\def\csname ms@k@species.stephensi.exploration_step_m\endcsname{50.0}
\expandafter\def\csname ms@k@species.stephensi.first_cycle_target_ul\endcsname{0.8}
\expandafter\def\csname ms@k@species.stephensi.flight_speed_ms\endcsname{0.3375}
\expandafter\def\csname ms@k@species.stephensi.gonotrophic_arrest_prob_dry\endcsname{0.0}
\expandafter\def\csname ms@k@species.stephensi.male_daily_survival\endcsname{0.773}
\expandafter\def\csname ms@k@species.stephensi.male_ratio\endcsname{0.5}
\expandafter\def\csname ms@k@species.stephensi.meal_size_min_ul\endcsname{0.1}
\expandafter\def\csname ms@k@species.stephensi.olfactory_radius_m\endcsname{50.0}
\expandafter\def\csname ms@k@species.stephensi.plume_memory_s\endcsname{0.7}
\expandafter\def\csname ms@k@species.stephensi.pref.cattle\endcsname{0.15}
\expandafter\def\csname ms@k@species.stephensi.pref.goat\endcsname{0.05}
\expandafter\def\csname ms@k@species.stephensi.pref.human\endcsname{0.7}
\expandafter\def\csname ms@k@species.stephensi.pref.sheep\endcsname{0.05}
\expandafter\def\csname ms@k@species.stephensi.pref.wildlife\endcsname{0.05}
\expandafter\def\csname ms@k@species.stephensi.probing_success_rate\endcsname{0.7275}
\expandafter\def\csname ms@k@species.stephensi.salinity_hi_tol_ppt\endcsname{20.0}
\expandafter\def\csname ms@k@species.stephensi.salinity_opt_ppt\endcsname{0.0}
\expandafter\def\csname ms@k@species.stephensi.salinity_slope\endcsname{0.15}
\expandafter\def\csname ms@k@species.stephensi.terminal_radius_m\endcsname{0.15}
\expandafter\def\csname ms@k@species.stephensi.visual_radius_m\endcsname{0.485}
\expandafter\def\csname ms@k@species.stephensi.weights.crop_pluvial\endcsname{0.05}
\expandafter\def\csname ms@k@species.stephensi.weights.open_pluvial\endcsname{0.4}
\expandafter\def\csname ms@k@species.stephensi.weights.surface_persistent\endcsname{0.25}
\expandafter\def\csname ms@k@species.stephensi.weights.urban_persistent\endcsname{0.15}
\expandafter\def\csname ms@k@species.stephensi.weights.wetland_persistent\endcsname{0.15}
\expandafter\def\csname ms@k@species.stephensi.weights.wetland_pluvial\endcsname{0.0}
\expandafter\def\csname ms@k@species.stephensi.wind_block_threshold_ms\endcsname{0.3375}
\expandafter\def\csname ms@k@species.stephensi.wind_reference\endcsname{"u100/v100"}
\expandafter\def\csname ms@k@thermal.efd_max_eggs_per_day\endcsname{0.5}
\expandafter\def\csname ms@k@thermal.efd_tmax_c\endcsname{34.0}
\expandafter\def\csname ms@k@thermal.efd_tmin_c\endcsname{16.0}
\expandafter\def\csname ms@k@thermal.efd_topt_c\endcsname{25.0}
\expandafter\def\csname ms@k@thermal.egg_development_scale\endcsname{0.00062}
\expandafter\def\csname ms@k@thermal.egg_development_tmax_c\endcsname{40.0}
\expandafter\def\csname ms@k@thermal.egg_development_tmin_c\endcsname{10.0}
\expandafter\def\csname ms@k@thermal.larva_development_scale\endcsname{0.001}
\expandafter\def\csname ms@k@thermal.larva_development_tmax_c\endcsname{40.0}
\expandafter\def\csname ms@k@thermal.larva_development_tmin_c\endcsname{10.0}
\expandafter\def\csname ms@k@thermal.pupa_development_scale\endcsname{0.00066}
\expandafter\def\csname ms@k@thermal.pupa_development_tmax_c\endcsname{40.0}
\expandafter\def\csname ms@k@thermal.pupa_development_tmin_c\endcsname{12.0}
\expandafter\def\csname ms@k@transmission.beta_hv\endcsname{0.4}
\expandafter\def\csname ms@k@transmission.beta_vh\endcsname{0.5}
\expandafter\def\csname ms@k@transmission.human_immunity_days\endcsname{180}
\expandafter\def\csname ms@k@transmission.human_incubation_days\endcsname{12}
\expandafter\def\csname ms@k@transmission.human_infectious_days\endcsname{20}
\expandafter\def\csname ms@k@transmission.initial_human_prevalence\endcsname{0.05}
\expandafter\def\csname ms@k@transmission.initial_vector_infected_frac\endcsname{0.0}
\expandafter\def\csname ms@k@urban.beta_urban\endcsname{0.75}
\expandafter\def\csname ms@k@urban.rain_pools_built_damage\endcsname{0.3}
\expandafter\def\csname ms@k@version\endcsname{3}
\expandafter\def\csname ms@k@wildlife.class_suitability.bare\endcsname{0.2}
\expandafter\def\csname ms@k@wildlife.class_suitability.built\endcsname{0.05}
\expandafter\def\csname ms@k@wildlife.class_suitability.crops\endcsname{0.3}
\expandafter\def\csname ms@k@wildlife.class_suitability.flooded_vegetation\endcsname{0.45}
\expandafter\def\csname ms@k@wildlife.class_suitability.grass\endcsname{0.6}
\expandafter\def\csname ms@k@wildlife.class_suitability.shrub_and_scrub\endcsname{0.7}
\expandafter\def\csname ms@k@wildlife.class_suitability.snow_and_ice\endcsname{0.1}
\expandafter\def\csname ms@k@wildlife.class_suitability.trees\endcsname{0.9}
\expandafter\def\csname ms@k@wildlife.class_suitability.water\endcsname{0.0}
\expandafter\def\csname ms@k@wildlife.weights.habitat\endcsname{0.5}
\expandafter\def\csname ms@k@wildlife.weights.remoteness\endcsname{0.2}
\expandafter\def\csname ms@k@wildlife.weights.water\endcsname{0.3}

% The values that live in the code (module constants, function defaults) are
% emitted by the same generator, under the second macro family. An unknown
% key makes \codeconst fail the build.
% Generated by tools/gen_concepts_values.py -- do not edit.
% Values are read from the CODE that implements each stage, by AST.
% A key that no longer exists makes \codeconst fail the LaTeX build
% with an error, in any document that references it.
% Each line carries the file:line the value was read from.

\makeatletter
\newcommand{\codeconst}[1]{%
  \expandafter\ifx\csname cc@k@#1\endcsname\relax
    \errmessage{Concept document: unknown code constant `#1'}%
  \else
    \csname cc@k@#1\endcsname%
  \fi}
\makeatother

% download.wildlife_habitat_fallback = 0.5  (mal-commonlib/src/mal_commonlib/data/loaders/wildlife.py:205)
\expandafter\def\csname cc@k@download.wildlife_habitat_fallback\endcsname{0.5}
% download.wildlife_remoteness_neutral = 1.0  (mal-commonlib/src/mal_commonlib/data/loaders/wildlife.py:161)
\expandafter\def\csname cc@k@download.wildlife_remoteness_neutral\endcsname{1.0}
% ingest.bbox_east = 1.5  (mal-commonlib/src/mal_commonlib/aois.yaml:18)
\expandafter\def\csname cc@k@ingest.bbox_east\endcsname{1.5}
% ingest.bbox_north = 11.5  (mal-commonlib/src/mal_commonlib/aois.yaml:18)
\expandafter\def\csname cc@k@ingest.bbox_north\endcsname{11.5}
% ingest.bbox_south = 4.5  (mal-commonlib/src/mal_commonlib/aois.yaml:18)
\expandafter\def\csname cc@k@ingest.bbox_south\endcsname{4.5}
% ingest.bbox_west = -3.5  (mal-commonlib/src/mal_commonlib/aois.yaml:18)
\expandafter\def\csname cc@k@ingest.bbox_west\endcsname{-3.5}
% ingest.cell_size_m = 1000  (mal-commonlib/src/mal_commonlib/aois.yaml:20)
\expandafter\def\csname cc@k@ingest.cell_size_m\endcsname{1000}
% ingest.env_daily_channels = 2  (mal-core/src/mal_core/ingest/daily_nc.py:414)
\expandafter\def\csname cc@k@ingest.env_daily_channels\endcsname{2}
% ingest.env_static_planes = 5  (mal-core/src/mal_core/ingest/daily_nc.py:414)
\expandafter\def\csname cc@k@ingest.env_static_planes\endcsname{5}
% ingest.fallback_building_rural = 0.1  (mal-core/src/mal_core/ingest/hosts.py:142)
\expandafter\def\csname cc@k@ingest.fallback_building_rural\endcsname{0.1}
% ingest.fallback_building_urban = 0.6  (mal-core/src/mal_core/ingest/hosts.py:142)
\expandafter\def\csname cc@k@ingest.fallback_building_urban\endcsname{0.6}
% ingest.fallback_wildlife_proxy = 0.3  (mal-core/src/mal_core/ingest/hosts.py:165)
\expandafter\def\csname cc@k@ingest.fallback_wildlife_proxy\endcsname{0.3}
% ingest.jrc_occurrence_threshold_pct = 80  (mal-commonlib/src/mal_commonlib/data/loaders/jrc_gsw.py:72)
\expandafter\def\csname cc@k@ingest.jrc_occurrence_threshold_pct\endcsname{80}
% ingest.k_ref_diag_per_km2 = 1329900.0  (mal-core/src/mal_core/habitat/capacity.py:30)
\expandafter\def\csname cc@k@ingest.k_ref_diag_per_km2\endcsname{1329900.0}
% ingest.mobility_beta_day = 0.05  (mal-core/src/mal_core/ingest/mobility.py:43)
\expandafter\def\csname cc@k@ingest.mobility_beta_day\endcsname{0.05}
% ingest.mobility_beta_livestock = 0.1  (mal-core/src/mal_core/ingest/mobility.py:45)
\expandafter\def\csname cc@k@ingest.mobility_beta_livestock\endcsname{0.1}
% ingest.mobility_beta_night = 0.5  (mal-core/src/mal_core/ingest/mobility.py:44)
\expandafter\def\csname cc@k@ingest.mobility_beta_night\endcsname{0.5}
% ingest.mobility_cell_size_km = 1.0  (mal-core/src/mal_core/ingest/mobility.py:42)
\expandafter\def\csname cc@k@ingest.mobility_cell_size_km\endcsname{1.0}
% ingest.mobility_max_distance_km = 50.0  (mal-core/src/mal_core/ingest/mobility.py:46)
\expandafter\def\csname cc@k@ingest.mobility_max_distance_km\endcsname{50.0}
% ingest.mobility_sparsity_threshold = 0.0005  (mal-core/src/mal_core/ingest/mobility.py:47)
\expandafter\def\csname cc@k@ingest.mobility_sparsity_threshold\endcsname{0.0005}
% ingest.urban_class_rural = 50  (mal-core/src/mal_core/ingest/hosts.py:126)
\expandafter\def\csname cc@k@ingest.urban_class_rural\endcsname{50}
% ingest.urban_class_urban = 30  (mal-core/src/mal_core/ingest/hosts.py:125)
\expandafter\def\csname cc@k@ingest.urban_class_urban\endcsname{30}
% ingest.urban_class_water = 20  (mal-core/src/mal_core/ingest/hosts.py:127)
\expandafter\def\csname cc@k@ingest.urban_class_water\endcsname{20}
% prediction.block_size = 128  (mal-core/src/mal_core/prediction/predictor.py:126)
\expandafter\def\csname cc@k@prediction.block_size\endcsname{128}
% prediction.contract_version_manifest = "1.1"  (mal-core/src/mal_core/prediction/registry.py:31)
\expandafter\def\csname cc@k@prediction.contract_version_manifest\endcsname{"1.1"}
% prediction.contract_version_sidecar = "1.1"  (mal-core/src/mal_core/prediction/predictor.py:146)
\expandafter\def\csname cc@k@prediction.contract_version_sidecar\endcsname{"1.1"}
% prediction.in_channels = 10  (mal-core/src/mal_core/prediction/registry.py:32)
\expandafter\def\csname cc@k@prediction.in_channels\endcsname{10}
% prediction.manifest_version = "1.0"  (mal-core/src/mal_core/prediction/registry.py:30)
\expandafter\def\csname cc@k@prediction.manifest_version\endcsname{"1.0"}
% prediction.out_channels = 6  (mal-core/src/mal_core/prediction/registry.py:33)
\expandafter\def\csname cc@k@prediction.out_channels\endcsname{6}
% prediction.scale_continental_m = 5000  (mal-core/src/mal_core/prediction/aggregator.py:29)
\expandafter\def\csname cc@k@prediction.scale_continental_m\endcsname{5000}
% prediction.scale_national_m = 2000  (mal-core/src/mal_core/prediction/aggregator.py:28)
\expandafter\def\csname cc@k@prediction.scale_national_m\endcsname{2000}
% prediction.scale_regional_m = 1000  (mal-core/src/mal_core/prediction/aggregator.py:27)
\expandafter\def\csname cc@k@prediction.scale_regional_m\endcsname{1000}
% scoring.d11_larva_band_max = 0.9  (mal-core/src/mal_core/scoring/scorers/d11_aquatic_structure.py:35)
\expandafter\def\csname cc@k@scoring.d11_larva_band_max\endcsname{0.9}
% scoring.d11_larva_band_min = 0.65  (mal-core/src/mal_core/scoring/scorers/d11_aquatic_structure.py:35)
\expandafter\def\csname cc@k@scoring.d11_larva_band_min\endcsname{0.65}
% scoring.d11_pupa_band_max = 0.15  (mal-core/src/mal_core/scoring/scorers/d11_aquatic_structure.py:36)
\expandafter\def\csname cc@k@scoring.d11_pupa_band_max\endcsname{0.15}
% scoring.d11_pupa_band_min = 0.02  (mal-core/src/mal_core/scoring/scorers/d11_aquatic_structure.py:36)
\expandafter\def\csname cc@k@scoring.d11_pupa_band_min\endcsname{0.02}
% scoring.d15_expansion_fraction = 0.2  (mal-core/src/mal_core/scoring/scorers/d15_persistence.py:125)
\expandafter\def\csname cc@k@scoring.d15_expansion_fraction\endcsname{0.2}
% scoring.d15_min_horizon_days = 30  (mal-core/src/mal_core/scoring/scorers/d15_persistence.py:97)
\expandafter\def\csname cc@k@scoring.d15_min_horizon_days\endcsname{30}
% scoring.d15_survival_fraction = 0.3  (mal-core/src/mal_core/scoring/scorers/d15_persistence.py:105)
\expandafter\def\csname cc@k@scoring.d15_survival_fraction\endcsname{0.3}
% scoring.d16_coverage_target = 0.5  (mal-core/src/mal_core/scoring/scorers/d16_detection_coverage.py:44)
\expandafter\def\csname cc@k@scoring.d16_coverage_target\endcsname{0.5}
% scoring.d16_radius_km_default = 5.0  (mal-core/src/mal_core/scoring/scorers/d16_detection_coverage.py:43)
\expandafter\def\csname cc@k@scoring.d16_radius_km_default\endcsname{5.0}
% scoring.d16_sensitivity_radius_far_km = 10.0  (mal-core/src/mal_core/scoring/scorers/d16_detection_coverage.py:45)
\expandafter\def\csname cc@k@scoring.d16_sensitivity_radius_far_km\endcsname{10.0}
% scoring.d16_sensitivity_radius_mid_km = 5.0  (mal-core/src/mal_core/scoring/scorers/d16_detection_coverage.py:45)
\expandafter\def\csname cc@k@scoring.d16_sensitivity_radius_mid_km\endcsname{5.0}
% scoring.d16_sensitivity_radius_near_km = 2.0  (mal-core/src/mal_core/scoring/scorers/d16_detection_coverage.py:45)
\expandafter\def\csname cc@k@scoring.d16_sensitivity_radius_near_km\endcsname{2.0}
% scoring.d19_activation_floor = 0.05  (mal-core/src/mal_core/scoring/scorers/d19_pool_persistence.py:81)
\expandafter\def\csname cc@k@scoring.d19_activation_floor\endcsname{0.05}
% scoring.d19_check_drying_weight = 0.2  (mal-core/src/mal_core/scoring/scorers/d19_pool_persistence.py:130)
\expandafter\def\csname cc@k@scoring.d19_check_drying_weight\endcsname{0.2}
% scoring.d19_check_floor_weight = 0.3  (mal-core/src/mal_core/scoring/scorers/d19_pool_persistence.py:134)
\expandafter\def\csname cc@k@scoring.d19_check_floor_weight\endcsname{0.3}
% scoring.d19_check_monotone_weight = 0.2  (mal-core/src/mal_core/scoring/scorers/d19_pool_persistence.py:132)
\expandafter\def\csname cc@k@scoring.d19_check_monotone_weight\endcsname{0.2}
% scoring.d19_check_persistence_weight = 0.3  (mal-core/src/mal_core/scoring/scorers/d19_pool_persistence.py:128)
\expandafter\def\csname cc@k@scoring.d19_check_persistence_weight\endcsname{0.3}
% scoring.d19_dry_days = 10  (mal-core/src/mal_core/scoring/scorers/d19_pool_persistence.py:84)
\expandafter\def\csname cc@k@scoring.d19_dry_days\endcsname{10}
% scoring.d19_pulse_days = 3  (mal-core/src/mal_core/scoring/scorers/d19_pool_persistence.py:78)
\expandafter\def\csname cc@k@scoring.d19_pulse_days\endcsname{3}
% scoring.d19_rain_pulse_mm = 30.0  (mal-core/src/mal_core/scoring/scorers/d19_pool_persistence.py:79)
\expandafter\def\csname cc@k@scoring.d19_rain_pulse_mm\endcsname{30.0}
% scoring.d1_band_max = 0.7  (mal-core/src/mal_core/scoring/scorers/d1_habitat_occupancy.py:33)
\expandafter\def\csname cc@k@scoring.d1_band_max\endcsname{0.7}
% scoring.d1_band_min = 0.2  (mal-core/src/mal_core/scoring/scorers/d1_habitat_occupancy.py:33)
\expandafter\def\csname cc@k@scoring.d1_band_min\endcsname{0.2}
% scoring.d1_center = 0.45  (mal-core/src/mal_core/scoring/scorers/d1_habitat_occupancy.py:34)
\expandafter\def\csname cc@k@scoring.d1_center\endcsname{0.45}
% scoring.d1_sigma = 0.15  (mal-core/src/mal_core/scoring/scorers/d1_habitat_occupancy.py:35)
\expandafter\def\csname cc@k@scoring.d1_sigma\endcsname{0.15}
% scoring.d20_cap_probe_mm = 80.0  (mal-core/src/mal_core/scoring/scorers/d20_washout_response.py:56)
\expandafter\def\csname cc@k@scoring.d20_cap_probe_mm\endcsname{80.0}
% scoring.d20_check_cap_weight = 0.2  (mal-core/src/mal_core/scoring/scorers/d20_washout_response.py:71)
\expandafter\def\csname cc@k@scoring.d20_check_cap_weight\endcsname{0.2}
% scoring.d20_check_no_washout_weight = 0.3  (mal-core/src/mal_core/scoring/scorers/d20_washout_response.py:69)
\expandafter\def\csname cc@k@scoring.d20_check_no_washout_weight\endcsname{0.3}
% scoring.d20_check_threshold_weight = 0.2  (mal-core/src/mal_core/scoring/scorers/d20_washout_response.py:73)
\expandafter\def\csname cc@k@scoring.d20_check_threshold_weight\endcsname{0.2}
% scoring.d20_check_washout_weight = 0.3  (mal-core/src/mal_core/scoring/scorers/d20_washout_response.py:67)
\expandafter\def\csname cc@k@scoring.d20_check_washout_weight\endcsname{0.3}
% scoring.d20_min_washout_fraction = 0.15  (mal-core/src/mal_core/scoring/scorers/d20_washout_response.py:48)
\expandafter\def\csname cc@k@scoring.d20_min_washout_fraction\endcsname{0.15}
% scoring.d20_no_washout_probe_mm = 30.0  (mal-core/src/mal_core/scoring/scorers/d20_washout_response.py:51)
\expandafter\def\csname cc@k@scoring.d20_no_washout_probe_mm\endcsname{30.0}
% scoring.d20_threshold_max_fraction = 0.01  (mal-core/src/mal_core/scoring/scorers/d20_washout_response.py:63)
\expandafter\def\csname cc@k@scoring.d20_threshold_max_fraction\endcsname{0.01}
% scoring.d20_threshold_probe_mm = 40.0  (mal-core/src/mal_core/scoring/scorers/d20_washout_response.py:61)
\expandafter\def\csname cc@k@scoring.d20_threshold_probe_mm\endcsname{40.0}
% scoring.d20_washout_probe_mm = 50.0  (mal-core/src/mal_core/scoring/scorers/d20_washout_response.py:46)
\expandafter\def\csname cc@k@scoring.d20_washout_probe_mm\endcsname{50.0}
% scoring.d21_horizon_day = 30  (mal-core/src/mal_core/scoring/scorers/d21_dispersal_kernel.py:112)
\expandafter\def\csname cc@k@scoring.d21_horizon_day\endcsname{30}
% scoring.d21_p50_band_max = 0.7  (mal-core/src/mal_core/scoring/scorers/d21_dispersal_kernel.py:39)
\expandafter\def\csname cc@k@scoring.d21_p50_band_max\endcsname{0.7}
% scoring.d21_p50_band_min = 0.15  (mal-core/src/mal_core/scoring/scorers/d21_dispersal_kernel.py:39)
\expandafter\def\csname cc@k@scoring.d21_p50_band_min\endcsname{0.15}
% scoring.d21_p50_center = 0.4  (mal-core/src/mal_core/scoring/scorers/d21_dispersal_kernel.py:41)
\expandafter\def\csname cc@k@scoring.d21_p50_center\endcsname{0.4}
% scoring.d21_p50_sigma = 0.2  (mal-core/src/mal_core/scoring/scorers/d21_dispersal_kernel.py:41)
\expandafter\def\csname cc@k@scoring.d21_p50_sigma\endcsname{0.2}
% scoring.d21_p90_band_max = 3.0  (mal-core/src/mal_core/scoring/scorers/d21_dispersal_kernel.py:40)
\expandafter\def\csname cc@k@scoring.d21_p90_band_max\endcsname{3.0}
% scoring.d21_p90_band_min = 0.7  (mal-core/src/mal_core/scoring/scorers/d21_dispersal_kernel.py:40)
\expandafter\def\csname cc@k@scoring.d21_p90_band_min\endcsname{0.7}
% scoring.d21_p90_center = 1.6  (mal-core/src/mal_core/scoring/scorers/d21_dispersal_kernel.py:42)
\expandafter\def\csname cc@k@scoring.d21_p90_center\endcsname{1.6}
% scoring.d21_p90_sigma = 0.8  (mal-core/src/mal_core/scoring/scorers/d21_dispersal_kernel.py:42)
\expandafter\def\csname cc@k@scoring.d21_p90_sigma\endcsname{0.8}
% scoring.d25_r_target = 0.6  (mal-core/src/mal_core/scoring/scorers/d25_cases_ghana.py:45)
\expandafter\def\csname cc@k@scoring.d25_r_target\endcsname{0.6}
% scoring.d2_lifetime_centre = 11.0  (mal-core/src/mal_core/scoring/scorers/d2_survival.py:70)
\expandafter\def\csname cc@k@scoring.d2_lifetime_centre\endcsname{11.0}
% scoring.d2_lifetime_width = 4.0  (mal-core/src/mal_core/scoring/scorers/d2_survival.py:70)
\expandafter\def\csname cc@k@scoring.d2_lifetime_width\endcsname{4.0}
% scoring.d2_window_max = 15  (mal-core/src/mal_core/scoring/scorers/d2_survival.py:47)
\expandafter\def\csname cc@k@scoring.d2_window_max\endcsname{15}
% scoring.d2_window_min = 8  (mal-core/src/mal_core/scoring/scorers/d2_survival.py:47)
\expandafter\def\csname cc@k@scoring.d2_window_min\endcsname{8}
% scoring.d3_band_max = 0.5  (mal-core/src/mal_core/scoring/scorers/d3_eip.py:69)
\expandafter\def\csname cc@k@scoring.d3_band_max\endcsname{0.5}
% scoring.d3_band_min = 0.2  (mal-core/src/mal_core/scoring/scorers/d3_eip.py:69)
\expandafter\def\csname cc@k@scoring.d3_band_min\endcsname{0.2}
% scoring.d3_center = 0.35  (mal-core/src/mal_core/scoring/scorers/d3_eip.py:72)
\expandafter\def\csname cc@k@scoring.d3_center\endcsname{0.35}
% scoring.d3_day = 30  (mal-core/src/mal_core/scoring/scorers/d3_eip.py:67)
\expandafter\def\csname cc@k@scoring.d3_day\endcsname{30}
% scoring.d3_sigma = 0.15  (mal-core/src/mal_core/scoring/scorers/d3_eip.py:72)
\expandafter\def\csname cc@k@scoring.d3_sigma\endcsname{0.15}
% scoring.d4_collapse_max = 0.8  (mal-core/src/mal_core/scoring/scorers/d4_stability.py:70)
\expandafter\def\csname cc@k@scoring.d4_collapse_max\endcsname{0.8}
% scoring.d4_collapse_min = 0.1  (mal-core/src/mal_core/scoring/scorers/d4_stability.py:70)
\expandafter\def\csname cc@k@scoring.d4_collapse_min\endcsname{0.1}
% scoring.d4_explosion_max = 5.0  (mal-core/src/mal_core/scoring/scorers/d4_stability.py:77)
\expandafter\def\csname cc@k@scoring.d4_explosion_max\endcsname{5.0}
% scoring.d5_band_max = 0.85  (mal-core/src/mal_core/scoring/scorers/d5_morans.py:95)
\expandafter\def\csname cc@k@scoring.d5_band_max\endcsname{0.85}
% scoring.d5_band_min = 0.2  (mal-core/src/mal_core/scoring/scorers/d5_morans.py:95)
\expandafter\def\csname cc@k@scoring.d5_band_min\endcsname{0.2}
% scoring.d5_center = 0.5  (mal-core/src/mal_core/scoring/scorers/d5_morans.py:98)
\expandafter\def\csname cc@k@scoring.d5_center\endcsname{0.5}
% scoring.d5_sigma = 0.25  (mal-core/src/mal_core/scoring/scorers/d5_morans.py:98)
\expandafter\def\csname cc@k@scoring.d5_sigma\endcsname{0.25}
% scoring.d8_min_score = 0.4  (mal-core/src/mal_core/scoring/scorers/d8_coupling.py:35)
\expandafter\def\csname cc@k@scoring.d8_min_score\endcsname{0.4}
% scoring.d8_pass_threshold = 0.3  (mal-core/src/mal_core/scoring/scorers/d8_coupling.py:76)
\expandafter\def\csname cc@k@scoring.d8_pass_threshold\endcsname{0.3}
% scoring.d8_score_divisor = 0.5  (mal-core/src/mal_core/scoring/scorers/d8_coupling.py:70)
\expandafter\def\csname cc@k@scoring.d8_score_divisor\endcsname{0.5}
% scoring.gate_census_reference = 27000000  (mal-core/src/mal_core/scoring/scorers/gates.py:191)
\expandafter\def\csname cc@k@scoring.gate_census_reference\endcsname{27000000}
% scoring.gate_ratio_band_max = 1.2  (mal-core/src/mal_core/scoring/scorers/gates.py:214)
\expandafter\def\csname cc@k@scoring.gate_ratio_band_max\endcsname{1.2}
% scoring.gate_ratio_band_min = 0.8  (mal-core/src/mal_core/scoring/scorers/gates.py:214)
\expandafter\def\csname cc@k@scoring.gate_ratio_band_min\endcsname{0.8}
% scoring.gate_row_sum_threshold = 0.99  (mal-core/src/mal_core/scoring/scorers/gates.py:271)
\expandafter\def\csname cc@k@scoring.gate_row_sum_threshold\endcsname{0.99}
% scoring.gate_row_sum_tolerance = 0.01  (mal-core/src/mal_core/scoring/scorers/gates.py:260)
\expandafter\def\csname cc@k@scoring.gate_row_sum_tolerance\endcsname{0.01}
% scoring.gate_rural_class = 50  (mal-core/src/mal_core/scoring/scorers/gates.py:317)
\expandafter\def\csname cc@k@scoring.gate_rural_class\endcsname{50}
% scoring.gate_sample_files = 48  (mal-core/src/mal_core/scoring/scorers/gates.py:59)
\expandafter\def\csname cc@k@scoring.gate_sample_files\endcsname{48}
% scoring.gate_urban_class = 30  (mal-core/src/mal_core/scoring/scorers/gates.py:316)
\expandafter\def\csname cc@k@scoring.gate_urban_class\endcsname{30}
% scoring.gate_urban_ratio_max = 0.7  (mal-core/src/mal_core/scoring/scorers/gates.py:336)
\expandafter\def\csname cc@k@scoring.gate_urban_ratio_max\endcsname{0.7}
% scoring.gate_wall_time_limit_s = 30.0  (mal-core/src/mal_core/scoring/scorers/gates.py:158)
\expandafter\def\csname cc@k@scoring.gate_wall_time_limit_s\endcsname{30.0}
% scoring.regression_threshold = -0.05  (mal-core/src/mal_core/scoring/feedback.py:51)
\expandafter\def\csname cc@k@scoring.regression_threshold\endcsname{-0.05}
% scoring.verdict_borderline = 0.5  (mal-core/src/mal_core/scoring/runner.py:24)
\expandafter\def\csname cc@k@scoring.verdict_borderline\endcsname{0.5}
% scoring.verdict_regressed = 0.3  (mal-core/src/mal_core/scoring/runner.py:24)
\expandafter\def\csname cc@k@scoring.verdict_regressed\endcsname{0.3}
% scoring.verdict_viable = 0.7  (mal-core/src/mal_core/scoring/runner.py:24)
\expandafter\def\csname cc@k@scoring.verdict_viable\endcsname{0.7}
% training.ckpt_period = 10  (mal-core/src/mal_core/training/trainer.py:102)
\expandafter\def\csname cc@k@training.ckpt_period\endcsname{10}
% training.dice_count_threshold = 0.5  (mal-core/src/mal_core/training/model.py:110)
\expandafter\def\csname cc@k@training.dice_count_threshold\endcsname{0.5}
% training.dice_eps = 1e-06  (mal-core/src/mal_core/training/model.py:97)
\expandafter\def\csname cc@k@training.dice_eps\endcsname{1e-06}
% training.dice_tau = 0.5  (mal-core/src/mal_core/training/model.py:87)
\expandafter\def\csname cc@k@training.dice_tau\endcsname{0.5}
% training.dice_thresh = 0.5  (mal-core/src/mal_core/training/model.py:87)
\expandafter\def\csname cc@k@training.dice_thresh\endcsname{0.5}
% training.dice_weight = 0.5  (mal-core/src/mal_core/training/model.py:105)
\expandafter\def\csname cc@k@training.dice_weight\endcsname{0.5}
% training.env_channels = 4  (mal-core/src/mal_core/training/dataset.py:22)
\expandafter\def\csname cc@k@training.env_channels\endcsname{4}
% training.fallback_in_channels = 6  (mal-core/src/mal_core/training/wrapper.py:62)
\expandafter\def\csname cc@k@training.fallback_in_channels\endcsname{6}
% training.fallback_out_channels = 2  (mal-core/src/mal_core/training/wrapper.py:66)
\expandafter\def\csname cc@k@training.fallback_out_channels\endcsname{2}
% training.pad_multiple = 16  (mal-core/src/mal_core/training/wrapper.py:100)
\expandafter\def\csname cc@k@training.pad_multiple\endcsname{16}
% training.patch_size = 128  (mal-core/src/mal_core/training/dataset.py:19)
\expandafter\def\csname cc@k@training.patch_size\endcsname{128}
% training.sub_seed = 42  (mal-core/src/mal_core/training/dataset.py:210)
\expandafter\def\csname cc@k@training.sub_seed\endcsname{42}
% training.transmission_channels = 4  (mal-core/src/mal_core/training/dataset.py:21)
\expandafter\def\csname cc@k@training.transmission_channels\endcsname{4}
% training.unet_channels = 32, 64, 128, 256  (mal-core/src/mal_core/training/model.py:46)
\expandafter\def\csname cc@k@training.unet_channels\endcsname{32, 64, 128, 256}
% training.vector_channels = 2  (mal-core/src/mal_core/training/dataset.py:20)
\expandafter\def\csname cc@k@training.vector_channels\endcsname{2}



\AtBeginEnvironment{longtable}{\footnotesize}
\usepackage[hidelinks]{hyperref}

% -- code / parameter names ------------------------------------------------
\newcommand{\code}[1]{\path{#1}}
\newcommand{\bcode}[1]{\code{#1}}

% -- pseudo-code style -----------------------------------------------------
\lstset{
  basicstyle=\ttfamily\footnotesize,
  columns=fullflexible,
  keepspaces=true,
  breaklines=true,
  frame=single,
  framesep=3pt,
  xleftmargin=2pt,
  xrightmargin=2pt,
  aboveskip=6pt,
  belowskip=6pt,
  showstringspaces=false
}

% -- headers ---------------------------------------------------------------
\pagestyle{fancy}
\fancyhf{}
\fancyhead[L]{\footnotesize ODD 2010 --- MalariaSentinel vector ABM}
\fancyhead[R]{\footnotesize \code{mal_abm_fast}}
\fancyfoot[C]{\thepage}
\renewcommand{\headrulewidth}{0.4pt}

\setlength{\parskip}{0.35em}
\setlength{\emergencystretch}{3em}
\renewcommand{\arraystretch}{1.14}
\newcolumntype{L}[1]{>{\raggedright\arraybackslash}p{#1}}

\title{ODD model description of the MalariaSentinel\\
Anopheles--malaria agent-based model\\[2pt]
\large (\code{mal_abm_fast}, state contract 4.0, env-NC contract 3.0)}
\author{David Fl\'orez-Mazuera\\MalariaSentinel project}
\date{28 October 2026}

\begin{document}
\maketitle

The model description follows the ODD (Overview, Design concepts,
Details) protocol for describing individual- and agent-based models
(Grimm et al.\ 2006, 2010).

\subsection*{Scope and provenance of this description}

This document describes the vector agent-based model \code{mal_abm_fast}
as implemented in the repository \emph{MalariaSentinel} on 28 October
2026 (source tree at \code{HEAD 5e28644}, including the working-tree
modifications present on that date). It describes the model, not its
command-line interface; the executable and its flags are documented in
the repository. All model constants are read at run time from
\code{mal-core/config/model_constants.yaml} (schema version 2), which is
the single source for both the Python and the C++ sides. No model
constant has a compiled default: a missing or unknown key aborts the run.
The state output contract is version 4.0 and the environmental NetCDF
input contract is version 3.0; the generator string is
\code{proxy_habitat_v1}. The engine identifies itself as version 0.1.0.

The seven numbered elements below are the ODD 2010 elements, with their
official titles. Appendix~A records implementation details that matter
for exact replication but are not part of the protocol. Appendix~B
records provenance, versioning, and testing status. Throughout, a
``parameter'' is a fixed model constant, while a ``state variable'' is
what would have to be saved to restart a run exactly.

\subsection*{Sources of truth and their precedence}

This description quotes three kinds of authority, and they do not carry
the same weight. \textbf{Values} --- every number that is a model
constant --- come from \code{mal-core/config/model_constants.yaml} and
reach this document through the generated macros and the generated
parameter table at the end of this document; where the text here and the constants file
disagree, the constants file wins and this text is the defect.
\textbf{Contract values and names} --- the version strings, the band
names, and the key names of the state sidecar and of the environmental
NetCDF --- come from the generated contract table, which is derived from
the code that writes them; where the text here and that table disagree,
the table wins and this text is the defect. \textbf{Semantics} --- what
a state variable means, which process a submodel implements, what a
formula stands for --- are owned by this document; where the text here
and the code disagree, this text is the defect, because the code is the
model and the description is a description of it. On every contradiction
the lower authority is the bug and the higher one is corrected. Each of
these rules is stated by the document that owns that content, so this
document states the rule for its own content and there is no shared page
that arbitrates between the documents.

\subsection*{Notation}

\begin{center}
\small
\begin{tabular}{@{}ll@{}}
\toprule
Symbol & Meaning \\
\midrule
$d$ & simulation day index, 0-based; one step is one day \\
$i$ or $(r,c)$ & grid cell (row $r$, column $c$) \\
$p$ & habitat patch; in the current implementation one patch per candidate cell \\
$g$ & habitat type; configured types and engine slots differ (Section~7.1) \\
$k$ & aquatic stage (egg, larva instar 1--4, pupa) \\
$T_{i,d}$ & daily mean air temperature at cell $i$ ($^\circ$C) \\
$R_{i,d}$ & daily rainfall at cell $i$ (mm\,d$^{-1}$) \\
$\sigma_{i,d}$ & water salinity at cell $i$ (psu; 0 = freshwater) \\
$H_{i,g}$ & static habitat potential of type $g$ in cell $i$, in $[0,1]$ \\
$w_g$ & species affinity weight for habitat type $g$ \\
$z_{i,d}$ & pluvial (rain-driven) habitat state, in $[0,1]$ \\
$S_{i,d}$ & daily habitat availability (``accuracy'' in the code), in $[0,1]$ \\
$K_{i,d}$ & daily larval capacity of cell $i$ (individuals) \\
$N_{p,k,d}$ & number of immatures in cohort $(p,k)$ \\
$p_d$ & adult daily survival probability \\
$E$ & extrinsic incubation progress (degree-days) \\
$H_j$ & number of humans in cell $j$ \\
\bottomrule
\end{tabular}
\end{center}

% =====================================================================
\section{Purpose}

The model is a daily, spatially explicit simulator of \emph{Anopheles}
population dynamics and malaria transmission over a regional
area of interest (AOI). It has three intended uses.

\begin{enumerate}
\item \textbf{Ground-truth generator for a surrogate model.} The engine
  produces the state fields (adult counts per cell, human infection
  fields) that train, validate, and stress-test a U-Net surrogate of the
  vector dynamics. The ABM is the reference process; the surrogate is an
  approximation to be fitted to it.
\item \textbf{Simulation engine of a calibration and scoring pipeline.}
  The engine is run under alternative climate, habitat, host, and
  seeding configurations; the outputs are scored against patterns and
  against the surrogate. The purpose is comparative and
  pattern-oriented, not the production of an official forecast.
\item \textbf{Mechanistic scenario analysis.} The model is used to ask
  where and when a vector population can establish, how the aquatic
  stages respond to rainfall and temperature, how host distribution and
  daily human mobility reshape biting risk, and whether a seeded human
  outbreak can start a transmission chain.
\end{enumerate}

The model expands an earlier thin-slice formulation into a richer
representation of the same system: individual adult mosquitoes with a
gonotrophic cycle, cohort-based aquatic stages with thermal development
and density-dependent mortality, static remote-sensing habitat proxies
with a rain-driven pluvial response, a spatial host landscape with
phase-specific human mobility, and a Ross--Macdonald-type SEIR--SEI
transmission layer. The expansion is deliberate ``from basic principles
to richer representation'': each added layer is a mechanism that the
thin slice either omitted or aggregated.

The model is \emph{not} intended for the following, and the current code
does not implement them: vector control (insecticide-treated nets,
indoor residual spraying, larviciding), drugs or vaccination, mosquito
genetics or insecticide resistance, evolution or learning, human
demography (births, deaths, migration of the SEIR population), water
balance or depth-resolved hydrology, or sub-daily mosquito flight.
Human movement is represented only as a change in the effective host
distribution for biting, not as movement of the infection compartments.
A single mosquito species is active per run (the default is
\emph{An.\ coluzzii}); parameter rows for six species exist, of which
two (\emph{An.\ melas}, \emph{An.\ stephensi}) are dormant.

The patterns that the project intends to use when judging whether the
model is realistic enough for these purposes are seasonal adult
abundance, the persistence and spatial location of adult foci, the
fraction of vectors that become infectious, and the timing and location
of a seeded human outbreak. These patterns are calibration targets, not
outputs that the model is forced to reproduce: the current parameter set
is a set of literature values and explicit engineering priors. The
project states that a formal calibration to Ghanaian field data has not
yet been completed. The code records no fit provenance for any
parameter, so it can neither confirm nor refute that statement; it is
reported here as a project claim, not as a model fact
(Appendix~B.4).

% =====================================================================
\section{Entities, state variables, and scales}

\subsection*{2.1 Entities}

\begin{description}
\item[Adult mosquitoes (individual agents).] The only true individuals in
  the model. Each adult has a unique identity, a position, a sex, a
  species identity, a transmission state, and (for females) a
  gonotrophic state. Adults are stored in a structure-of-arrays (SoA)
  layout; dead agents are removed by swapping the last live agent into
  the vacated slot. Population size is not fixed; it is bounded only by
  the optional operational \code{--max-population} guard, which aborts a
  run rather than imposing a biological ceiling.

\item[Aquatic immatures (cohorts, not individuals).] Eggs, larvae
  (instars L1--L4), and pupae are not individual agents. They are
  aggregated as one cohort per (patch, stage, instar) with an integer
  count and a development fraction. A cohort is a bookkeeping
  super-individual, not an entity that can be identified or tracked.

\item[Habitat patches (modeler-defined spatial units).] A patch is a
  candidate breeding cell. The patch set is the union of the features in
  the habitat GeoPackage (one point per candidate cell) and any grid cell
  whose static total potential clears the engine floor
  (\code{candidate_min_engine} = \mskey{habitat.candidate_min_engine}).
  Three counts must be kept apart. In the delivered Ghana 2024 product
  the GeoPackage holds 382,127 point features; 368,761 of those features
  have \code{habitat_potential} at or above the floor; and 370,061 grid
  cells have a static total potential at or above the floor. The
  operational patch set is the \emph{union} of all GeoPackage cells and
  all floor cells (\code{coordinator.cpp:156-174}), so it is at least
  382,127 cells; in the delivered
  product the floor cells are a subset of the GeoPackage cells and the
  union is exactly 382,127 cells. The three numbers are readings of the
  delivered artefacts, not model constants
  (Appendix~B.4).
  Each patch has a stable \code{patch_id}, a cell $(r,c)$, a static
  amplitude, and (diagnostic only) a dominant habitat type.

\item[Grid cells and the environment.] The AOI is a regular grid of
  1\,km$^2$ cells in EPSG:4326. Cells carry the daily rainfall, air
  temperature and (optional) salinity, the five static habitat planes,
  the pluvial state $z$, a dry-day counter, and the daily capacity $K$.
  The default Ghana AOI is 551 columns $\times$ 779 rows $=429{,}229$
  cells.

\item[Human hosts (compartments, not agents).] Humans are not
  individuals. Each cell carries susceptible, exposed, infectious, and
  recovered \emph{counts}, plus three discrete age-of-infection cohort
  queues (exposed, infectious, recovered). The human population is
  constant during a run. In the delivered Ghana product the host grid
  integrates to about 32.0 million persons over the bounding box (which
  includes neighbouring land and sea, not only Ghana).

\item[Host-landscape cells.] Each cell carries counts of humans, cattle,
  goats, sheep, pigs, and chickens, a wildlife proxy, a building
  fraction, an indoor fraction, a residential fraction, and a binary
  urbanicity. These are attributes of the environment, not agents.

\item[Collectives.] There are no emergent collectives: no herds, swarms,
  households, or social groups. ``Patch'', ``cohort'', ``grid cell'', and
  ``phase-specific host grid'' are modeler-defined aggregates.
\end{description}

\subsection*{2.2 State variables}

Table~\ref{tab:state} lists the state variables. The test applied is the
one recommended by the protocol: if the model were stopped and restarted,
which values would have to be saved? The answer is the union of the rows
marked ``state'', plus the random-number-generator states. Rows whose note
begins ``derived'' or ``forcing'' are not independent state: they are
computed each day from the state, the static inputs, or the day's
environmental forcing, and they are listed because the engine exposes them
and because their units matter for replication. (The engine does not
currently implement a checkpoint writer; the list is the conceptual
restart state.)

\begin{longtable}{@{}L{2.9cm}L{5.6cm}L{6.5cm}@{}}
\caption{State variables by entity type. Counts are integers unless
stated; ``persons'' and ``individuals'' are absolute counts, not
densities.}\label{tab:state}\\
\toprule
\textbf{Entity} & \textbf{State variable} & \textbf{Unit, range, and notes} \\
\midrule
\endfirsthead
\toprule
\textbf{Entity} & \textbf{State variable} & \textbf{Unit, range, and notes} \\
\midrule
\endhead
\bottomrule
\endlastfoot
Adult mosquito & \code{uid} & integer; stable for life \\
 & \code{patch_id} & integer; patch association, not necessarily the current cell \\
 & \code{row}, \code{col} & cell index; current position \\
 & \code{lon}, \code{lat} & decimal degrees, EPSG:4326; current position \\
 & \code{sex} & 0 = male, 1 = female; fixed at emergence \\
 & \code{stage} & always 1 (adult) for live agents; immatures are cohorts \\
 & \code{vector_state} & 0 = susceptible, 1 = exposed, 2 = infectious \\
 & \code{parasite_eip_progress} & degree-days $\ge 0$; meaningful for adult females \\
 & \code{gonotrophic_state} & enum; live subset is TENERAL, MATE\_SEEKING, HOST\_SEEKING, BLOOD\_FED, RESTING, EGG\_MATURING, GRAVID, OVIPOSITION\_SEEKING, OVIPOSITING \\
 & \code{gonotrophic_timer} & days in the current gonotrophic state \\
 & \code{gonotrophic_cycles} & completed cycles; distinguishes first from later cycles \\
 & \code{feeding_success} & daily flag: 1 if fed today, 0 otherwise; reset each day (not restart state) \\
 & \code{species_id} & 0--5; one active species per run \\
 & \code{last_patch_update_row/col} & cache used to skip the patch-snap scan; not biological \\
Aquatic cohort & \code{patch_id} & integer \\
 & \code{stage} & EGG, LARVA, PUPA \\
 & \code{larval_instar} & 1--4 for LARVA; 0 otherwise \\
 & \code{development} & fraction of the current stage completed, $[0,1]$ \\
 & \code{count} & individuals; integer $\ge 0$ \\
 & \code{mean_age_days} & days; diagnostic counter, incremented only on days when the cohort's patch is active, so it does not advance while the patch is frozen and is not a running age (\code{aquatic_cohort_bank.cpp:346-367}) \\
Habitat patch & \code{activated} & derived: $S\ge$ activation floor and salinity below the species' high tolerance \\
 & \code{habitat_avail} ($S$) & derived: $[0,1]$; species-neutral daily availability \\
 & \code{k_habitat} ($K$) & derived: individuals; absolute Beverton--Holt capacity \\
 & \code{habitat_dry_days} & days since $S$ last fell below the activation floor \\
 & \code{rain_d}, \code{temp_d}, \code{salinity_ppt} & forcing: mm\,d$^{-1}$, $^\circ$C, psu; copied from the cell \\
Grid cell & $z$ & state: $[0,1]$; pluvial state \\
 & $H_{i,g}$ & static input: \codeconst{ingest.env_static_planes} planes in $[0,1]$: JRC, built, wetland, open+water, crop \\
 & \code{dry_days} & state: integer; consecutive days with $S<$ activation floor \\
 & $K$ & derived: individuals; same value exposed through the patch state \\
 & rainfall, air temperature, salinity & forcing: mm\,d$^{-1}$, $^\circ$C, psu; daily \\
Human grid & $S_H,E_H,I_H,R_H$ & persons per cell; $R_H$ is zero when immunity is disabled \\
 & population & static input: persons per cell; constant \\
 & incidence & daily output: new $E_H$ per day, persons \\
 & prevalence & derived: $I_H/$population, dimensionless \\
 & exposed cohorts & persons by day since exposure; \mskey{transmission.human_incubation_days} bins \\
 & infectious cohorts & persons by day since onset; \mskey{transmission.human_infectious_days} bins \\
 & recovered cohorts & persons by day since recovery; \mskey{transmission.human_immunity_days} bins if immunity is on \\
Host cell & humans, cattle, goats, sheep, pigs, chickens & animals or persons per cell \\
 & wildlife proxy & index in $[0,1]$; 0 when no data \\
 & building fraction, indoor fraction, residential fraction & $[0,1]$ \\
 & urbanicity, urban class & 1 = urban / 0 = rural; GHS-SMOD class code \\
\end{longtable}

\subsection*{2.3 Temporal and spatial resolution and extent}

\textbf{Time.} One time step is one day. The engine accepts 1 to 731
daily steps in a single invocation; the delivered Ghana environmental
product covers 366 days (2024). Time is discrete: there is no continuous
clock and no event queue. Within a day the processes run in a fixed
order (Element~3). The four daily activity phases (DAY 06:00--18:00,
EVENING 18:00--21:00, NIGHT 21:00--05:00, DAWN 05:00--06:00) are labels
used to build phase-specific host landscapes during the host-seeking
pass; they are not sub-steps of an hourly clock. The only genuinely
sub-daily input is the optional 6-hourly wind field, which is queried at
one fixed hour per day (\mskey{dispersal.flight_hour}:00) when windborne migration is active.

\textbf{Space.} One grid cell represents 1\,km$^2$ on the ground, in
EPSG:4326. Widths are converted with the cosine of the AOI centroid
latitude; cell area for capacity is computed row by row as
$10^{-6}\,(\mskey{ingest.cell_size_hint_m})^2\cos(\phi_i)$\,km$^2$ (about 0.99\,km$^2$ at 8$^\circ$N).
The default Ghana AOI spans 5$^\circ$ of longitude by 7$^\circ$ of
latitude, giving 551 $\times$ 779 cells (a bounding box of about
428,600\,km$^2$, which includes sea and neighbouring land; it is not the
land area of Ghana). The model has no sub-cell geometry: a host, a
patch, and an adult are all located at a cell, with continuous
lon/lat carried for dispersal and wind displacement.

\textbf{Biological resolution.} Adults are individuals. Immatures are
cohorts. Humans are compartments. Hosts are counts. This mixture is
deliberate: the questions asked are about adult vector abundance and
spatial foci, for which individual adults matter, while the immature
stages are numerous and their individual identities are irrelevant to
the questions.

\textbf{Extent.} The model is regional, not global. The default AOI is
Ghana; any AOI with a matching manifest and habitat product can be run.
The environmental extent is the AOI bounding box; the biological extent
is the set of candidate patches within it. Adults can disperse outside
the candidate patch set but not outside the AOI grid: positions are
clamped to the grid at the cell-snapping step.

% =====================================================================
\section{Process overview and scheduling}

Time is modelled as a sequence of discrete daily updates. The order of
processes within a day is fixed and is the same in every run. The order
of entities within a process is deterministic given the random seed, but
it is \emph{not} a random permutation: adults are visited in SoA index
order, and removals use swap-with-last, which changes the index of the
agent moved into the vacated slot. Cohorts are visited in the order of
the cohort vector, which is the order of creation and is rebuilt after
compaction. Patches are emitted in row-major $(r,c)$ order.

The daily schedule is given below in pseudo-code. Names in
\textsc{small caps} are the processes that become the submodels of
Element~7.

\begin{lstlisting}
for d = 0 .. D-1:
    climate.set_day(d)                       # select the daily band
    UPDATE HABITAT                           # all cells, row-major
        z <- f(z, rain_d); S <- f(H, w, z); K <- k_ref * area * S
        dry_days <- update(S); patch_states <- build_patch_states()
    mosquito.advance_day(patch_states):
        ADVANCE AQUATIC STAGES               # cohort bank
            desiccation -> washout -> egg/pupa mortality
            -> larval density mortality -> development + promotion
        EMERGE ADULTS                        # pupae with dev >= 1
        bite_ledger.reset_day()
        for each adult female in SoA order:
            RUN GONOTROPHIC CYCLE            # one state transition
            if HOST_APPROACH: move toward host (if hosts loaded)
            if HOST_SEEKING: increment starvation timer; die at \mskey{host_seeking.starvation_days} d
            SEEK OVIPOSITION SITE AND OVIPOSIT
            SEEK HOSTS AND FEED              # non-phase-aware fallback
        if effective hosts loaded:           # phase-aware pass
            for phase in DAY, EVENING, NIGHT, DAWN:
                if species active in phase:
                    sample H_eff(day, phase) for humans and livestock
                    for each HOST_SEEKING female in SoA order:
                        move one step toward best host cell
                        SEEK HOSTS AND FEED
        DISPERSE ADULTS                      # wind first, then Gaussian
        snap every adult to nearest active patch within \mskey{dispersal.snap_radius_m} m
        SURVIVE ADULTS                       # daily survival draw
    if transmission enabled:
        aggregate adult counts; check and trigger human outbreak
        advance vector EIP                   # exposed females only
        advance human SEIR from today's bites
        record daily statistics
    date <- date + 1 day
\end{lstlisting}

Table~\ref{tab:sched} states, for each process, which entity performs it,
in what order the entities are visited, and whether updates are
asynchronous (a new value is visible immediately to later entities in the
same process) or synchronous (values are collected and applied
together).

\begin{longtable}{@{}L{4.0cm}L{3.0cm}L{4.4cm}L{3.6cm}@{}}
\caption{Process order within one simulated day.}\label{tab:sched}\\
\toprule
\textbf{Process} & \textbf{Performed by} & \textbf{Entity order} & \textbf{Update mode} \\
\midrule
\endfirsthead
\toprule
\textbf{Process} & \textbf{Performed by} & \textbf{Entity order} & \textbf{Update mode} \\
\midrule
\endhead
\bottomrule
\endlastfoot
Update habitat & coordinator, all cells & row-major $(r,c)$ & synchronous within a cell; cells independent \\
Advance aquatic stages & cohort bank & cohort-vector order & asynchronous: each cohort is updated in place \\
Emerge adults & pupae with dev $\ge1$ & cohort-vector order & synchronous event list, applied after the scan \\
Gonotrophic transition & adult females & SoA index order & asynchronous; a state change is visible in the same iteration \\
Oviposition seeking and laying & gravid females & SoA index order & asynchronous; eggs are added immediately \\
Host seeking and feeding & host-seeking females & SoA index order, then phase order & asynchronous; a fed female is skipped in later phases \\
Dispersal & all adults & SoA index order & asynchronous; positions updated in place \\
Patch snap & all adults & SoA index order & asynchronous; nearest active patch within \mskey{dispersal.snap_radius_m}\,m \\
Adult survival & all adults & reverse SoA index order & asynchronous; survivors are compacted by swap-with-last \\
Transmission & vector females and human cells & SoA index order; cells row-major & vector EIP then human SEIR; human cohorts shifted per cell \\
\end{longtable}

Several consequences of this schedule are worth stating explicitly
because they affect results.

\begin{enumerate}
\item \textbf{No within-day feedback from adults to habitat.} The habitat
  state $S$ and capacity $K$ are computed once, from the day's rainfall
  and the static planes, before any mosquito acts. Adults and immatures
  cannot change $S$ or $K$ during the day. The only feedback is indirect
  and delayed: adults lay eggs, and the resulting larvae experience next
  day's $K$.
\item \textbf{Aquatic promotions do not cascade.} Development is
  accumulated for every cohort first; promotions are then applied. A
  cohort that reaches the end of the egg stage today becomes an L1
  cohort and does not also advance through L1 on the same day.
\item \textbf{A female can oviposit and feed on the same day.} The
  gonotrophic transition runs first; a female that reaches
  OVIPOSITING lays eggs and is reset to HOST\_SEEKING; the host-seeking
  block then runs in the same iteration (or, in the phase-aware path, in
  the phase pass later in the same day). She can therefore take a blood
  meal on the day she lays.
\item \textbf{A female can make several feeding attempts in one day.} In
  the phase-aware path the pass re-examines every female that is still
  in HOST\_SEEKING in each active phase, and each phase makes one
  host-selection attempt (one uniform draw over the attraction field).
  With the four default phases active, a female that does not feed can
  make up to four attempts in a day. Once she reaches BLOOD\_FED she
  leaves HOST\_SEEKING and is skipped by the remaining phases. The
  single-phase fallback path makes one attempt per female per day.
\item \textbf{Death is applied immediately.} When a female starves or an
  adult fails its survival draw, the last live agent is swapped into its
  slot and the live range shrinks. This changes the order in which the
  remaining agents are visited, but the order remains deterministic for
  a given seed.
\item \textbf{Transmission is a separate block after the mosquito day.}
  The vector EIP and the human SEIR update therefore use the adult
  positions and the bite ledger \emph{after} dispersal, patch snapping,
  and adult mortality. A mosquito that died today does not contribute to
  today's transmission update, but the bites it recorded before dying
  remain in the ledger.
\end{enumerate}

The model is stochastic but reproducible: for a given source state, a
given environmental input, and a given master seed, the sequence of
draws is fixed. The seed structure is described in Element~5 and
Appendix~A.

% =====================================================================
\section{Design concepts}

The eleven design concepts are treated as a checklist. Two of them,
learning and prediction, are not implemented; they are declared here
rather than silently omitted, so that a reader does not have to infer
their absence. The model is rule-based: the adaptive traits are fixed
species traits and fixed behavioural rules, not optimized strategies.

\subsection{Basic principles}

The model is built on five principles.

\begin{enumerate}
\item \textbf{Individual-based ecology.} System-level dynamics (population
  growth, spatial spread, transmission) are expected to emerge from the
  fates of individual adults and from the environment they experience
  (Grimm \& Railsback 2005). Adults are the only individuals; the
  immature stages are aggregated because their individual identities are
  not needed for the questions asked.
\item \textbf{Metabolic and thermal physiology.} Development of eggs,
  larvae, and pupae, adult survival, egg production, and the extrinsic
  incubation period are temperature-dependent. The functional forms are
  taken from the mosquito thermal-biology literature: Brière-1 curves for
  eggs and pupae, a quadratic curve for larvae (Bayoh \& Lindsay 2003;
  Mordecai et al.\ 2013), a Gaussian survival curve around an optimum,
  and a degree-day EIP.
\item \textbf{Ross--Macdonald transmission.} The transmission layer
  decomposes vectorial capacity into bites on humans, the probability of
  infection per bite in each direction, vector survival, and the
  extrinsic incubation period (Macdonald 1957; Smith \& McKenzie 2004).
  It is spatialized: every quantity is a per-cell field, and the bite
  process is explicit rather than a mean biting rate.
\item \textbf{Open remote-sensing proxies instead of a hydrological
  model.} Breeding habitat is represented by static potential planes
  derived from open satellite products (JRC Global Surface Water and
  Dynamic World land cover), multiplied by a rain-driven pluvial state.
  There is no digital elevation model, no flow accumulation, no water
  balance in millimetres, and no depth or volume. This is a deliberate
  structural choice: the available open data support a presence and
  areal-fraction proxy, not a measured water budget.
\item \textbf{Pattern-oriented modelling as the validation strategy.} The
  model is judged by whether it reproduces a set of observed patterns at
  the individual and system levels (seasonal abundance, focus
  persistence, infectious fraction, outbreak timing), not by fitting a
  single aggregate curve. The patterns are listed in Element~1.
\end{enumerate}

The model uses prior theory (thermal curves, Ross--Macdonald) at the
submodel level and a system-level construction of its own (static
potential $\times$ pluvial availability, phase-specific host landscape,
bite-ledger transmission). It is not designed to test one of these
theories against another; it is a synthetic laboratory in which their
combined consequences can be examined.

\subsection{Emergence}

The results expected to emerge from individual behaviour and spatial
heterogeneity are: the seasonal trajectory of adult abundance; the
location, size, and persistence of adult foci; local extinction and
re colonization; the age and gonotrophic structure of the female
population; the fraction of vectors that become infectious; the
probability and timing of a human outbreak after seeding; and the
sensitivity of all of these to rainfall, temperature, habitat, and host
distribution.

Other results are tightly imposed by the rules and are therefore
``built in'' rather than emergent: the shape of the thermal development
and survival curves, the absolute capacity field $K$, the force of
infection given a bite field, the host preference ordering, and the
dispersal kernel. A change in those parameters changes the results
directly; it is not mediated by individual adaptation.

\subsection{Adaptation}

Adults have no learning and no strategic choice. Their ``adaptive
traits'' are fixed rules:

\begin{itemize}
\item a gonotrophic state machine that alternates host seeking, blood
  feeding, egg maturation, and oviposition;
\item host selection by a stochastic roulette over an attraction field
  that combines host counts, species host preferences, indoor and urban
  modifiers, and exponential distance decay;
\item directed movement toward the strongest currently sensed host cell
  and toward the strongest currently sensed breeding cell;
\item skip-oviposition: a gravid female that cannot find an acceptable
  breeding site within \mskey{oviposition.timeout_days} days abandons the
  batch and returns to host seeking;
\item starvation: a female that cannot find a host within
  \mskey{host_seeking.starvation_days} days dies;
\item local dispersal and, when a wind field is supplied, windborne
  migration by gravid females.
\end{itemize}

These rules cause individuals to reproduce observed behaviours
(``move toward the strongest plume'', ``lay in the nearest acceptable
patch'') rather than to maximize an explicit measure of success. The
traits are species parameters; they do not change within a run or across
generations.

\subsection{Objectives}

The model has no explicit objective function and no ranking of
alternatives by expected fitness. The implicit objective of an adult
female is to survive and reproduce: find a host, obtain a blood meal,
mature eggs, and lay them in a suitable patch. The criteria used in the
two ``decisions'' are:

\begin{itemize}
\item \textbf{Host choice:} the attraction $A_j$ of a cell (Equation~17),
  which increases with the number and preferred type of hosts and
  decreases exponentially with distance. The female does not evaluate
  the expected number of offspring, the risk of being killed, or the
  quality of the blood meal.
\item \textbf{Oviposition choice:} the attraction
  $H^{\text{total}}_p e^{-d_p/\mskey{oviposition.scale_m}}$ of an active
  patch (Equation~16). The
  female does not forecast larval survival, competition, or future
  drying.
\end{itemize}

The concept is included to make this explicit: the model is
rule-based, not objective-driven. Any fitness-like quantity (for
example, the approximate effective reproduction number reported in the
transmission log) is a \emph{diagnostic} computed from the emergent
output, not an objective that agents maximize.

\subsection{Learning}

Not implemented. Individuals do not change their behavioural rules with
experience; there is no within-run learning, no memory of past hosts or
breeding sites, and no evolutionary updating of parameters. The only
state that accumulates with experience is the gonotrophic cycle counter
(which distinguishes the first, slower egg maturation from later ones)
and the EIP clock (which is physiological, not cognitive).

\subsection{Prediction}

Not implemented. Agents do not estimate future environmental or internal
conditions. Movement is a response to the current attraction field; a
female does not anticipate where hosts will be at the next phase, where
rain will fall, or whether a breeding site will still be wet when her
eggs hatch. The EIP progress is a physiological accumulator, not a
prediction.

\subsection{Sensing}

Sensing is local and perfect within a radius:

\begin{itemize}
\item \textbf{Hosts:} a female computes an attraction field over cells
  within \mskey{species.coluzzii.host_seeking_radius_m}\,m (the species
  attraction radius) and, for directed
  movement, detects the best host cell within
  \mskey{host_seeking.detect_range_m}\,m (the engine
  detection range). The field uses the true host counts (or the
  phase-specific effective counts when mobility is loaded). There is no
  perceptual error, no occlusion, and no memory.
\item \textbf{Breeding sites:} a gravid female detects active patches
  within \mskey{oviposition.detect_range_m}\,m, with attraction equal to
  the patch's static
  potential times an exponential decay with a \mskey{oviposition.scale_m}\,m
  scale.
\item \textbf{Environment:} temperature, rainfall, and salinity are
  experienced at the cell; the aquatic stages use the cell temperature
  of their patch. There is no sensing of conspecifics, no pheromone or
  sound signalling, and no social network.
\item \textbf{Hosts' perception:} humans and livestock do not perceive or
  respond to mosquitoes. Host counts are a static (or phase-mobile)
  resource; they are not depleted by biting.
\end{itemize}

\subsection{Interaction}

The model has three kinds of interaction:

\begin{enumerate}
\item \textbf{Indirect competition among larvae.} All larvae in a patch
  share the same daily capacity $K_{p,d}$; survival declines with the
  total larval count $N_p$ through a Beverton--Holt function
  (Equation~9). This is exploitation competition for a shared,
  unmodelled resource, represented as a density-dependent survival
  penalty.
\item \textbf{Vector--host contact.} A feeding female bites one host in
  the selected cell. The interaction is one-sided: the host does not
  respond, does not become aware, and does not change behaviour. The bite
  is recorded in a ledger and can infect the vector (human to vector) or
  the human (vector to human).
\item \textbf{Shared attraction.} Adults do not deplete hosts, so there
  is no direct exploitation competition among adults. They interact
  indirectly through the shared attraction field: a cell with many hosts
  attracts many females, and the resulting bites are concentrated there.
\end{enumerate}

There is no direct mosquito--mosquito interaction, no competition for
mating (mating is instantaneous and not modelled), no communication, and
no human behavioural response. Human mobility changes the effective host
distribution and therefore the interaction surface, but it does not
change the infection compartments (Appendix~A.12).

\subsection{Stochasticity}

Stochasticity is used where the model does not represent the causal
details of a process, and to produce demographic variability in a
finite population. The stochastic draws are:

\begin{itemize}
\item sex at emergence and at seeding: Bernoulli(0.5);
\item adult daily survival: Bernoulli($p_d$) per adult;
\item local dispersal: two independent normal draws per dispersing adult,
  clipped to a circle;
\item windborne migration: Bernoulli(\mskey{dispersal.wind_migration_prob})
  per eligible female per day,
  then Bernoulli(\mskey{dispersal.wind_survival}) survival;
\item egg batch: Binomial(\mskey{gonotrophic.egg_batch_n},
  $\mskey{thermal.efd_max_eggs_per_day}\,m_{\text{EFD}}(T)$) per
  oviposition;
\item egg and pupa mortality: Binomial($N$, $m_k$) per cohort per day;
\item larval survival: Binomial($N$, $p_L$) per larval cohort per day;
\item washout: Binomial($N$, $f(R)$) per cohort on heavy-rain days;
\item host selection: one uniform draw per female per feeding attempt,
  used in a weighted roulette over attraction scores;
\item human infection: Binomial($\operatorname{round}(S_H)$, $\lambda_H$)
  per cell per day, with a deterministic fallback $S_H\lambda_H$ when
  $\operatorname{round}(S_H)=0$; human cohort transitions are
  deterministic shifts of the
  queues;
\item phase-specific human and livestock movement: a multinomial flow
  per origin cell, per phase, per day;
\item outbreak focus selection and, in random-viable mode, the choice of
  seeding patches: uniform draws without replacement;
\item host-weighted seeding: a weighted roulette without replacement.
\end{itemize}

The stochasticity is demographic and behavioural, not
``measurement noise''. Environmental variability enters through the
input series (rainfall, temperature, salinity, wind), not through a
random weather generator. The model is not run as an ensemble by
default: one rollout is one trajectory. Multiple rollouts are obtained
by running the engine with different master seeds; each rollout is an
independent trajectory from the same initial conditions and the same
environmental forcing.

A reproducibility caveat is important for replication: the three
aquatic mortality processes (washout, egg/pupa mortality, larval
density mortality) do not use the rollout's random stream. Each uses a
separate \code{std::mt19937_64} re-seeded every day with a fixed
constant (\mskey{rng.aquatic_seed_washout}, \mskey{rng.aquatic_seed_stage_mortality}, and
\mskey{rng.aquatic_seed_larva_mortality}, respectively). Their draws are therefore a
deterministic function of cohort order and cohort counts, identical
across rollouts with the same state. The adult, host-selection, and
transmission draws do use the rollout stream. See Appendix~A.5.

\subsection{Collectives}

No collectives are represented. The model has no herds, swarms, schools,
households, villages, or social groups as entities. A ``patch'' is a
modeler-defined cell; a ``cohort'' is a bookkeeping aggregate of
immatures; a ``grid cell'' is a modeler-defined spatial unit; a
``phase-specific host grid'' is a derived field. None of these is an
emergent aggregation of individuals with its own state and behaviour.
The adult population of a cell is an output aggregate, not an entity.

\subsection{Observation}

The model is observed at full state: every adult, every cohort, and
every cell is known to the observer. There is no sampling, no detection
error, and no virtual ecologist. The observation layer is:

\begin{itemize}
\item \textbf{State COG} (GeoTIFF, float32, 2 bands, NoData $\mskey{ingest.nodata_sentinel}$):
  band 1 \code{adult_patch_count} (adults associated with each cell via
  their patch, after dispersal), band 2 \code{adult_cell_count} (adults
  at their current position). Both are raw counts of all adults, males
  included; there is no normalization by capacity.
\item \textbf{State sidecar} (JSON): fifteen keys, written in the
  emission order of \code{write_state_sidecar}
  (\code{output_contract.cpp:248-271}) ---
  \code{crs} (the EPSG string of the AOI grid),
  \code{transform} (the six-element GDAL affine tuple),
  \code{aoi\_slug} (the AOI identifier),
  \code{scale} (the AOI scale),
  \code{year} and \code{month} (the run's start date),
  \code{seed},
  \code{n\_rollouts} (how many rollouts the run has) and
  \code{rollout\_index} (which one this is, zero-based),
  \code{generator\_version},
  \code{abm\_params\_hash} (a placeholder: the code writes
  \code{"sha256:pending"}, not a computed hash --- Appendix~A),
  \code{contract\_version},
  \code{band\_names} (the two state band names),
  \code{nodata} (the NoData sentinel of the COG), and
  \code{shape} (a three-element array $[2, h, w]$: the number of bands,
  the number of rows, and the number of columns).
\item \textbf{Transmission COG} (optional, 4 bands, float32): human
  prevalence, human incidence per person, infectious vector pressure
  (infectious bites delivered per cell per day), and a binary active-focus
  band (prevalence at or above the focus threshold, or any incidence).
\item \textbf{Aquatic diagnostics} (JSON): total aquatic count, counts by
  stage and instar, top ten patches by count, and mean development by
  stage.
\item \textbf{Cohort log} (JSON, optional): daily adult, larva, egg, and pupa
  counts, emergence and death counts, and the legacy \code{eip_frac} field (Appendix~A.8).
\item \textbf{Transmission log} (JSON, optional): daily human and vector
  compartment counts, bites on humans, infectious bites, and an
  approximate $R_{\text{eff}}$.
\item \textbf{Foci log} (JSON, optional): the cells, cases, cell
  population, mosquito density, and truncation flag of the seeded
  outbreak.
\item \textbf{Animated GIF} (optional): a diagnostic rendering of the
  live state, including individual adults.
\end{itemize}

The state is written at the end of the requested horizon and, optionally,
every $N$ days. The sidecar's year and month are the run's start date,
not the date of the snapshot; the number of elapsed days is implicit in
the file name and the snapshot schedule. The observation is at cell
resolution and daily; the only individual-level visual output is the
optional GIF, which is diagnostic and not part of the data contract.

% =====================================================================
\section{Initialization}

\subsection*{5.1 State at time $t=0$}

A run starts on the first day of the requested year and month. Before
the first call to \code{step()}, the following holds.

\begin{itemize}
\item The environment is loaded: daily rainfall and air temperature for
  the requested horizon (plus any warm-up bands), the five static habitat
  planes, and optional salinity. The active day is the first simulation
  band.
\item The coordinate and capacity fields are \emph{not} yet computed:
  $z=0$ in every cell, the dry-day counters are 0, and $K=0$. They are
  filled by the first \code{update_habitat_day()} call, which consumes
  the first day's rainfall. The pluvial state at the first mosquito
  update is therefore the response to the start date's rain, not the
  cold-start zero.
\item The patch states are not yet built. In the default detection-based
  modes, the first \code{to_dataframe()} call in \code{step()} builds the
  candidate union and assigns stable patch identifiers; the uniform mode
  builds the day-0 union during construction.
\item The adult and aquatic populations are those produced by the
  seeding mode (Section~5.2).
\item The human compartments are initialized from the host grid
  (Section~5.3).
\item No wind field, mobility schedule, or transmission model is active
  unless the corresponding input was supplied or the transmission flag
  was set.
\end{itemize}

If \code{--warmup-bands} $n>0$ is supplied, the pluvial recursion is
integrated over the $n$ environmental bands immediately preceding the
start band, starting from $z=0$ at the cold end, before the simulation
begins. This is a causal warm-up: it uses real preceding rainfall, not a
synthetic fixed window. The default is $n=0$ (cold start).

\subsection*{5.2 Mosquito seeding}

The engine supports four seeding modes. The default is
\code{random-viable}.

\begin{longtable}{@{}L{2.3cm}L{3.3cm}L{4.1cm}L{4.9cm}@{}}
\caption{Seeding modes and their initial populations. Counts are per
selected patch unless stated.}\label{tab:seed}\\
\toprule
\textbf{Mode} & \textbf{Selection} & \textbf{Adults} & \textbf{Immatures} \\
\midrule
\endfirsthead
\toprule
\textbf{Mode} & \textbf{Selection} & \textbf{Adults} & \textbf{Immatures} \\
\midrule
\endhead
\bottomrule
\endlastfoot
\code{random-viable} (default) & 3 patches drawn uniformly without replacement from the viable set (gpkg potential $\ge\mskey{habitat.candidate_min_engine}$) & 50 per patch & 30 per patch, distributed over stages \\
\code{uniform} & every candidate cell of the day-0 union & 30 per patch & 270 per patch, distributed over stages \\
\code{explicit} & user-supplied points; nearest viable patch within 5\,km; points with no patch in range are dropped & 50 per point (fixed by the parser) & 30 per point (fixed by the parser) \\
\code{host-weighted} & 3 patches sampled without replacement with weight $\sum_j \text{host score}_j e^{-d_{pj}/s_s}$ within 5\,km, $s_s$ the species plume scale (\mskey{species.coluzzii.host_seeking_scale_m}\,m for the default species) & 50 per patch & 30 per patch \\
\end{longtable}

Every seeded adult is created with \code{stage}=1, a sex drawn
Bernoulli(0.5), gonotrophic state TENERAL, vector state SUSCEPTIBLE,
EIP progress 0, zero completed cycles, and the cell-centre lon/lat of its
patch. Adults are immediately able to disperse, but a TENERAL female
must pass through MATE\_SEEKING before she can seek a host: her first
blood meal occurs on the second daily update after seeding (or later, if
no host is found).

The immature count of a detection-seeded patch is distributed as
follows. The integer allocation is computed from the requested count
$n$ as
\[
n_{\text{pupa}}=\max(1,\lfloor0.05n\rfloor),\quad
n_{L4}=\max(1,\lfloor0.15n\rfloor),\quad
n_{L3}=n_{L2}=\lfloor0.25n\rfloor,\quad
n_{L1}=\lfloor0.20n\rfloor,
\]
and the remainder is eggs. The pre-set development fractions are 0.58
for pupae, 0.47 for L4, 0.38 for L3, 0.29 for L2, 0.20 for L1, and 0.00
for eggs. These fractions are fractions of the current stage, not ages;
they stagger the first emergence wave instead of releasing all immatures
at once. For the default $n=30$ the allocation is 5 eggs, 6 L1, 7 L2,
7 L3, 4 L4, and 1 pupa; for the uniform $n=270$ it is 29 eggs, 54 L1,
67 L2, 67 L3, 40 L4, and 13 pupae. The small-$n$ rounding is therefore
visible: the nominal 10/20/25/25/15/5 percentages become 16.7/20/23.3/
23.3/13.3/3.3 at $n=30$.

If the candidate set is empty, or if no seeding instruction is produced
(no viable patch, no host within range in host-weighted mode, no
detection point in range in explicit mode), the population starts empty.
An empty population cannot produce a new generation, because there are
no adults to lay eggs; the run then reports zero vectors throughout.

\subsection*{5.3 Humans, transmission, hosts, mobility, and wind}

\textbf{Human compartments.} The human population of each cell is the
value of the \code{human} variable in the host NetCDF (WorldPop counts
aggregated to the 1\,km cell). If no host file is supplied but
transmission is enabled, every cell is assigned one person. The initial
compartments are $S_H=$ population and $E_H=I_H=R_H=0$, except in the
legacy \code{uniform-legacy} seeding mode, where a fraction
\code{initial_human_prevalence} $=\mskey{transmission.initial_human_prevalence}$ of each cell is infectious and is
spread evenly over the \mskey{transmission.human_infectious_days} infectious cohort bins. The human population
does not change during the run.

\textbf{Outbreak seeding.} If transmission is enabled and the human
seeding mode is \code{random-viable} (the default), the model waits until
the configured outbreak day (default 0). It first checks that the maximum
cell density reaches \code{human_outbreak_min_density} (default 0, which
disables the gate), then looks for a cell with population greater than or equal to
\code{human_min_cell_pop} (default 50; the test is
\code{population_ >= min_pop}, \code{human_state.cpp:112}) and vector presence in its
$3\times3$ neighbourhood (in cluster mode, at least
\code{human_outbreak_min_density}). It seeds
\code{human_outbreak_foci} (default 3) foci of
\code{human_outbreak_cases} (default 50) infectious humans, truncated by
the susceptible population of each cell. If no cell qualifies, it retries
on the next day. In cluster mode (a positive cluster radius), it seeds
one cluster instead: the highest-population qualifying core cell, with
cases split across the Chebyshev neighbourhood in proportion to
population times $(0.05+\text{local density})$. The \code{explicit} mode
uses user coordinates instead of the viability search; \code{none} and
\code{uniform-legacy} do not seed at trigger time. The default vector
seeding fraction is \mskey{transmission.initial_vector_infected_frac}; if it is positive, that fraction of adult females
is created already INFECTIOUS with EIP progress at the threshold.

\textbf{Host landscape.} If a host file is supplied, the per-cell human
and livestock counts, building fraction, indoor fraction, and urbanicity
are loaded. If it is not supplied, there is no host landscape, no host
seeking, and the fallback feeding rule is used (Section~7.6). The host
landscape is static except for the phase-specific redistribution by
mobility.

\textbf{Mobility.} If human and/or livestock origin--destination matrices
are present, the phase-specific effective host grids are built; if they
are absent, everyone stays in their residential cell. The phase-aware
host-seeking pass requires mobility data: without them, the single-phase
residential fallback is used.

\textbf{Wind.} If a wind field is supplied, windborne migration is active
in the months and at the hour described in Section~7.7; otherwise it is
off. The delivered Ghana product has no wind field in the manifest, so
wind migration is off by default.

\textbf{Transmission.} The transmission layer is off unless
\code{--enable-transmission} is set. When it is off, the bite ledger is
still maintained (it is cheap) but no EIP, human infection, or
transmission output is produced.

\subsection*{5.4 Randomness and reproducibility}

The master seed is \code{--seed} (default 1). Rollout $i$ uses
$\text{seed}+i$. From that stream the engine derives, in order, the
coordinator seed (first state word), one uniform draw, the submodel seed
(second state word), and, if transmission is enabled, the transmission
seed (a further state word plus one uniform draw). The coordinator stream is
used for seeding selection; the submodel stream is used for all daily
adult, host-selection, and egg-batch draws; the transmission stream is
used for outbreak and human binomial draws. The phase-specific host
flows use a named substream derived by splitmix64 from the day, phase,
and host type, so they do not consume the main stream. As noted in
Section~4.9, the three aquatic mortality processes use fixed daily seeds
and are independent of the rollout seed.

Two runs with the same seed, the same inputs, and the same binary produce
byte-identical state outputs. Two runs with different seeds share the
environmental forcing but differ in seeding selection and in every
stochastic draw. Initialization is never varied by a random environment:
the climate and habitat inputs are fixed files.

\subsection*{5.5 What is data-based and what is a scenario choice}

The environmental drivers, the static habitat planes, the host counts,
the building fraction, the urban class, and the optional mobility and
wind fields are data-based. The initial mosquito counts, the number of
detection patches, the detection radius, the outbreak size, and the
cluster radius are scenario choices, not estimates. The model constants
are literature values or explicit engineering priors, as recorded in the
constants file and in the parameter tables of Element~7. The project states
that no parameter has been fitted to Ghanaian mosquito abundance in the
current release; the code records no fit provenance, so this is a project
statement rather than a code-verifiable fact
(Appendix~B.4).

% =====================================================================
\section{Input data}

The model uses input data to represent time-varying and spatially varying
processes. It does \emph{not} use input data for parameter values or for
initial state variables: those are the subject of Elements~5 and~7. All
inputs are prepared by the project's ingest stage from openly available
products and are read as files at run time; the engine does not download
data and does not call another model.

\begin{longtable}{@{}L{2.4cm}L{3.2cm}L{2.6cm}L{2.6cm}L{3.2cm}@{}}
\caption{Input data. ``1 km'' means the AOI grid cell; the native
resolution of the source is given in the preparation column.}\label{tab:input}\\
\toprule
\textbf{Input} & \textbf{Variable} & \textbf{Units and resolution} & \textbf{Source} & \textbf{Preparation and uncertainty} \\
\midrule
\endfirsthead
\toprule
\textbf{Input} & \textbf{Variable} & \textbf{Units and resolution} & \textbf{Source} & \textbf{Preparation and uncertainty} \\
\midrule
\endhead
\bottomrule
\endlastfoot
Daily rainfall & \bcode{rainfall} & mm\,d$^{-1}$; 1 km, daily & CHIRPS v2.0 (UCSB Climate Hazards Center), native 0.05$^\circ$ ($\sim$5\,km) & Nearest-neighbour reprojection to the AOI grid; NoData filled from the nearest valid cell for that day; open sea set to 0; negative values clipped. A 5\,km product cannot resolve local convective cells; daily totals are smeared over the 1\,km grid. \\
Daily air temperature & \bcode{air_temp_c} & $^\circ$C; 1 km, daily mean & ERA5-Land 2\,m temperature, native $\sim$9\,km & Bilinear reprojection; open sea set to 0; finite values below $-900$ treated as NoData and replaced by an eight-neighbour wavefront imputation from land. Daily means, not diurnal cycles. \\
Permanent surface water & \bcode{jrc_persistent} & fraction $[0,1]$; 1 km, static & JRC Global Surface Water v1.4, native 30\,m & Fraction of pixels with $\ge\codeconst{ingest.jrc_occurrence_threshold_pct}\%$ occurrence, averaged to 1\,km; open-sea cells zeroed, coastal lagoons and river mouths inside the configured coastal band retained (\mskey{ingest.coastline_buffer_m}\,m in the current constants; the tree records no value for the setting used when the delivered NC was built, so it is unknown whether the delivered product used this band). Occurrence is not depth, volume, or breeding suitability. \\
Built land & \bcode{lc_built} & fraction $[0,1]$; 1 km, static & Dynamic World V1 (Google), native 10\,m & Per-year mode-of-year label composite, then areal fraction per 1\,km cell (average resampling); mean over available years. Label 6. Mode-of-year can lose episodic construction. \\
Wetland & \bcode{lc_wetland} & fraction $[0,1]$; 1 km, static & Dynamic World V1 & Areal fraction of label 3 (flooded vegetation) only; mangroves are not separately represented. Tidal and deciduous pixels can be assigned to trees in the mode composite. \\
Open land and water & \bcode{lc_open_plus_water} & fraction $[0,1]$; 1 km, static & Dynamic World V1 & Areal fraction of labels 2 (grass), 5 (shrub), 7 (bare), and 0 (water). In the engine, open habitat adds $\mskey{urban.rain_pools_built_damage}\times$ built land to represent rain pools on sealed surfaces. \\
Cropland & \bcode{lc_crop} & fraction $[0,1]$; 1 km, static & Dynamic World V1 & Areal fraction of label 4. \\
Salinity & \bcode{salinity_ppt} & psu; 1 km, daily, optional & SMAP monthly sea-surface salinity & Monthly field broadcast to days; missing values and negatives set to 0 (freshwater). Absent from the delivered Ghana product, so every cell is freshwater there. \\
Habitat patches & \bcode{habitat_potential}, \bcode{hab_type}, \bcode{row}, \bcode{col} & points; 1 km & Derived from JRC and Dynamic World by ingest & 382,127 point features in the Ghana 2024 product, one per grid cell with any static habitat potential; 368,761 of the features are at or above the \mskey{habitat.candidate_min_engine} engine floor, and 370,061 grid cells clear that floor on their static total potential. The engine's patch set is the union of the file's cells and the floor cells, i.e.\ 382,127 cells in this product. The amplitude is the sum of the six $H$ values, capped at 1. \\
Human population & \bcode{human} & persons per cell; 1 km, static & WorldPop 2020 constrained counts & Sum-preserving aggregation from $\sim$100\,m to 1\,km. A modelled count, not a census; about 32.0 million persons in the delivered bounding box. \\
Livestock & \bcode{cattle}, \bcode{goats}, \bcode{sheep}, \bcode{pigs}, \bcode{chickens} & animals per cell; 1 km, static & FAO GLW4-2020 & Sum-preserving aggregation from $\sim$10\,km. About 3.5\,M cattle, 14.8\,M goats, 8.6\,M sheep, 2.7\,M pigs, and 121\,M chickens in the delivered grid. \\
Urban class & \bcode{urban_class} & class code; 1 km, static & GHS-SMOD R2023A & Nearest-neighbour aggregation. The engine's binary rule treats class $\le40$ as urban (indoor fraction \mskey{host_seeking.indoor_fraction_urban}, urbanicity 1) and $>40$ as rural (\mskey{host_seeking.indoor_fraction_rural}, 0). \\
Building fraction & \bcode{building_fraction} & $[0,1]$; 1 km, static & Overture Maps / OpenStreetMap & Sum-preserving aggregation. Used only to derive the residential fraction. \\
Wildlife proxy & \bcode{wildlife_host_proxy} & index $[0,1]$; 1 km, static & Dynamic World class suitability, JRC water, remoteness & Weighted combination \mskey{wildlife.weights.habitat} habitat suitability, \mskey{wildlife.weights.water} water, \mskey{wildlife.weights.remoteness} remoteness. A proxy for wild-host availability, not a count. \\
Human mobility & \bcode{*_mobility_day.csr}, \bcode{*_mobility_night.csr} & row-stochastic fractions; cell-to-cell & Built by ingest from population and attractiveness & Sparse origin--destination matrices; no within-run update. Absent $\Rightarrow$ identity (everyone at home). \\
Livestock mobility & \bcode{*_livestock_mobility.csr} & row-stochastic fractions; cell-to-cell & Built by ingest & One seasonal matrix reused for all phases except NIGHT (identity). \\
Wind & \bcode{u100}, \bcode{v100} & m\,s$^{-1}$; 6-hourly, optional & ERA5 & Nearest-neighbour spatial lookup and linear temporal interpolation; queried at one hour per day. Absent from the delivered Ghana product. \\
\end{longtable}

\textbf{Uncertainties and biases.} The habitat layers are proxies, not
measurements of breeding sites. JRC occurrence records where water was
observed, not how deep or how long it persists; the \codeconst{ingest.jrc_occurrence_threshold_pct}\% threshold turns
a continuous occurrence field into a persistent-water mask. Dynamic
World's mode-of-year composite removes within-year phenology and can
misclassify tidal wetlands and mangrove margins. CHIRPS at 5\,km and
ERA5-Land at 9\,km cannot resolve the local rain cells and temperature
gradients that drive pool filling and larval development at the 1\,km
scale; the model's daily series are therefore smoother than the real
drivers. WorldPop and GLW4 are modelled surfaces with their own
uncertainty, and the 1\,km aggregation mixes residential and
non-residential land. The static habitat planes do not change within a
year, so the model cannot represent land-cover change, irrigation
expansion, or the seasonal drying of a specific pool beyond the shared
pluvial state. The optional salinity field is monthly and absent from the
delivered product, so coastal brackish habitat is effectively freshwater
in the default Ghana runs. Open-sea cells are not excluded from the
open+water static plane, and their air temperature is the 0.0 sentinel
rather than NoData; they therefore act as cold, non-developing habitat
rather than being masked out (Appendix~A.14). The coastal band applied when the delivered product was built is not
recorded in the tree, so it cannot be verified from the repository whether
the delivered JRC mask used the current buffer. These limitations are
structural, not numerical: more compute does not remove them.

% =====================================================================
\section{Submodels}

The submodels below correspond one-to-one to the processes listed in
Element~3. For each, the factual description (equations and algorithm)
comes first, followed by the origin and rationale, and then a table of
parameters with units, reference values, and provenance. All values are
read from \code{model_constants.yaml}; the tables quote that file. A
value marked \emph{engineering prior} has no measurement or fit behind
it: it is a design choice made to keep the model runnable in the absence
of local data. A value marked \emph{literature anchor} is a published
value or a published functional form, but it has not been re-fitted to
Ghanaian data for this model.

The model has been tested for software correctness (unit tests, golden
fixtures, and a Python--C++ parity check of the constants file) but not
validated against field observations. The distinction matters: a passing test suite says that
the code does what this document says, not that the parameter values are
right for a particular district. Appendix~B summarises the test and
calibration status.

\subsection{Update habitat availability and capacity}

For every cell, once per day and before any mosquito acts, the model
updates the pluvial state from the day's rainfall,
\begin{equation}\label{eq:z}
z_{i,d}=z^{\text{dry}}_{i,d}+\left(1-z^{\text{dry}}_{i,d}\right)
\left(1-e^{-R_{i,d}/\rho}\right),\qquad
z^{\text{dry}}_{i,d}=z_{i,d-1}\,e^{-\Delta/\tau},
\end{equation}
where $\rho$ is the rainfall scale (\mskey{pluvial.rain_scale_mm}\,mm), $\tau$ the dry-memory time constant
(\mskey{pluvial.dry_memory_days}\,d), and $\Delta=\mskey{pluvial.delta_days}$\,d.

\textbf{Two vocabularies for habitat types.} The configuration file and
the Python contract layer name six habitat types ---
\mskey{habitat_types.types} --- of which three are persistent and three
pluvial (\code{habitat_types.dynamics}). The engine carries six static
slots of its own, in the fixed order surface, urban, wetland, open,
crop, forest, and the forest slot is identically zero
(\code{habitat_availability.hpp:55-85}; \code{zone_id.hpp:13-21}). The
engine maps five of the six configured types onto those slots by name:
\code{surface_persistent} onto the surface slot,
\code{urban_persistent} onto the urban slot,
\code{wetland_persistent} onto the wetland slot,
\code{open_pluvial} onto the open slot and \code{crop_pluvial} onto
the crop slot (\code{config.cc:472-483}). The configured type
\code{wetland_pluvial} has no engine slot: in the engine the wetland
potential enters once, through the persistent wetland slot. The Python
mirror is different: it defines six potentials, with
$H_{\text{wetland\_pluvial}}=W$, so the wetland fraction is counted
twice there (\code{mal_core/habitat/potential.py:74}). The two
implementations do not agree on the number of habitat types; the
engine's behaviour is what the rest of this section describes.

The static potential of the engine's slots is assembled from the five
environmental planes
($J$ = JRC persistent fraction, $B$ = built fraction, $W$ = wetland
fraction, $O$ = open+water fraction, $C$ = crop fraction):
\begin{equation}\label{eq:H}
H_{i,g}=\begin{cases}
\min(1,\max(0,J_i)), & g=\text{surface},\\
\min(1,\mskey{urban.beta_urban}\,B_i), & g=\text{urban},\\
\min(1,\max(0,W_i)), & g=\text{wetland},\\
\min(1,\max(0,O_i)+\mskey{urban.rain_pools_built_damage}\,B_i), & g=\text{open},\\
\min(1,\max(0,C_i)), & g=\text{crop},\\
0, & g=\text{forest}.
\end{cases}
\end{equation}
The daily availability is a weighted sum of the static potential, with
persistent types at full availability and pluvial types at $z$:
\begin{equation}\label{eq:S}
S_{i,d}=\min\left(1,\sum_g w_g H_{i,g}A_{g,d}\right),\qquad
A_{g,d}=\begin{cases}1,&g\text{ persistent},\\ z_{i,d},&g\text{ pluvial}.\end{cases}
\end{equation}
The persistent/pluvial split is not positional: it is read from the
configured dynamics of each type by name and mapped onto the engine's
slots, so the surface, urban and wetland slots are persistent and the
open and crop slots are pluvial; the forest slot has no configured type
and is always zero.
The absolute daily capacity is
\begin{equation}\label{eq:K}
K_{i,d}=k_{\text{ref}}\,a_i\,S_{i,d},\qquad
a_i=10^{-6}\,(\mskey{ingest.cell_size_hint_m})^2\cos\phi_i,
\end{equation}
with $k_{\text{ref}}=\mskey{species.k_ref}$ and $\phi_i$ the latitude of
the cell. A cell is an active patch when
$S_{i,d}\ge S_{\text{act}}=\mskey{habitat.activation_min_engine}$ and its salinity is below the active
species' high tolerance; otherwise it is inactive. Separately, a dry-day
counter $D_{i,d}$ is reset to zero on days when $S_{i,d}\ge S_{\text{act}}$
and incremented otherwise; it drives desiccation in the aquatic submodel.
The counter therefore tracks $S$, not the salinity gate: a cell that is
deactivated by salinity can still have $D_{i,d}=0$. The candidate patch
union is the habitat GeoPackage plus any cell whose static total
potential clears \mskey{habitat.candidate_min_engine}; the engine's total
potential is the sum over its five non-zero slots, so the wetland
fraction enters it once; patches are emitted in row-major order with
stable identifiers.

\begin{longtable}{@{}L{3.3cm}L{4.4cm}L{1.6cm}L{1.6cm}L{3.6cm}@{}}
\caption{Parameters of the habitat update.}\label{tab:habparams}\\
\toprule
\textbf{Parameter} & \textbf{Meaning} & \textbf{Unit} & \textbf{Value} & \textbf{Provenance} \\
\midrule
\endfirsthead
\toprule
\textbf{Parameter} & \textbf{Meaning} & \textbf{Unit} & \textbf{Value} & \textbf{Provenance} \\
\midrule
\endhead
\bottomrule
\endlastfoot
\code{rain_scale_mm} & rainfall scale $\rho$ of the pluvial fill & mm & \mskey{pluvial.rain_scale_mm} & engineering prior, aligned with the d19/d20 scorers \\
\code{dry_memory_days} & dry-memory time constant $\tau$ & d & \mskey{pluvial.dry_memory_days} & engineering prior \\
\code{delta_days} & daily step $\Delta$ & d & \mskey{pluvial.delta_days} & model definition \\
\code{cold_start_z} & initial pluvial state & -- & \mskey{pluvial.cold_start_z} & model definition \\
\code{activation_min_engine} & $S_{\text{act}}$, active-patch floor & -- & \mskey{habitat.activation_min_engine} & engineering prior \\
\code{candidate_min_engine} & engine patch-set floor & -- & \mskey{habitat.candidate_min_engine} & engineering prior \\
\code{beta_urban} & coefficient of built land in $H_{\text{urban}}$ & -- & \mskey{urban.beta_urban} & engineering prior; unchanged from the deployed calibration \\
\code{rain_pools_built_damage} & built-land term added to $H_{\text{open}}$ & -- & \mskey{urban.rain_pools_built_damage} & engineering prior \\
\code{persistent_availability} & $A$ for persistent types & -- & \mskey{capacity.persistent_availability} & model definition \\
\code{k_ref} & capacity scale & ind.\,km$^{-2}$ & \mskey{species.k_ref} & engineering prior; see Appendix~A.9 on the unit annotation \\
\code{cell_size_hint_m} & default AOI cell size & m & \mskey{ingest.cell_size_hint_m} & model definition \\
\end{longtable}

\textbf{Origin and rationale.} The submodel replaces the retired pool
hydrology (water balance in millimetres, topographic wetness index, urban
floor, capacity clamp). The replacement is deliberately simpler: the
open products support a static areal fraction and a rain-driven
availability, not a measured water depth. The urban term encodes the
premise that leaks, seepage, drainage, and human water storage maintain
some dry-season breeding supply in built-up areas; it is a proxy, not a
measurement. The pluvial state is shared by all pluvial habitat types;
the model does not maintain a separate water balance per type.

\textbf{Behaviour on the delivered landscape.} Evaluating
Equations~\eqref{eq:H}--\eqref{eq:K} directly on the Ghana 2024 static
planes gives a useful characterization of the initial habitat. With the
species-neutral weights of Appendix~A.1, 11.0\% of cells are above the
\mskey{habitat.activation_min_engine} activation floor at $z=0$
(persistent supply only), and 82.9\% are
above it at $z=1$; the mean capacity rises from 0.61 to 10.4 individuals
per cell, with a maximum of about 32 at $z=1$. These numbers are computed
from the delivered planes and the equations above; they are not the
result of an engine run, they are readings of the delivered artefact
(Appendix~B.4), and they will change with a different AOI
or year.

\subsection{Advance aquatic stages}

Immatures live in cohorts indexed by (patch, stage, instar). Each cohort
has a count $N$ and a development fraction $x\in[0,1)$ relative to its
current stage. Development is temperature-dependent. For eggs and pupae
the rate is Brière-1; for larvae it is quadratic:
\begin{equation}\label{eq:thermal}
r_{\text{egg/pupa}}(T)=a\,T\,(T-T_L)\sqrt{T_H-T},\qquad
r_{\text{larva}}(T)=a\,(T-T_L)(T_H-T),
\end{equation}
for $T_L<T<T_H$ and zero otherwise, clipped to $[0,1]$ per day. Each
active-patch cohort accumulates
\begin{equation}\label{eq:dev}
x_{k,d}=x_{k,d-1}+r_k(T_{p,d});
\end{equation}
a cohort with $x\ge1$ is promoted to the next instar or stage, and its
development is reset to zero. Promotions are collected first and applied
after the scan, so they do not cascade within a day. Cohorts on inactive
patches do not develop; their development is frozen.

Larval survival is density-dependent and shared by all four instars of a
patch. With $N_p$ the total larval count of the patch before mortality,
\begin{equation}\label{eq:BH}
p_L=\begin{cases}
S_0\,\dfrac{K_{p,d}}{K_{p,d}+\alpha N_p}, & K_{p,d}>0,\\[6pt]
0, & K_{p,d}=0,
\end{cases}
\qquad p_L\leftarrow p_L\,s(\sigma_{p,d}),
\end{equation}
where $S_0=\mskey{aquatic_stages.larva_bh_s0}$ and $\alpha=\mskey{aquatic_stages.larva_bh_alpha}$. If the patch has zero capacity, all
larvae of that patch die on that day; eggs and pupae are not touched by
this operation. The salinity multiplier is
\begin{equation}\label{eq:sal}
s(\sigma)=\begin{cases}
1, & \sigma\le\sigma_{\text{opt}},\\
\left(1-\dfrac{\sigma-\sigma_{\text{opt}}}{\sigma_{\text{hi}}-\sigma_{\text{opt}}}\right)
e^{-b(\sigma-\sigma_{\text{opt}})}, & \sigma_{\text{opt}}<\sigma<\sigma_{\text{hi}},\\
0, & \sigma\ge\sigma_{\text{hi}}.
\end{cases}
\end{equation}
Eggs and pupae have constant daily mortality; washout and desiccation are
applied before development:
\begin{equation}\label{eq:stress}
m_{\text{des}}(D)=\begin{cases}0,&D\le\mskey{aquatic_stress.desicc_grace_days},\\ \mskey{aquatic_stress.desicc_base_daily}\left(1-e^{-\mskey{aquatic_stress.desicc_ramp_decay}(D-\mskey{aquatic_stress.desicc_grace_days})}\right),&D>\mskey{aquatic_stress.desicc_grace_days},\end{cases}
\qquad
f(R)=\begin{cases}0,&R<\mskey{aquatic_stress.wash_threshold_mm},\\ \min\left(\mskey{aquatic_stress.wash_fraction_max},\,\mskey{aquatic_stress.wash_fraction_max}\,\dfrac{R-\mskey{aquatic_stress.wash_threshold_mm}}{\mskey{aquatic_stress.wash_threshold_mm}}\right),&R\ge\mskey{aquatic_stress.wash_threshold_mm}.\end{cases}
\end{equation}
Desiccation applies only to eggs and to larvae in instars L1--L2
(\code{vulnerable_instar_max}=\mskey{aquatic_stress.vulnerable_instar_max}); L3, L4, and pupae are not desiccated.
Washout applies to all aquatic stages when daily rainfall reaches the
\mskey{aquatic_stress.wash_threshold_mm}\,mm threshold; the flushed fraction is binomial. The daily order is:
desiccation, washout, egg/pupa mortality, larval density mortality,
development and promotion, compaction.

\begin{longtable}{@{}L{4.2cm}L{4.0cm}L{1.4cm}L{1.4cm}L{3.4cm}@{}}
\caption{Parameters of the aquatic-stage submodel.}\label{tab:aquaparams}\\
\toprule
\textbf{Parameter} & \textbf{Meaning} & \textbf{Unit} & \textbf{Value} & \textbf{Provenance} \\
\midrule
\endfirsthead
\toprule
\textbf{Parameter} & \textbf{Meaning} & \textbf{Unit} & \textbf{Value} & \textbf{Provenance} \\
\midrule
\endhead
\bottomrule
\endlastfoot
\code{egg_development_tmin_c} & $T_L$, egg Brière-1 & $^\circ$C & \mskey{thermal.egg_development_tmin_c} & fit limit on Bayoh \& Lindsay (2003) \\
\code{egg_development_tmax_c} & $T_H$, egg Brière-1 & $^\circ$C & \mskey{thermal.egg_development_tmax_c} & fit limit \\
\code{egg_development_scale} & $a$, egg & d$^{-1}$\,$^\circ$C$^{-2.5}$ & \mskey{thermal.egg_development_scale} & fit to Bayoh \& Lindsay (2003) \\
\code{larva_development_tmin_c} & $T_L$, larva quadratic & $^\circ$C & \mskey{thermal.larva_development_tmin_c} & fit limit \\
\code{larva_development_tmax_c} & $T_H$, larva quadratic & $^\circ$C & \mskey{thermal.larva_development_tmax_c} & fit limit \\
\code{larva_development_scale} & $a$, larva & d$^{-1}$\,$^\circ$C$^{-2}$ & \mskey{thermal.larva_development_scale} & calibrated on Bayoh \& Lindsay (2003); validated on Agyekum (2022) and Ou\'edraogo (2024) per the constants file \\
\code{pupa_development_tmin_c} & $T_L$, pupa Brière-1 & $^\circ$C & \mskey{thermal.pupa_development_tmin_c} & fit limit \\
\code{pupa_development_tmax_c} & $T_H$, pupa Brière-1 & $^\circ$C & \mskey{thermal.pupa_development_tmax_c} & fit limit \\
\code{pupa_development_scale} & $a$, pupa & d$^{-1}$\,$^\circ$C$^{-2.5}$ & \mskey{thermal.pupa_development_scale} & fit to Bayoh \& Lindsay (2003) \\
\code{egg_daily_mort} & daily egg mortality & d$^{-1}$ & \mskey{aquatic_stages.egg_daily_mort} & engineering prior; the constants file notes that Depinay (2004) implies about 0.01 \\
\code{pupa_daily_mort} & daily pupa mortality & d$^{-1}$ & \mskey{aquatic_stages.pupa_daily_mort} & engineering prior \\
\code{larva_bh_s0} & zero-density larval survival $S_0$ & -- & \mskey{aquatic_stages.larva_bh_s0} & engineering prior \\
\code{larva_bh_alpha} & competition coefficient $\alpha$ & -- & \mskey{aquatic_stages.larva_bh_alpha} & engineering prior; CLI-overridable \\
\code{last_larval_instar} & last larval instar & instar & \mskey{aquatic_stages.last_larval_instar} & model definition \\
\code{desicc_grace_days} & desiccation grace period & d & \mskey{aquatic_stress.desicc_grace_days} & engineering prior \\
\code{desicc_base_daily} & desiccation ceiling & d$^{-1}$ & \mskey{aquatic_stress.desicc_base_daily} & engineering prior \\
\code{desicc_ramp_decay} & desiccation ramp & d$^{-1}$ & \mskey{aquatic_stress.desicc_ramp_decay} & engineering prior \\
\code{wash_threshold_mm} & washout threshold & mm & \mskey{aquatic_stress.wash_threshold_mm} & engineering prior \\
\code{wash_fraction_max} & maximum washed fraction & -- & \mskey{aquatic_stress.wash_fraction_max} & engineering prior \\
\code{vulnerable_instar_max} & last desiccation-vulnerable instar & instar & \mskey{aquatic_stress.vulnerable_instar_max} & model definition \\
\code{salinity_missing_psu} & salinity used when data are absent & psu & \mskey{aquatic_stages.salinity_missing_psu} & model definition (freshwater) \\
\code{salinity_opt_ppt}, \code{salinity_hi_tol_ppt}, \code{salinity_slope} & species salinity response & psu, psu, psu$^{-1}$ & \mskey{species.coluzzii.salinity_opt_ppt}, \mskey{species.coluzzii.salinity_hi_tol_ppt}, \mskey{species.coluzzii.salinity_slope} & \emph{An.\ coluzzii} row; literature-anchored but not locally fitted \\
\code{aquatic_seed_washout}, \code{aquatic_seed_stage_mortality}, \code{aquatic_seed_larva_mortality} & fixed daily RNG seeds & -- & \mskey{rng.aquatic_seed_washout}, \mskey{rng.aquatic_seed_stage_mortality}, \mskey{rng.aquatic_seed_larva_mortality} & reproducibility constants \\
\end{longtable}

\textbf{Origin and rationale.} The cohort formulation keeps the immature
stages cheap while preserving their temperature-dependent timing and
their density-dependent loss. The thermal curves are the standard
mosquito development forms; the values are fitted to laboratory
development data, not to Ghanaian field populations. The Beverton--Holt
form is a standard compensatory survival function: at low density,
survival approaches $S_0$; at high density, survival declines
hyperbolically, with half of the density-dependent loss at
$N_p=K/\alpha=20K$. The desiccation and washout stressors are the
retained biological stresses of the retired pool model, re-driven by the
new habitat state. The salinity multiplier is an independent gate: it
multiplies larval survival, so brackish water reduces, and above the
high tolerance eliminates, larval survival for a freshwater-adapted
species.

\textbf{Consequences worth stating.} At 25$^\circ$C, the rates give
about 1.1\,d for eggs, 4.4\,d per larval instar (17.8\,d for four
instars), and 1\,d for pupae, i.e.\ roughly 20\,d from egg to adult;
colder and hotter temperatures lengthen the larval stage because the
quadratic curve peaks at 25$^\circ$C. A dry patch with $S$ below the
activation floor freezes development but does not automatically kill
L3, L4, or pupae; those cohorts persist until the patch reactivates or
until a day with $K=0$ removes the larvae. Eggs and L1--L2 are the
stages exposed to desiccation, consistent with the configured
vulnerability window.

\subsection{Emerge adults}

When a pupa cohort reaches $x\ge1$ in the promotion pass, its count is
converted to adults in a synchronous event list applied after the cohort
scan. Each new adult is assigned a sex by a Bernoulli(0.5) draw, the
gonotrophic state TENERAL, vector state SUSCEPTIBLE, EIP progress 0, zero
gonotrophic cycles, and the cell centre of its patch. The new adults are
appended to the SoA and are subject to adult survival in the same day's
survival pass. There are no parameters specific to this submodel beyond
the 0.5 sex ratio.

\textbf{Rationale.} Emergence is the only route by which the adult
population grows; there is no immigration and no adult release except
the initial seeding and, optionally, windborne migration. The 0.5 sex
ratio is a neutral default; it is not a fitted value.

\subsection{Run the gonotrophic cycle}

Each adult female advances one step of a deterministic state machine on
each daily update:
\begin{center}
\small
TENERAL $\to$ MATE\_SEEKING $\to$ HOST\_SEEKING $\to$ BLOOD\_FED
$\to$ RESTING $\to$ EGG\_MATURING $\to$ GRAVID
$\to$ OVIPOSITION\_SEEKING $\to$ OVIPOSITING $\to$ HOST\_SEEKING.
\end{center}
Mating is instantaneous (MATE\_SEEKING lasts one update). The female
remains in HOST\_SEEKING until the host-seeking submodel gives her a
blood meal. After feeding, BLOOD\_FED lasts one daily update (the
configured \mskey{gonotrophic.resting_duration_days}\,d is rounded down to zero completed days), then RESTING
lasts one update, and EGG\_MATURING lasts
$\lceil\mskey{gonotrophic.first_cycle_days}\rceil=\mskey{gonotrophic.first_cycle_days}$ days on the
first cycle and $\lceil\mskey{gonotrophic.cycle_duration_days}\rceil=\mskey{gonotrophic.cycle_duration_days}$ days on later cycles. GRAVID lasts one
update; OVIPOSITION\_SEEKING is resolved by the oviposition submodel;
OVIPOSITING resets the female to HOST\_SEEKING after laying. In the
absence of delays, the interval from a blood meal to the next oviposition
is therefore about 5 days for later cycles and about 7 days for the first
cycle. The two states HOST\_APPROACH and PROBING exist in the enum and in
the transition function but are not entered by the current feeding code;
the female goes directly from HOST\_SEEKING to BLOOD\_FED when a host is
found (Appendix~A.2).

The egg batch is drawn at oviposition:
\begin{equation}\label{eq:eggs}
n_{\text{eggs}}\sim\text{Binomial}\left(\mskey{gonotrophic.egg_batch_n},\;\mskey{thermal.efd_max_eggs_per_day}\,m_{\text{EFD}}(T)\right),
\qquad \mskey{gonotrophic.egg_batch_min}\le n_{\text{eggs}}\le\mskey{gonotrophic.egg_batch_max},
\end{equation}
with the thermal modifier
\begin{equation}\label{eq:efd}
\text{EFD}(T)=\mskey{thermal.efd_max_eggs_per_day}\,\frac{(T-\mskey{thermal.efd_tmin_c})(\mskey{thermal.efd_tmax_c}-T)}{(\mskey{thermal.efd_tmax_c}-\mskey{thermal.efd_tmin_c})^2/4},\qquad
m_{\text{EFD}}(T)=\frac{\text{EFD}(T)}{\text{EFD}(\mskey{thermal.efd_topt_c})},
\end{equation}
zero outside $[\mskey{thermal.efd_tmin_c},\mskey{thermal.efd_tmax_c}]$\,$^\circ$C and clipped to $[0,1]$. At
\mskey{thermal.efd_topt_c}$^\circ$C the mean batch is \mskey{gonotrophic.egg_batch_mean} eggs (binomial standard deviation
7.75). At the thermal thresholds the modifier is zero, so the binomial
draws zero and the batch is raised to the \mskey{gonotrophic.egg_batch_min}-egg floor: the floor, not
the thermal response, sets the minimum batch at extreme temperatures.

\begin{longtable}{@{}L{3.3cm}L{4.4cm}L{1.4cm}L{1.5cm}L{3.9cm}@{}}
\caption{Parameters of the gonotrophic submodel.}\label{tab:gonoparams}\\
\toprule
\textbf{Parameter} & \textbf{Meaning} & \textbf{Unit} & \textbf{Value} & \textbf{Provenance} \\
\midrule
\endfirsthead
\toprule
\textbf{Parameter} & \textbf{Meaning} & \textbf{Unit} & \textbf{Value} & \textbf{Provenance} \\
\midrule
\endhead
\bottomrule
\endlastfoot
\code{cycle_duration_days} & egg maturation, later cycles & d & \mskey{gonotrophic.cycle_duration_days} & engineering prior; rounded up to 2 \\
\code{first_cycle_days} & egg maturation, first cycle & d & \mskey{gonotrophic.first_cycle_days} & engineering prior; rounded up to 4 \\
\code{feeding_success_rate} & PROBING $\to$ BLOOD\_FED & -- & \mskey{gonotrophic.feeding_success_rate} & literature anchor; \emph{not consumed at run time} (Appendix~A.2) \\
\code{egg_batch_mean} & mean of the binomial draw & eggs & \mskey{gonotrophic.egg_batch_mean} & engineering prior \\
\code{egg_batch_n} & trials of the binomial draw & trials & \mskey{gonotrophic.egg_batch_n} & engineering prior \\
\code{egg_batch_min} & lower clamp & eggs & \mskey{gonotrophic.egg_batch_min} & engineering prior \\
\code{egg_batch_max} & upper clamp & eggs & \mskey{gonotrophic.egg_batch_max} & engineering prior \\
\code{resting_duration_days} & post-feed rest & d & \mskey{gonotrophic.resting_duration_days} & engineering prior; rounded down, one update \\
\code{efd_max_eggs_per_day} & EFD scale & eggs female$^{-1}$ d$^{-1}$ & \mskey{thermal.efd_max_eggs_per_day} & calibration scale, not absolute EFD \\
\code{efd_tmin_c}, \code{efd_tmax_c}, \code{efd_topt_c} & EFD thresholds and optimum & $^\circ$C & \mskey{thermal.efd_tmin_c}, \mskey{thermal.efd_tmax_c}, \mskey{thermal.efd_topt_c} & Mordecai et al.\ (2013) form \\
\end{longtable}

\textbf{Origin and rationale.} The gonotrophic cycle is the reproductive
clock of the model. The state machine makes the female's current
behaviour depend on her recent blood meal, which is the minimum needed to
couple biting, egg production, and habitat quality. The first cycle is
slower than later ones, reflecting the initial maturation of the
reproductive system. The batch size is a binomial draw rather than a
fixed mean so that individual females vary; the thermal modifier ties
fecundity to the same per-patch temperature that drives development and
survival. The clamp keeps a batch within a plausible range but, as noted,
it also imposes a \mskey{gonotrophic.egg_batch_min}-egg minimum when the thermal modifier is zero.

\subsection{Seek an oviposition site and oviposit}

A female in OVIPOSITION\_SEEKING looks for active patches within
\mskey{oviposition.detect_range_m}\,m and scores them by
\begin{equation}\label{eq:ovip}
A_p=H^{\text{tot}}_p\,e^{-d_p/\mskey{oviposition.scale_m}\,\text{m}},
\end{equation}
where $H^{\text{tot}}_p$ is the patch's static potential (the GeoPackage
amplitude, or the cell's $S$ when the patch is not in the GeoPackage).
The best-scoring patch is accepted only if $A_p>\mskey{oviposition.min_attraction}$; otherwise the
female has no target. If a target exists, she takes one step of
\mskey{oviposition.step_m}\,m toward it, snaps to a cell, and checks whether her new cell is an active
patch. If it is, she enters OVIPOSITING and lays; if not, she remains in
OVIPOSITION\_SEEKING and increments her search timer. If no target is
found, she stays in place and increments the timer. After
\mskey{oviposition.timeout_days} days without
laying, she abandons the batch, returns to HOST\_SEEKING, and does not
lay. When she lays, the eggs are added to the patch identifier she holds
at that moment, and the engine re-snaps her to the nearest \emph{active}
patch within \mskey{dispersal.snap_radius_m}\,m of her cell before laying; if no active patch is in
range, she keeps her previous patch association and the eggs go there
(\code{mosquito_submodel.cpp:704,726};
\code{mosquito_submodel.cpp:965-1006}). The patch that was originally
detected is therefore not necessarily the one that receives the eggs.

\begin{longtable}{@{}L{3.3cm}L{4.8cm}L{1.5cm}L{1.5cm}L{3.4cm}@{}}
\caption{Parameters of the oviposition submodel.}\label{tab:oviparams}\\
\toprule
\textbf{Parameter} & \textbf{Meaning} & \textbf{Unit} & \textbf{Value} & \textbf{Provenance} \\
\midrule
\endfirsthead
\toprule
\textbf{Parameter} & \textbf{Meaning} & \textbf{Unit} & \textbf{Value} & \textbf{Provenance} \\
\midrule
\endhead
\bottomrule
\endlastfoot
\code{scale_m} & plume decay scale & m & \mskey{oviposition.scale_m} & engineering prior; mechanism after Okal (2013) \\
\code{detect_range_m} & detection radius & m & \mskey{oviposition.detect_range_m} & engineering prior \\
\code{step_m} & daily step toward site & m & \mskey{oviposition.step_m} & engineering prior \\
\code{min_attraction} & acceptance threshold & -- & \mskey{oviposition.min_attraction} & engineering prior \\
\code{timeout_days} & skip-oviposition timeout & d & \mskey{oviposition.timeout_days} & engineering prior \\
\end{longtable}

\textbf{Origin and rationale.} Gravid females are not assumed to lay
randomly: they follow a decaying plume toward habitat. The \mskey{oviposition.scale_m}\,m decay
makes the choice strongly local; the \mskey{oviposition.detect_range_m}\,m detection radius lets a
female find a site across a few cells, and the \mskey{oviposition.step_m}\,m daily step means
that a distant site takes several days to reach. The timeout implements
skip-oviposition: a female without a suitable site resorbs or retains
eggs rather than laying in a hostile cell. The acceptance threshold
prevents a female from chasing a negligible potential.

\textbf{Consequence of the timeout.} A female walks at most \mskey{oviposition.step_m}\,m per day
for \mskey{oviposition.timeout_days} days before abandoning the batch, so the effective reachable radius
is about 350\,m, not the \mskey{oviposition.detect_range_m}\,m detection radius. A site detected between
350\,m and 2\,km can only be reached if the female is displaced by local
dispersal or wind while she is searching. The detection radius is
therefore an upper bound on perception, not on successful location.

\subsection{Seek hosts and feed}

Host seeking has two implementations, selected by the inputs. If a host
landscape is loaded and a mobility schedule is present, the engine runs
the \emph{phase-aware} path: for each of the four daily phases it builds
a phase-specific effective host landscape and lets every female still in
HOST\_SEEKING search and feed. If a host landscape is loaded but no
mobility schedule is present, it runs a single-phase residential path, in
which every HOST\_SEEKING female makes exactly one host-selection
attempt per day. If
no host landscape is loaded, it uses a deterministic fallback in which
every HOST\_SEEKING female takes a HUMAN blood meal at her current cell.

\textbf{Attraction.} For a female at cell $i$, the attraction of a host
cell $j$ within the species radius is
\begin{equation}\label{eq:attraction}
A_j=\left(\sum_k N_{jk}\pi_k\right)
\left(1+\mskey{host_seeking.indoor_attraction_add}\,I_j\right)\left(1+\mskey{host_seeking.urbanicity_add}\,U_j\right)
e^{-d_j/s_s},
\end{equation}
where $N_{jk}$ is the count of host type $k$ in cell $j$, $\pi_k$ the
species preference for that type, $I_j$ the cell's indoor fraction, $U_j$
its urbanicity (1 urban, 0 rural), $d_j$ the centre-to-centre distance,
and $s_s$ the species plume scale (\mskey{species.coluzzii.host_seeking_scale_m}\,m for \emph{An.\ coluzzii};
\mskey{species.arabiensis.host_seeking_scale_m}\,m for \emph{An.\ arabiensis}). Pigs are weighted as cattle and
chickens as the ``other'' (wildlife) class. The preference row of
\emph{An.\ coluzzii} is human \mskey{species.coluzzii.pref.human}, cattle \mskey{species.coluzzii.pref.cattle}, goat \mskey{species.coluzzii.pref.goat}, sheep \mskey{species.coluzzii.pref.sheep},
wildlife \mskey{species.coluzzii.pref.wildlife}; the other species rows are in the constants file.

\textbf{Effective hosts.} With mobility data, the residential host counts
are redistributed by a row-stochastic origin--destination matrix,
\begin{equation}\label{eq:heff}
H^{\text{eff}}_j(d,\phi)=\sum_i P_{ij}(\phi)\,H^{\text{res}}_i,
\end{equation}
realized stochastically each day and phase by a multinomial flow per
origin cell (with stochastic rounding of fractional populations). Humans
and cattle are moved; goats, sheep, pigs, and chickens remain at their
residential cells. The human day matrix is used for DAY, EVENING, and
DAWN; the night matrix for NIGHT. The livestock day matrix is used for
all phases except NIGHT, where the identity is used. When no mobility
data are present, the effective grid is the residential grid.

\textbf{Phase pass.} For each phase whose species activity weight exceeds
0.001 (with the default row, all four phases qualify), the engine samples
the human and livestock effective grids once and shares them across all
females. For each female still in HOST\_SEEKING, in SoA order, it:
(i) finds the best host cell within \mskey{host_seeking.detect_range_m}\,m whose attraction exceeds
0.5 and takes one step of \mskey{dispersal.host_flight_step_m}\,m toward it (if no such cell, she does not
move); (ii) recomputes the attraction field within
\mskey{species.coluzzii.host_seeking_radius_m}\,m at her new
position; (iii) selects one cell by a roulette proportional to
attraction; (iv) if a cell is selected, feeds with certainty, sets
BLOOD\_FED, records the meal at the selected host cell, and uses the
dominant host type of that cell as the host type. A female that feeds in
an early phase is skipped by later phases. A female that does not feed
remains in HOST\_SEEKING and may try again in the next phase. No random
number is consumed when the attraction field is empty: the selection
returns nothing, the attempt is still counted, and the female stays in
HOST\_SEEKING. The same holds when the total attraction is zero, in which
case the first cell of the field is taken without a draw; that branch is
unreachable in practice, because a cell whose attraction is zero is
dropped when the field is built.

\textbf{Starvation.} The HOST\_SEEKING timer is incremented once per daily
update. A female that reaches \mskey{host_seeking.starvation_days} days without a blood meal dies. The timer
is reset when she feeds and when she returns from oviposition.

\textbf{Transmission hazard on a bite.} If the female is SUSCEPTIBLE, the
host is HUMAN, and transmission is enabled, the bite infects her with
probability
\begin{equation}\label{eq:pinf}
P_{\text{inf}}=\beta_{hv}\,p^{\text{prev}}_j,
\end{equation}
where $j$ is the host cell and $p^{\text{prev}}_j$ is its infectious
fraction $I_H(j)/H_j$ \emph{as of the end of the previous day}: the human
grid keeps one daily prevalence field, written at the end of each SEIR
update, and the bite reads the value left there by the last update
(\code{mosquito_submodel.cpp:105,111};
\code{human_state.hpp:106};\code{human_state.cpp:330}). The hazard is
therefore lagged by one day relative to the bites recorded on the same
day. When the stored value is zero the code draws no random number at
all. On success she becomes EXPOSED with EIP
progress 0.

\begin{longtable}{@{}L{5.2cm}L{3.3cm}L{1.2cm}L{1.3cm}L{3.3cm}@{}}
\caption{Parameters of the host-seeking submodel.}\label{tab:hostparams}\\
\toprule
\textbf{Parameter} & \textbf{Meaning} & \textbf{Unit} & \textbf{Value} & \textbf{Provenance} \\
\midrule
\endfirsthead
\toprule
\textbf{Parameter} & \textbf{Meaning} & \textbf{Unit} & \textbf{Value} & \textbf{Provenance} \\
\midrule
\endhead
\bottomrule
\endlastfoot
\bcode{host_seeking.starvation_days} & maximum days without a host & d & \mskey{host_seeking.starvation_days} & engineering prior \\
\bcode{host_seeking.indoor_attraction_add} & indoor modifier coefficient & -- & \mskey{host_seeking.indoor_attraction_add} & engineering prior; no recorded source \\
\bcode{host_seeking.urbanicity_add} & urban modifier coefficient & -- & \mskey{host_seeking.urbanicity_add} & engineering prior \\
\bcode{host_seeking.indoor_fraction_urban} & indoor fraction, urban class & -- & \mskey{host_seeking.indoor_fraction_urban} & literature anchor for \emph{An.\ gambiae} endophily \\
\bcode{host_seeking.indoor_fraction_rural} & indoor fraction, rural class & -- & \mskey{host_seeking.indoor_fraction_rural} & literature anchor \\
\bcode{host_seeking.detect_range_m} & directed-movement detection range & m & \mskey{host_seeking.detect_range_m} & engineering prior \\
\bcode{host_seeking.approach_step_m} & HOST\_APPROACH step & m & \mskey{host_seeking.approach_step_m} & engineering prior; path not entered \\
\bcode{species.<row>.host_seeking_scale_m} & plume decay $s_s$ & m & \mskey{species.coluzzii.host_seeking_scale_m} (\mskey{species.arabiensis.host_seeking_scale_m} for \emph{arabiensis}) & engineering prior; mechanism recorded as Giraldo (2023) \\
\bcode{species.<row>.host_seeking_radius_m} & attraction-field radius & m & \mskey{species.coluzzii.host_seeking_radius_m} & engineering prior; no local plume measurement \\
\bcode{dispersal.host_flight_step_m} & directed flight step per phase & m & \mskey{dispersal.host_flight_step_m} & engineering prior \\
\bcode{species.<row>.activity_weights} & phase activity weights [DAY, EVENING, NIGHT, DAWN] & -- & \mskey{species.coluzzii.activity_weights} & \emph{An.\ coluzzii} row; engineering prior with a nocturnal shape \\
\bcode{kPhaseActivityEpsilon} & phase gate & -- & 0.001 & model definition \\
\bcode{species.<row>.pref} & host preference $\pi_k$ & -- & \mskey{species.coluzzii.pref.human}, \mskey{species.coluzzii.pref.cattle}, \mskey{species.coluzzii.pref.goat}, \mskey{species.coluzzii.pref.sheep}, \mskey{species.coluzzii.pref.wildlife} & \emph{An.\ coluzzii} row; model attraction weights, not observed HBI \\
\end{longtable}

\textbf{Origin and rationale.} The attraction field is a minimal
representation of host odour and CO$_2$ plumes: host abundance, species
preference, indoor/urban concentration, and exponential distance decay.
It is not a plume-transport model; the decay scale is a prior, not a
measurement. The phase pass exists because human and livestock hosts move
during the day, so the same mosquito experiences a different host
landscape at night and at midday. The multinomial flow preserves the
total host count and introduces realistic day-to-day variation in who is
where. The fallback without a host landscape is deliberately crude: it
guarantees a blood meal at the current cell so that a run without host
data still produces a gonotrophic cycle and eggs. It should not be used
for a quantitative host-partitioning analysis.

\textbf{Consequences worth stating.} With a \mskey{dispersal.host_flight_step_m}\,m step in each of four
active phases, a female can move up to 2\,km per day toward hosts,
independently of the local dispersal kernel; the phase weights gate
activity but are not used as probabilities, so the daytime weight of
\mskey{species.coluzzii.activity_weights} still produces a full daytime search-and-feed opportunity (Appendix~A.4).
Because the host type is the \emph{dominant} host of the selected cell,
and the selected cell is chosen by total attraction, the realized host
partition is not a direct draw from $\pi_k$; in a mixed cell, the most
attractive host type dominates the recorded bites.

\subsection{Disperse adults}

Every adult makes a local dispersal draw once per day. With probability
\mskey{dispersal.adult_prob_day} it moves by an isotropic Gaussian offset,
\begin{equation}\label{eq:disp}
(\Delta x,\Delta y)\sim\mathcal{N}(0,\sigma^2),\qquad
\sigma=\mskey{dispersal.adult_sigma_m}\,\text{m},
\end{equation}
clipped to a circle of radius \mskey{dispersal.adult_max_m}\,m and converted to degrees with
the cosine of the local latitude. Otherwise it stays. The draw is
independent of sex, age, and gonotrophic state.

If a wind field is supplied, gravid females and females in
OVIPOSITION\_SEEKING may additionally migrate by wind. The wind path is
active only in the months July--October and December--March, and only at
the flight hour (\mskey{dispersal.flight_hour}:00). For each eligible female, with probability \mskey{dispersal.wind_migration_prob},
the model reads the wind vector at her position and the current day/hour
and displaces her downwind by
\begin{equation}\label{eq:wind}
\Delta=\left(|\mathbf{u}|+\mskey{dispersal.wind_self_speed_ms}\,\text{m s}^{-1}\right)\cdot \mskey{dispersal.wind_flight_hours}\,\text{h},
\end{equation}
where $|\mathbf{u}|$ is the wind speed and the \mskey{dispersal.wind_self_speed_ms}\,m\,s$^{-1}$ is the
female's own flight speed. She survives the flight with probability \mskey{dispersal.wind_survival};
on survival she is placed at the destination and reset to
OVIPOSITION\_SEEKING; on failure she is removed. If the wind speed is at
most 0.1\,m\,s$^{-1}$, the displacement is zero but the survival draw and
the state reset still occur. Local dispersal is applied afterwards, so a
wind-migrated female can also make a local jump on the same day. Lon/lat
may leave the AOI during wind displacement; the cell index is clamped to
the AOI border at the next snap.

After dispersal, every adult is snapped to the nearest active patch
within \mskey{dispersal.snap_radius_m}\,m. If no active patch is within that radius, the previous
patch association is kept. The snap is skipped when the adult's
$(r,c)$ has not changed since the previous day (a cache in
\code{update_patch_id}), so a stationary adult keeps its previous
association even if that patch has since deactivated; the association is
recomputed when the adult changes cell. The patch association is what
band~1 of the state output counts; band~2 counts the current cell.

\begin{longtable}{@{}L{3.4cm}L{4.5cm}L{1.4cm}L{1.5cm}L{3.8cm}@{}}
\caption{Parameters of the dispersal submodel.}\label{tab:dispparams}\\
\toprule
\textbf{Parameter} & \textbf{Meaning} & \textbf{Unit} & \textbf{Value} & \textbf{Provenance} \\
\midrule
\endfirsthead
\toprule
\textbf{Parameter} & \textbf{Meaning} & \textbf{Unit} & \textbf{Value} & \textbf{Provenance} \\
\midrule
\endhead
\bottomrule
\endlastfoot
\code{adult_prob_day} & daily local dispersal probability & d$^{-1}$ & \mskey{dispersal.adult_prob_day} & engineering prior \\
\code{adult_sigma_m} & Gaussian scale $\sigma$ & m & \mskey{dispersal.adult_sigma_m} & prior inside the 350--650\,m\,d$^{-1}$ range recorded as Costantini (1996); not the midpoint \\
\code{adult_max_m} & clip radius & m & \mskey{dispersal.adult_max_m} & compromise between the p95 values recorded as Thomas (2013) \\
\code{wind_migration_prob} & daily wind-migration probability & d$^{-1}$ & \mskey{dispersal.wind_migration_prob} & engineering prior \\
\code{wind_flight_hours} & flight duration & h & \mskey{dispersal.wind_flight_hours} & engineering prior; mechanism recorded as Huestis (2019) and Faiman (2020) \\
\code{wind_survival} & survival during wind flight & -- & \mskey{dispersal.wind_survival} & engineering prior; mechanism recorded as Sanogo (2021) \\
\code{wind_self_speed_ms} & own flight speed & m\,s$^{-1}$ & \mskey{dispersal.wind_self_speed_ms} & engineering prior \\
\code{flight_hour} & hour used for the wind lookup & h & \mskey{dispersal.flight_hour} & model definition \\
\code{snap_radius_m} & patch-snap radius & m & \mskey{dispersal.snap_radius_m} & engineering prior \\
\code{night_start_hour}, \code{night_end_hour} & night window for wind flight & h & \mskey{dispersal.night_start_hour}, \mskey{dispersal.night_end_hour} & model definition \\
\end{longtable}

\textbf{Origin and rationale.} Local dispersal is a diffusion
approximation: most adults stay, a minority make a Gaussian jump, and the
tail is clipped to avoid implausible long jumps. The scale is within the
range of mark--release--recapture estimates for \emph{An.\ gambiae}
complex populations, but the value is a prior, not a local estimate. The
wind path represents the well-documented long-distance transport of
mosquitoes in the Sahel; it is optional because it requires a wind field
and is not part of the delivered Ghana product. The \mskey{dispersal.snap_radius_m}\,m snap reflects
the resolution of the habitat product: an adult between two patches is
associated with the nearest active one, not with the nearest cell.

\subsection{Survive adults}

Each adult draws a uniform random number and survives with probability
\begin{equation}\label{eq:survival}
p_d(T)=\mskey{adults.daily_mort_basal}\,\exp\left(-\frac{(T-\mskey{adults.opt_c})^2}{2\cdot\mskey{adults.sigma_c}^2}\right),
\end{equation}
clamped to $[\mskey{adults.mort_floor},\mskey{adults.mort_cap}]$. The temperature $T$ is that of
the adult's \emph{current cell}: the engine builds a cell-indexed
temperature grid from the patch states and looks it up at the adult's
$(r,c)$ (\code{mosquito_submodel.cpp:1104-1153}). An adult whose cell is
not in that table --- it is not a candidate patch --- gets the bare basal
value \mskey{adults.daily_mort_basal} with no Gaussian term, and the clamp
leaves that value unchanged. The survival draw is independent of sex,
age, gonotrophic
state, and infection state; there is no senescence, no insecticide, and
no host-defense mortality. The function is named
\code{daily_mort_basal} in the constants file, but the value \mskey{adults.daily_mort_basal} is a
daily \emph{survival} probability, as the file itself warns.

\begin{longtable}{@{}L{3.4cm}L{4.6cm}L{1.4cm}L{1.5cm}L{3.7cm}@{}}
\caption{Parameters of the adult-survival submodel.}\label{tab:survparams}\\
\toprule
\textbf{Parameter} & \textbf{Meaning} & \textbf{Unit} & \textbf{Value} & \textbf{Provenance} \\
\midrule
\endfirsthead
\toprule
\textbf{Parameter} & \textbf{Meaning} & \textbf{Unit} & \textbf{Value} & \textbf{Provenance} \\
\midrule
\endhead
\bottomrule
\endlastfoot
\code{daily_mort_basal} & basal daily survival & d$^{-1}$ & \mskey{adults.daily_mort_basal} & literature anchor recorded as Lardeux; see Appendix~A.10 \\
\code{opt_c} & thermal optimum & $^\circ$C & \mskey{adults.opt_c} & literature anchor \\
\code{sigma_c} & thermal width & $^\circ$C & \mskey{adults.sigma_c} & literature anchor \\
\code{mort_cap} & upper clamp of $p_d$ & -- & \mskey{adults.mort_cap} & $\le$ basal, so dormant for the current value \\
\code{mort_floor} & lower clamp of $p_d$ & -- & \mskey{adults.mort_floor} & binds only outside about 12.1--37.9$^\circ$C \\
\end{longtable}

\textbf{Origin and rationale.} The Gaussian form encodes a thermal
optimum with a broad tolerance: survival is highest at 25$^\circ$C
and declines symmetrically away from it. The basal value implies a mean
adult lifetime of $-1/\ln(\mskey{adults.daily_mort_basal})=7.2$\,d at the optimum; at the lower end
of the Ghana temperature field (about 21$^\circ$C) it is 0.84
(5.8\,d), and at the upper end (about 36.9$^\circ$C) it is 0.64
(2.2\,d). Because the Gaussian factor is at most 1 and the basal value is
0.87, the \mskey{adults.mort_cap} cap can never bind; the \mskey{adults.mort_floor} floor binds only below about 12.1$^\circ$C or above about 37.9$^\circ$C, which the delivered
Ghana daily means do not reach. The cap is therefore dormant in the
current calibration, and the floor is dormant in Ghana 2024.

\subsection{Transmit parasites}

The transmission layer is a spatial Ross--Macdonald SEIR--SEI model. It
is active only when \code{--enable-transmission} is set.

\textbf{Vector EIP.} For each adult female in state EXPOSED, the model
adds the day's degree-day increment at her current cell,
\begin{equation}\label{eq:eip}
E\leftarrow E+\max(0,T-\mskey{eip.base_c}),
\end{equation}
and changes her to INFECTIOUS when $E\ge\mskey{eip.threshold_gd}$ degree-days. At
25$^\circ$C this takes 12.2 days; at 30$^\circ$C, 7.9 days; at
21$^\circ$C, 22 days. The EIP block runs after the mosquito day, so a
female infected by today's bite accrues today's degree-days.

\textbf{Vector infection.} A susceptible female that takes a blood meal
on a human is infected with probability given by
Equation~\eqref{eq:pinf}; the prevalence read at the host cell is the
one of the previous day. On
success she becomes EXPOSED with $E=0$.

\textbf{Human force of infection.} The bite ledger aggregates, per cell,
the infectious bites delivered to humans during the day. The force of
infection on the human population of cell $j$ is
\begin{equation}\label{eq:foi}
\lambda_j=1-\exp\left(-\mskey{transmission.beta_vh}\frac{B^{\text{inf}}_j}{H_j}\right),
\end{equation}
and the number of new exposures is
\begin{equation}\label{eq:newexp}
E^{\text{new}}_j\sim\text{Binomial}\left(\operatorname{round}(S_j),\lambda_j\right),
\end{equation}
with a deterministic fallback $E^{\text{new}}_j=S_j\lambda_j$ when
$\operatorname{round}(S_j)=0$, and
clamped to the susceptible count. The new exposures enter the exposed
cohort queue at age zero. The queue is shifted by one day per update:
a cohort exposed on day $d$ becomes infectious on day $d+\mskey{transmission.human_incubation_days}$ and recovers
on day $d+\mskey{transmission.human_infectious_days}$ after onset. With immunity disabled (the default), recovery
returns the cohort directly to $S_H$ (an SIS-like human model) and $R_H$
stays zero; with immunity enabled, recovered cohorts spend \mskey{transmission.human_immunity_days} days in
$R_H$ before returning to $S_H$.

\textbf{Outbreak seeding.} The seeding rule is described in
Section~5.3. Two timing details matter. First, the outbreak is checked
after the mosquito day, so the seeded infectious humans are present only
after that day's bites; they can infect vectors from the next day onward.
Second, the check is retried daily until a qualifying cell exists, so a
run can spend many days with zero human infections while the vector
population is still too sparse or too far from a settlement.

\textbf{Observation.} The transmission COG (optional) writes human
prevalence, per-person incidence, infectious vector pressure (infectious
bites per cell per day), and a binary active-focus band. The focus band is
1 when prevalence is at or above 0.01 or incidence is positive; the
threshold is an output classification, not a model rate.

\begin{longtable}{@{}L{3.4cm}L{4.5cm}L{1.4cm}L{1.5cm}L{3.8cm}@{}}
\caption{Parameters of the transmission submodel.}\label{tab:txparams}\\
\toprule
\textbf{Parameter} & \textbf{Meaning} & \textbf{Unit} & \textbf{Value} & \textbf{Provenance} \\
\midrule
\endfirsthead
\toprule
\textbf{Parameter} & \textbf{Meaning} & \textbf{Unit} & \textbf{Value} & \textbf{Provenance} \\
\midrule
\endhead
\bottomrule
\endlastfoot
\bcode{beta_hv} & human-to-vector probability per bite on an infectious human & -- & \mskey{transmission.beta_hv} & literature anchor; not locally fitted \\
\bcode{beta_vh} & vector-to-human probability per infectious bite & -- & \mskey{transmission.beta_vh} & literature anchor; not locally fitted \\
\bcode{human_incubation_days} & $E_H\to I_H$ & d & \mskey{transmission.human_incubation_days} & literature anchor \\
\bcode{human_infectious_days} & $I_H\to R_H$ (or $S_H$) & d & \mskey{transmission.human_infectious_days} & literature anchor \\
\bcode{human_immunity_days} & $R_H\to S_H$ when immunity is on & d & \mskey{transmission.human_immunity_days} & literature anchor; immunity off by default \\
\bcode{initial_human_prevalence} & initial infectious fraction, uniform-legacy only & -- & \mskey{transmission.initial_human_prevalence} & scenario default \\
\bcode{initial_vector_infected_frac} & initial infectious fraction of adult females & -- & \mskey{transmission.initial_vector_infected_frac} & scenario default \\
\bcode{eip.base_c} & EIP base temperature & $^\circ$C & \mskey{eip.base_c} & degree-day formulation \\
\bcode{eip.threshold_gd} & EIP threshold & degree-days & \mskey{eip.threshold_gd} & degree-day formulation \\
\bcode{human_outbreak_day} & first day the outbreak may seed & d & 0 & scenario default \\
\bcode{human_outbreak_foci} & number of foci & -- & 3 & scenario default \\
\bcode{human_outbreak_cases} & cases per focus (total in cluster mode) & persons & 50 & scenario default \\
\bcode{human_min_cell_pop} & minimum cell population for a focus & persons & 50 & scenario default \\
\bcode{human_outbreak_min_density} & adaptive density gate & adults\,cell$^{-1}$ & 0 (any positive) & scenario default \\
\bcode{human_cluster_radius_km} & cluster radius; 0 = independent foci & km & 0 & scenario default \\
\bcode{human_cluster_max_core_pop} & core population cap; 0 = none & persons & 0 & scenario default \\
\bcode{focus_threshold} & prevalence threshold of the active-focus band & -- & 0.01 & output classification \\
\end{longtable}

\textbf{Origin and rationale.} The model is the standard
Ross--Macdonald structure, spatialized and driven by explicit bites
rather than by a mean biting rate. The advantage is that every infection
event is traceable to a cell and a host type; the cost is that the
transmission parameters are per bite, and the bite process is itself
driven by the host-seeking submodel. The human durations are fixed, not
exponential, so the model avoids the dwell-time distortion of a Markov
SEIR. The default is an SIS-like human model: without immunity, a
recovered human is immediately susceptible again, which is a strong
assumption for a population with repeated exposure. The \mskey{transmission.initial_human_prevalence} initial
prevalence is used only in the legacy uniform mode; the default
random-viable mode starts with no human infections and seeds a focal
outbreak when vectors are present.

% =====================================================================
\appendix

\section{Implementation notes for replicators}

The notes below record behaviours of the current code that are not
apparent from the protocol-level description and that matter for exact
replication or for interpreting a run. They are factual descriptions of
the source tree, not recommendations. File and function names refer to
the engine under \code{mal-core/src/mal_core/abm/}.

\begin{enumerate}
\item \textbf{The habitat update uses the species-neutral affinity row.}
  \bcode{CoordinatorModel::update_habitat_day()} and
  \code{to_dataframe()} call \code{default_species_weights()} with no
  argument, which returns the neutral row (surface \mskey{species.neutral.weights.surface_persistent}, wetland \mskey{species.neutral.weights.wetland_persistent},
  urban \mskey{species.neutral.weights.urban_persistent}, wetland pluvial \mskey{species.neutral.weights.wetland_pluvial}, open \mskey{species.neutral.weights.open_pluvial}, crop \mskey{species.neutral.weights.crop_pluvial}), not the active
  species' row. The per-species \code{species.<name>.weights} rows are
  loaded and exposed by \bcode{default_species_weights(SpeciesId)} but
  are not consumed by the habitat update. Consequently $S$ and $K$ are
  species-neutral in the current release; the species affects salinity,
  host preference, phase activity, and plume scale, but not the habitat
  offer.

\item \textbf{Feeding success is not stochastic, and two states are
  unreachable.} \bcode{gonotrophic.feeding_success_rate} (\mskey{gonotrophic.feeding_success_rate}) is
  loaded and stored in \code{GonotrophicParams}, but no code path draws
  it. A female that finds a host feeds with certainty. The states
  HOST\_APPROACH and PROBING are defined and handled in the state
  function, but the feeding code goes directly from HOST\_SEEKING to
  BLOOD\_FED, so the HOST\_APPROACH block and the PROBING case are never
  entered.

\item \textbf{There are no hourly host-seeking sub-steps.} The header
  \bcode{multirate_scheduler.hpp} describes 12 nightly hourly sub-steps,
  and \code{MosquitoSubmodel} holds a \code{MultirateDayState} member, but
  that member is never advanced. Host seeking is a single daily loop plus
  the four-phase pass described in Section~7.6.

\item \textbf{The phase activity weights gate, but do not weight, the
  phase pass.} \code{phase_active()} returns true when the species
  activity weight exceeds 0.001. With the default weights
  (\mskey{species.coluzzii.activity_weights}) all four phases are active, and each phase
  gives a full search-and-feed opportunity. The weights are not used as
  probabilities, so the model has substantial daytime host seeking
  despite the nocturnal shape of the row.

\item \textbf{The three aquatic mortality streams are fixed and
  re-seeded daily.} \code{washout()}, \code{stage_mortality()}, and
  \bcode{larva_mortality_density()} each construct a fresh
  \code{std::mt19937_64} with a constant seed (\mskey{rng.aquatic_seed_washout}, \mskey{rng.aquatic_seed_stage_mortality}, \mskey{rng.aquatic_seed_larva_mortality}) on every
  call. Their draws depend on cohort order and counts but not on the
  rollout seed, and the same random sequence is reused every day. This is
  a deliberate reproducibility choice, but it means that two rollouts
  differ only in adult, host-selection, and transmission draws, not in
  the aquatic mortality draws.

\item \textbf{The GeoPackage \code{dynamics} column is a string, but the
  loader reads an integer.} The delivered Ghana file stores
  \code{dynamics} as a string (the habitat type's dynamic label), while
  \bcode{HabitatEngine::load_from_gpkg()} calls \code{GetIntField()} on
  it. The field is therefore 0 for every patch. It is diagnostic only:
  the coordinator's persistent/pluvial split comes from
  \code{habitat_types.dynamics} in the constants file, not from the
  patch record.

\item \textbf{The sidecar's \code{abm_params_hash} is a placeholder.}
  The state sidecar writes \code{"sha256:pending"} for this key; it is
  not a computed hash of the parameters. A consumer that relies on it to
  prove parameter identity will be misled.

\item \textbf{The cohort log's \code{eip_frac} is not an EIP fraction.}
  In \bcode{MosquitoSubmodel::advance_day()}, the field is set to
  \code{n_adults / n_alive} after the day's mortality, i.e.\ the adult
  fraction of the live population (which is always 1 when no immatures
  are in the SoA). The true EIP state is in \code{vector_state} and
  \bcode{parasite_eip_progress}, and the transmission log reports the
  vector compartments correctly.

\item \textbf{\code{k_ref} is annotated with a per-day unit but is
  consumed as a standing capacity.} The constants file labels
  \code{species.k_ref} as ind\,km$^{-2}$\,d$^{-1}$, but
  Equation~\eqref{eq:K} multiplies it by area and $S$ and the result is
  used directly as $K$ in the Beverton--Holt expression, where it must be
  a count. The implemented meaning is individuals per cell; no per-day
  dimension cancels. The unit annotation should be read as a provenance
  note, not as a dimensional statement.

\item \textbf{\code{daily_mort_basal} is a survival probability.} The
  key name says ``mortality'', but \mskey{adults.daily_mort_basal} is the daily probability of
  surviving, as the constants file's own comment warns. Equation
  \eqref{eq:survival} uses it as a survival multiplier.

\item \textbf{The round-robin submodel constructor is not used by the
  engine.} That overload (an integer patch count plus two integer counts)
  seeds every adult at row 0, column 0, and lon/lat 0, and seeds
  immatures as five equal egg cohorts. The engine always uses the
  instruction constructor, which carries the cell centre and the staged
  immature distribution. The round-robin path survives for tests and
  legacy callers; a run through the engine never uses it.

\item \textbf{Mobility moves hosts, not infections.} The phase-specific
  effective host grids redistribute humans and cattle for attraction and
  biting, but the human SEIR compartments remain at the residential
  cells. A bite recorded at a destination cell therefore applies its
  force of infection to the residential population of that cell. This is
  a structural assumption, not a bug in the queue update, but it means
  that the model does not represent the infection of a commuter at the
  commuter's home cell.

\item \textbf{\code{--max-population} is an operational abort, not a
  biological ceiling.} When set to a positive value, the engine throws and
  aborts the run if the live adult count exceeds it. It does not reserve
  memory, cap births, or alter the dynamics.

\item \textbf{Open water and open sea are not separated by the habitat
  planes.} The \code{lc_open_plus_water} plane includes Dynamic World's
  water class (label 0), and the ingest saltwater filter zeroes only the
  JRC persistent plane. An open-sea cell can therefore be a candidate
  patch with $H_{\text{open}}=1$, $S\approx0.40$ at $z=1$, and $K>0$. Its
  air temperature is the 0.0 sentinel (a finite value, so
  \code{impute_temp_nans} does not touch it), so aquatic development is
  zero and adults there survive at the \mskey{adults.mort_floor} floor. The \code{land_frac}
  plane, when present, is used by the diagnostic GIF basemap but is not
  part of the habitat feature bundle and does not gate the habitat update.
  In the delivered Ghana product a Gulf of Guinea cell (row 745, column
  275) has \code{jrc_persistent}=0, \code{lc_open_plus_water}=1,
  $H_{\text{total}}=1$, and a GeoPackage feature with potential 1.0,
  confirming its membership in the candidate set.
\end{enumerate}

% =====================================================================
\section{Provenance, versioning, and testing status}

\subsection*{B.1 Constants and contracts}

All model constants are in \bcode{mal-core/config/model_constants.yaml}
(schema version \mskey{version}). The file is parsed by a strict loader: every key is
mandatory, unknown keys are an error, wrong types are an error, and there
are no compiled defaults. The Python side reads the same file through a
validated loader, and a parity test compares the Python dump of every key
with the C++ dump of the same file, so a key added on one side without the
other fails loudly \emph{for the constants file}. The state output
contract is version \mskey{contracts.cogs_state}
(the two bands are raw adult counts, the \code{k_max} key is retired,
and the band names are \bcode{adult_patch_count} and
\code{adult_cell_count}); the environmental NetCDF contract is version
\mskey{contracts.env_nc} (daily rainfall and air temperature, five static habitat planes,
optional salinity). The two contracts are versioned independently. A
change to a frozen key requires a version bump and a contract review.

Two divergences between the Python habitat mirror and the engine must be
recorded, because they break the assumption that one implementation
mirrors the other. First, \bcode{mal_core/habitat/capacity.py:30}
hardcodes \code{K\_REF\_DIAG\_PER\_KM2} $=\codeconst{ingest.k_ref_diag_per_km2}$
as
the capacity scale, where the engine uses \code{species.k\_ref}
$=\mskey{species.k_ref}$ (\code{habitat_availability.hpp:191-193}); the two
are not the same quantity, and the Python one is not read from the
constants file at all. Second, \code{capacity.species_accuracy} defaults
to the \emph{coluzzii} weight row, while the engine's habitat update uses
the \emph{neutral} row (\code{coordinator.cpp:117}). The Python habitat
module is therefore a diagnostic implementation, not a parity reference
for the engine's behaviour.

\subsection*{B.2 Software testing}

The engine has C++ GoogleTest suites for climate reading, the
coordinator, dispersal, EIP, the engine facade, the env reader, the
gonotrophic cycle, and the effective host landscape; the Python side has
tests for the wrapper, the runner, the output contract, and the habitat
contract, plus golden fixtures that pin the per-day outputs of a small
synthetic scenario. The Python--C++ parity that does exist compares the
two loaders of the constants file. There is \emph{no} parity test of the
habitat implementation: the Python habitat tests
(\bcode{mal-core/tests/test_habitat_python_contract.py}) are
self-contained, they pin the Python diagnostic capacity constant and the
coluzzii weight row (\bcode{...:181-208}), and they never compare the
Python output with the engine's; the C++ counterpart named in that
module's docstring (\code{test_habitat_availability.cpp}) is not in the
tree. These tests establish that the code does
what this document says on the tested inputs. They do not establish that
the parameter values are biologically correct for a particular district.

\subsection*{B.3 Calibration and validation status}

The current parameter set is a mixture of literature values and explicit
engineering priors. The project states that no parameter has been fitted
to Ghanaian adult abundance, trap catches, sporozoite rates, or
incidence in this release. The code records no fit provenance for any
parameter, so this is a project statement that the code can neither
confirm nor refute (Appendix~B.4).
The model is used as a synthetic ground truth and as a scenario
laboratory; its outputs should be read as conditional consequences of the
stated rules and inputs, not as predictions of observed malaria. The
patterns listed in Element~1 are the intended calibration targets for a
future cycle. Until that cycle is done, the most defensible use of the
model is comparative: how do outputs change when an input or a prior is
changed, and which mechanisms are responsible?

% =====================================================================
\subsection*{B.4 Values that are readings of the tree}

Every model constant quoted in this document comes from the generated
macros, that is, from \code{mal-core/config/model_constants.yaml}
through \code{docs/concepts/generated/macros.tex}. The magnitudes below
have no key in that file. They are readings of the code or of a delivered
artefact, and they are listed here so that a reader can tell them apart
from the generated values.

\begin{itemize}
\item \textbf{Grid geometry of the delivered product.} 551 columns by
  779 rows, 429{,}229 cells, a bounding box of about 428{,}600\,km$^2$
  and about 0.99\,km$^2$ per cell at 8$^{\circ}$N: read from the shape
  and the geolocation of
  \code{data/ghana/ghana_regional_2024_2024_env.nc}.
\item \textbf{Habitat patch counts.} 382{,}127 point features in
  \code{data/ghana/ghana_habitat_patches.gpkg}; 368{,}761 of them
  with \code{habitat_potential} at or above 0.01; and 370{,}061 grid
  cells whose static total potential, computed with the engine's five
  non-zero slots, is at or above 0.01. The first two are read from the
  GeoPackage, the third is computed from the static planes of the
  environmental NetCDF.
\item \textbf{Population and livestock of the delivered grid.}
  32{,}048{,}206 humans, 3{,}457{,}920 cattle, 14{,}807{,}556 goats,
  8{,}591{,}670 sheep, 2{,}691{,}326 pigs and 121{,}353{,}624 chickens:
  read from \code{data/ghana/}.
\item \textbf{Habitat characterization of the delivered landscape.} The
  11.0\% and 82.9\% of cells above the activation floor at $z=0$ and
  $z=1$, and the mean capacities of 0.61 and 10.4 individuals per cell
  with a maximum of about 32: computed from the static planes of the
  environmental NetCDF with the equations of Section~7.1. They are not
  the output of an engine run.
\item \textbf{Engine constants with no key.} The Bernoulli(0.5) sex
  ratio at emergence; the seeding counts (50 adults and 30 immatures
  per selected patch, 3 foci of 50 cases, a minimum cell population of
  50, a focus prevalence threshold of 0.01); the stage-allocation
  fractions of the seeded immatures (0.05, 0.15, 0.25, 0.20) and the
  0.05 density term of the cluster seeding rule; the 0.5 attraction gate of
  the directed host flight; the phase gate
  \code{kPhaseActivityEpsilon} $=0.001$; the 0.1\,m\,s$^{-1}$ wind-speed
  threshold below which a wind displacement is zero; the 80\% occurrence
  threshold of the JRC product; the accepted range of 1 to 731 daily
  steps and the 366-day delivered horizon; the 5\,km radii of the
  explicit and host-weighted seeding modes; and the clock windows of the
  four activity phases.
\item \textbf{Statements the code cannot settle.} Whether a formal
  calibration to Ghanaian field data has been completed, and whether any
  parameter has been fitted to field observations. No fit provenance is
  recorded anywhere in the code, so these remain project statements
  (Element~1 and Section~B.3).
\end{itemize}

% =====================================================================
\section{Scenario harness}

The individualisation and movement redesign is verified against synthetic
scenarios, E00 to E18, run on a 10$\times$10 synthetic area of interest
through the same manifest route as a production run. This appendix is
generated from the scenario files and their catalogue
(\code{mal-core/src/mal_core/abm/tests/scenarios/}); every number a scenario
fixes is read from its JSON, and a scenario without a catalogue entry fails
the documentation build.

A scenario passing is easy to misread, so the test suite
(\code{mal-core/src/mal_core/abm/tests/test_scenarios.py}) separates four
kinds of check, and each entry below says which of them run today:

\begin{itemize}
\item \textbf{Seeding fidelity} (all scenarios): what the engine holds after
  reading the scenario matches the file. Adults obey the conservation
  identity (adults today equal adults yesterday minus deaths plus emerged),
  the raster output and the cohort log agree, and seeded pupae emerge.
\item \textbf{Baseline pins}: facts about the chain as it is now, recorded so
  that a change in behaviour is noticed. They are known failures, not
  expectations.
\item \textbf{Target gates}: the behaviour the design pre-registers (for
  example no bite at 200 m). They are written as strict expected failures:
  they fail today, and the day one passes the test turns red and must be
  promoted to a plain assertion.
\item \textbf{Determinism}: the same seed gives byte-identical artefacts and
  another seed does not.
\end{itemize}

A run that exits normally proves nothing about behaviour, and no check relies
on that alone. Scenarios whose behaviour is implemented by a later step run
only the seeding checks until that step lands; the field
\emph{Runs a behaviour test today} says so.

% Generated by tools/gen_concepts_values.py -- do not edit.
\subsection*{E00 --- Minimum harness case: one seeded female}
\begin{tabular}{@{}ll@{}}
\textit{Mode} & gate \\
\textit{Seed / days} & 20261117 / 1 \\
\textit{Mosquitoes seeded} & 1 female, 0 male \\
\textit{Closes} & Step 0 \\
\textit{Runs a behaviour test today} & full \\
\end{tabular}

\par\smallskip\noindent\textbf{Purpose.} Checks the scenario harness itself with the smallest possible world: one seeded female adult, one day, and no biology. It answers whether a seeded individual is placed, counted and kept by the chain before any behaviour is tested.
\par\smallskip\noindent\textbf{Setup.} One host-seeking female adult of age zero, seeded at a known position, with no humans, no aquatic habitat and no livestock. The run lasts one day. The scenario is the first step of the harness, so it also checks that the manifest declares exactly one scenario key.
\par\smallskip\noindent\textbf{What passing or failing implies.} If the seeded female is not where it was seeded, or the day-0 count differs from the seeded count, the harness is not reproducing its inputs and no later result can be trusted. A pass says only that seeding and the manifest route are faithful; it says nothing about movement or biology.
\par\smallskip\noindent\textbf{Expected behaviour.} Position, sex and state after one day match the seeded values. These are verifiable by construction, not from a biological band. The manifest declares exactly one scenario key, scenario:E00. This is a gate that closes Step 0 (harness and manifest route), with no tolerance band.
\par\smallskip\noindent\textbf{What the current chain does.} The current chain keeps the single adult in its seeded cell for the whole run and records one successful human bite on day 1 and none afterwards; the bite is drawn against the cell-aggregated human population (about 127,000 people), not against the seeded person
\par\smallskip\noindent\textbf{Checked today by the test suite.}
\begin{itemize}
\item Seeding fidelity: the adult population obeys the conservation identity on every day (adults today = adults yesterday - deaths + emerged, starting from the seeded count), so no seeded individual is dropped or invented; the cohort log and the raster output agree on the number of adults
\item The run completes through the manifest route ('validate\_completeness', resolve, binary) with exit code 0
\item Position pin: the agent remains in the seeded cell
\item GIF baseline pin: the transmission GIF of the one-day run has exactly one frame (a single slide); this is the visualiser state to be improved, not a result
\end{itemize}
\par\noindent\textbf{What the GIF must show.} Today the transmission GIF has one frame per simulated day and no scenario overlay. For E00 the informative version would show the single female marker at its seeded position, with the cell boundary drawn in every frame and a day label. A marker that stays inside its seeded cell confirms the position pin by eye; any jump out of the cell would be visible at once.

\subsection*{E01 --- Person beyond the activation radius}
\begin{tabular}{@{}ll@{}}
\textit{Mode} & gate \\
\textit{Seed / days} & 20261117 / 3 \\
\textit{Mosquitoes seeded} & 1 female, 0 male \\
\textit{People} & 1; distance from the first mosquito: 200 m \\
\textit{Closes} & Step 5 \\
\textit{Runs a behaviour test today} & seeding-only \\
\end{tabular}

\par\smallskip\noindent\textbf{Purpose.} Tests whether a host-seeking female with a person a short distance away, but outside her carbon dioxide activation radius, makes contact at all. It is the first check of the perception cascade and the baseline of the current chain, recorded at Step 1.
\par\smallskip\noindent\textbf{Setup.} One host-seeking female adult and one susceptible person placed just beyond the 100 m carbon dioxide activation radius of the female. There is no aquatic habitat, so the female has no breeding site to oviposit in. The run lasts three days.
\par\smallskip\noindent\textbf{What passing or failing implies.} A pass means the perception cascade respects its activation band: a person outside the activation radius is not contacted, and the chain does not bite people the sensory model should exclude. A fail means the chain still bites at distances the individual geometry rules out, so the transmission numbers would reflect the old cell-scale contact rule instead of individual distance. In the field a female does not find a host she cannot sense.
\par\smallskip\noindent\textbf{Expected behaviour.} Zero successful meals, because distance does not allow contact at this range. No oviposition, since there is no breeding site. The female starves by day three, which is the plan's baseline statement. The target is zero human bites, a strict expected failure until Step 5. The baseline pin records today's behaviour (one bite on day one) as a regression reference, not as an expectation.
\par\smallskip\noindent\textbf{What the current chain does.} The current chain keeps the single adult in its seeded cell for the whole run and records one successful human bite on day 1 and none afterwards; the bite is drawn against the cell-aggregated human population (about 127,000 people), not against the seeded person
\par\smallskip\noindent\textbf{Checked today by the test suite.}
\begin{itemize}
\item Seeding fidelity: the adult population obeys the conservation identity on every day (adults today = adults yesterday - deaths + emerged, starting from the seeded count), so no seeded individual is dropped or invented; the cohort log and the raster output agree on the number of adults
\item The run completes through the manifest route ('validate\_completeness', resolve, binary) with exit code 0
\item Baseline pin: the known-failure facts above are asserted, so any change in the chain's behaviour is flagged and must be acknowledged by updating the pin
\item Target gate (strict expected failure): the plan's target is asserted and is expected to fail until Step 5; when it starts passing the test fails and must be promoted
\item Target for E01: zero successful human bites over the run
\item GIF baseline pin: three simulated days give a three-frame GIF; the target of one frame per phase of the day is a strict expected failure until the visualiser is changed
\end{itemize}
\par\noindent\textbf{What the GIF must show.} Today the transmission GIF has one frame per simulated day and no scenario overlay, so the day-one bite is not visible. For E01 show the female marker, the person marker, and a circle of the 100 m activation radius around the female, redrawn in frames per phase of each night. A correct run shows the female wandering without ever entering the circle and no bite marker at any frame.

\subsection*{E02 --- Night trajectory versus net displacement}
\begin{tabular}{@{}ll@{}}
\textit{Mode} & report \\
\textit{Seed / days} & 20261117 / 3 \\
\textit{Mosquitoes seeded} & 1 female, 0 male \\
\textit{People} & 1; distance from the first mosquito: 2000 m \\
\textit{Closes} & none \\
\textit{Runs a behaviour test today} & seeding-only \\
\end{tabular}

\par\smallskip\noindent\textbf{Purpose.} Reports how far a female travels in one night and how far she ends from her start, when the person is far beyond any single night's reach. It checks the limits of the movement law, not contact, so the arrival of the female at the person is not a pass or fail criterion.
\par\smallskip\noindent\textbf{Setup.} One host-seeking female adult and one susceptible person placed far beyond what the female can cover in one night, with no aquatic habitat. The run lasts three days. The person stays out of reach for the whole run unless the walk happens to drift towards it.
\par\smallskip\noindent\textbf{What passing or failing implies.} The bounds say the movement law cannot produce a trajectory longer than the daily flight budget, nor a net shift beyond the exploration disc, so females cannot be carried across the landscape in one step. If either bound is exceeded, the macro-step or the exploration disc is wrong. A female reaching the person is not a failure, because the direction of the walk decides it.
\par\smallskip\noindent\textbf{Expected behaviour.} Report only. Per night, the trajectory length is at most 590 m ('daily\_move\_capacity\_m', [P], the mean distance travelled per day in the literature). Per night, the net displacement is at most 200 m, which is 2 x 'exploration\_range\_m' (100 m, [EP]). The earlier band of 1,410 to 2,000 m is withdrawn because it cannot be reached under these bounds.
\par\smallskip\noindent\textbf{What the current chain does.} The current chain keeps the single adult in its seeded cell for the whole run and records one successful human bite on day 1 and none afterwards; the bite is drawn against the cell-aggregated human population (about 127,000 people), not against the seeded person
\par\smallskip\noindent\textbf{Checked today by the test suite.}
\begin{itemize}
\item Seeding fidelity: the adult population obeys the conservation identity on every day (adults today = adults yesterday - deaths + emerged, starting from the seeded count), so no seeded individual is dropped or invented; the cohort log and the raster output agree on the number of adults
\item The run completes through the manifest route ('validate\_completeness', resolve, binary) with exit code 0
\item Baseline pin: the known-failure facts above are asserted, so any change in the chain's behaviour is flagged and must be acknowledged by updating the pin
\item Target gate (strict expected failure): the plan's target is asserted and is expected to fail until Step 5; when it starts passing the test fails and must be promoted
\item Target for E02: per-night net displacement bound of 200 m
\end{itemize}
\par\noindent\textbf{What the GIF must show.} Today the transmission GIF has one frame per simulated day and no scenario overlay. For E02 show the female track of each night in its own colour, the start point of the night, a 200 m ring around it, and the far person marker. A correct run shows each night's net shift inside its ring, while the track itself can wander and stays within the daily length budget.

\subsection*{E03 --- Several bites per night in a contact cluster}
\begin{tabular}{@{}ll@{}}
\textit{Mode} & gate \\
\textit{Seed / days} & 20261117 / 3 \\
\textit{Mosquitoes seeded} & 20 female, 0 male \\
\textit{People} & 5; distance from the first mosquito: 141 m, 424 m, 707 m, 728 m, 728 m \\
\textit{Closes} & Step 3 \\
\textit{Runs a behaviour test today} & seeding-only \\
\end{tabular}

\par\smallskip\noindent\textbf{Purpose.} Checks that one mosquito can make several blood-feeding attempts in one night and that each attempt is logged as its own event, not folded into one cell-level count. It is the first test of the multi-bite ledger.
\par\smallskip\noindent\textbf{Setup.} A cluster of host-seeking females with a few susceptible people placed among them, so that contacts occur in several places during the run, in an area with no aquatic habitat. The design freezes the probing success rate so that the expected count is a closed-form product. The JSON does not yet carry that override (its overrides are empty); the field enters with Step 3.
\par\smallskip\noindent\textbf{What passing or failing implies.} The ledger must count bites, not cell aggregates. If events per night equal the number of populated cells rather than the number of bites, transmission is understated in dense clusters and the per-person bookkeeping breaks. A pass confirms the multi-bite rule; a fail means the one-bite-per-night simplification is still present.
\par\smallskip\noindent\textbf{Expected behaviour.} Gate. The number of meals equals the number of females times the probing success rate times the probability of contact, computed by hand before the code is written. The ledger records one event per bite attempt, so a night shows several events, not one per cell.
\par\smallskip\noindent\textbf{What the current chain does.} Not meaningful yet: depends on Step 3 (blood-meal physiology and multi-bite pickup).
\par\smallskip\noindent\textbf{Checked today by the test suite.}
\begin{itemize}
\item Seeding fidelity: the adult population obeys the conservation identity on every day (adults today = adults yesterday - deaths + emerged, starting from the seeded count), so no seeded individual is dropped or invented; the cohort log and the raster output agree on the number of adults
\item The run completes through the manifest route ('validate\_completeness', resolve, binary) with exit code 0
\end{itemize}
\par\noindent\textbf{What the GIF must show.} One frame per phase of each night. Each female is a marker and each bite is a flash on the person it hits, with a running count of bites per person for the night. A correct run shows several flashes in one night from the same female, and the flashes land on different people as the females move.

\subsection*{E04 --- Partial meal does not start the cycle}
\begin{tabular}{@{}ll@{}}
\textit{Mode} & gate \\
\textit{Seed / days} & 20261117 / 3 \\
\textit{Mosquitoes seeded} & 1 female, 0 male \\
\textit{People} & 1; distance from the first mosquito: 200 m \\
\textit{Closes} & Step 3 \\
\textit{Runs a behaviour test today} & seeding-only \\
\end{tabular}

\par\smallskip\noindent\textbf{Purpose.} Checks that a female whose blood meal is below target does not start the gonotrophic cycle: she keeps host-seeking that night and nothing is laid.
\par\smallskip\noindent\textbf{Setup.} One host-seeking female and one susceptible person placed next to each other, so a meal is possible. The intent is that the female starts with a partial blood volume below the target. The JSON does not yet carry a blood-volume field, so the partial volume is seeded when Step 3 adds it.
\par\smallskip\noindent\textbf{What passing or failing implies.} Without this rule a single small feed would start the cycle, overstating egg production and the turnover of biting females. If the cycle advances after a partial meal, the physiology is not using the target volume. Real females need a sufficient meal before they mature eggs.
\par\smallskip\noindent\textbf{Expected behaviour.} Gate. After a partial meal the female remains host-seeking that night and the cycle does not start. The threshold is 'blood\_meal\_target\_ul' = 0.5 uL [EP].
\par\smallskip\noindent\textbf{What the current chain does.} Not meaningful yet: depends on Step 3 (blood-meal volume and the cycle clock).
\par\smallskip\noindent\textbf{Checked today by the test suite.}
\begin{itemize}
\item Seeding fidelity: the adult population obeys the conservation identity on every day (adults today = adults yesterday - deaths + emerged, starting from the seeded count), so no seeded individual is dropped or invented; the cohort log and the raster output agree on the number of adults
\item The run completes through the manifest route ('validate\_completeness', resolve, binary) with exit code 0
\end{itemize}
\par\noindent\textbf{What the GIF must show.} One frame per phase. The female is drawn in host-seeking colour, with a blood volume bar beside her that grows with each meal and a cycle timer beside it. A correct run keeps the female in host-seeking colour after the partial meal and the cycle timer stays at zero.

\subsection*{E05 --- Egg batch proportional to accumulated blood}
\begin{tabular}{@{}ll@{}}
\textit{Mode} & gate \\
\textit{Seed / days} & 20261117 / 3 \\
\textit{Mosquitoes seeded} & 2 female, 0 male \\
\textit{People} & 1; distance from the first mosquito: 200 m \\
\textit{Closes} & Step 3 \\
\textit{Runs a behaviour test today} & seeding-only \\
\end{tabular}

\par\smallskip\noindent\textbf{Purpose.} Checks that the egg batch scales with the accumulated blood volume a female actually took, not with a fixed binomial batch, by comparing a partially fed female with a fully fed one.
\par\smallskip\noindent\textbf{Setup.} Two host-seeking females next to the same susceptible person. One is meant to start with a partial blood volume and the other with a full meal. The JSON does not yet seed volumes, so the two females start identical until Step 3 adds the volume field.
\par\smallskip\noindent\textbf{What passing or failing implies.} Egg output must follow the blood volume the female took, which is the link the calibration rests on. If both females lay the same batch, the model loses the feeding-fecundity link that the literature measures, and the egg count no longer responds to host access.
\par\smallskip\noindent\textbf{Expected behaviour.} Gate. The egg count equals the blood volume times 223 eggs per uL [EP], rounded. This is calibrated so that a 0.5 uL meal gives about 112 eggs, close to the 111.67 plus or minus 5.36 measured in coluzzii [P]. The batch is proportional to the accumulated volume, so the fully fed female lays more than the partial one.
\par\smallskip\noindent\textbf{What the current chain does.} Not meaningful yet: depends on Step 3 (egg batch proportional to blood volume).
\par\smallskip\noindent\textbf{Checked today by the test suite.}
\begin{itemize}
\item Seeding fidelity: the adult population obeys the conservation identity on every day (adults today = adults yesterday - deaths + emerged, starting from the seeded count), so no seeded individual is dropped or invented; the cohort log and the raster output agree on the number of adults
\item The run completes through the manifest route ('validate\_completeness', resolve, binary) with exit code 0
\end{itemize}
\par\noindent\textbf{What the GIF must show.} Two females side by side, each with a blood volume bar and an egg count shown at the end of the run. A correct run shows the two bars in the same ratio as the two egg counts.

\subsection*{E06 --- First cycle needs two meals}
\begin{tabular}{@{}ll@{}}
\textit{Mode} & gate \\
\textit{Seed / days} & 20261117 / 3 \\
\textit{Mosquitoes seeded} & 1 female, 0 male \\
\textit{People} & 1; distance from the first mosquito: 200 m \\
\textit{Closes} & Step 3 \\
\textit{Runs a behaviour test today} & seeding-only \\
\end{tabular}

\par\smallskip\noindent\textbf{Purpose.} Checks that a teneral female, which has to complete a first gonotrophic cycle, needs two blood meals, because the first-cycle target is above the usual meal target.
\par\smallskip\noindent\textbf{Setup.} One teneral female adult, freshly emerged, next to one susceptible person. The JSON seeds the teneral state; the two-meal rule comes from the first-cycle and usual meal targets that Step 3 introduces.
\par\smallskip\noindent\textbf{What passing or failing implies.} If one meal completes the first cycle, the model underestimates the time to first oviposition and therefore the pace at which young females transmit. Requiring two meals keeps the young-female pace in line with the physiology the plan adopts.
\par\smallskip\noindent\textbf{Expected behaviour.} Gate. One meal is not enough. The female needs two, because 'first\_cycle\_target\_ul' = 0.8 uL [EP] is above 'blood\_meal\_target\_ul' = 0.5 uL [EP]. The first cycle closes only once the first-cycle target is reached.
\par\smallskip\noindent\textbf{What the current chain does.} Not meaningful yet: depends on Step 3 (first-cycle target of two meals).
\par\smallskip\noindent\textbf{Checked today by the test suite.}
\begin{itemize}
\item Seeding fidelity: the adult population obeys the conservation identity on every day (adults today = adults yesterday - deaths + emerged, starting from the seeded count), so no seeded individual is dropped or invented; the cohort log and the raster output agree on the number of adults
\item The run completes through the manifest route ('validate\_completeness', resolve, binary) with exit code 0
\end{itemize}
\par\noindent\textbf{What the GIF must show.} One frame per phase, with the female's blood volume bar and a marker for the first-cycle target. A correct run shows the bar crossing the usual target after the first meal but not the first-cycle target, and the cycle completes only after the second meal.

\subsection*{E07 --- Breeding site beats the host}
\begin{tabular}{@{}ll@{}}
\textit{Mode} & gate \\
\textit{Seed / days} & 20261117 / 3 \\
\textit{Mosquitoes seeded} & 1 female, 0 male \\
\textit{Aquatic cohorts} & 20 PUPA at (5100, 5200) m \\
\textit{People} & 1; distance from the first mosquito: 2000 m \\
\textit{Closes} & Step 7 and Step 9 \\
\textit{Runs a behaviour test today} & seeding-only \\
\end{tabular}

\par\smallskip\noindent\textbf{Purpose.} Two assertions in one scenario. Assertion A: an emerging adult appears at the position of its mother cohort, not at the centre of the cell. Assertion B: a gravid female chooses the breeding habitat, not the host, when both are on offer. Assertion B needs a gravid female, which the JSON does not seed yet.
\par\smallskip\noindent\textbf{Setup.} One host-seeking female, a nearly emerged pupal cohort with high development a short distance from her, and one person far beyond her reach. The layout puts the breeding site nearer than the host, so the choice is between them.
\par\smallskip\noindent\textbf{What passing or failing implies.} For A, emergence at the cell centre would leak the cell scale into the individual and lose the mother's location. For B, a female that goes to the person ignores the habitat signal, and biting is then driven by proximity rather than by the oviposition choice. Either failure would misplace the transmission the model reports.
\par\smallskip\noindent\textbf{Expected behaviour.} Gate with two assertions. A, closing Step 7: the emergence location equals the mother position, not the cell centre. B, closing Step 9: the oviposition site is the habitat site, not the host. Oviposition in natal habitat is not modelled (section 16.1), so B is about choosing habitat over host, not about the natal site.
\par\smallskip\noindent\textbf{What the current chain does.} Not meaningful yet: depends on Step 7 (emergence at the mother position) and Step 9 (oviposition site choice).
\par\smallskip\noindent\textbf{Checked today by the test suite.}
\begin{itemize}
\item Seeding fidelity: the adult population obeys the conservation identity on every day (adults today = adults yesterday - deaths + emerged, starting from the seeded count), so no seeded individual is dropped or invented; the cohort log and the raster output agree on the number of adults
\item The run completes through the manifest route ('validate\_completeness', resolve, binary) with exit code 0
\item Seeded pupae emerge: between one and the seeded number of adults emerge from the pupae seeded at the breeding site outside the candidate cells, so aquatic seeding at a metric position is live
\end{itemize}
\par\noindent\textbf{What the GIF must show.} One frame per phase. The pupal cohort is a marker with its emergence point, the female has a track, and the person is a far marker. For A the new adult should appear on the pupal marker, not at the cell centre. For B the female's track should end on the habitat marker rather than on the person.

\subsection*{E08 --- Two females, two persons, distinct identities}
\begin{tabular}{@{}ll@{}}
\textit{Mode} & gate \\
\textit{Seed / days} & 20261117 / 3 \\
\textit{Mosquitoes seeded} & 2 female, 0 male \\
\textit{People} & 2; distance from the first mosquito: 60 m, 7114 m \\
\textit{Closes} & Step 2 \\
\textit{Runs a behaviour test today} & seeding-only \\
\end{tabular}

\par\smallskip\noindent\textbf{Purpose.} Checks that each person has a unique, stable identifier and that each female is linked to the person she is matched with. The plan closes it at Step 2 with identifier checks: uniqueness, stability and the number of persons per cell.
\par\smallskip\noindent\textbf{Setup.} Two host-seeking females in two cells far apart from each other, each with a susceptible person well inside her activation radius. The pairs are separated so that no female is close to the other female's person.
\par\smallskip\noindent\textbf{What passing or failing implies.} If person identifiers collide or change between runs, the bite ledger cannot say who was bitten and the individual-level transmission cannot be traced. A pass means each person keeps one identity; it does not by itself prove the choice rule of the cascade.
\par\smallskip\noindent\textbf{Expected behaviour.} Gate at Step 2, verifiable by construction: each person has a unique human identifier that is deterministic from the seed and stable across runs (the earlier cell-times-slots formula was dropped by decision 2a), and the number of persons per cell matches the population. From Step 5 each female bites the person the weighted roulette picks among those inside her perception radius, with distance as the criterion of the stretch, and there is no winning cell.
\par\smallskip\noindent\textbf{What the current chain does.} Not meaningful yet: depends on Step 2 (individual humans with stable identifiers).
\par\smallskip\noindent\textbf{Checked today by the test suite.}
\begin{itemize}
\item Seeding fidelity: the adult population obeys the conservation identity on every day (adults today = adults yesterday - deaths + emerged, starting from the seeded count), so no seeded individual is dropped or invented; the cohort log and the raster output agree on the number of adults
\item The run completes through the manifest route ('validate\_completeness', resolve, binary) with exit code 0
\end{itemize}
\par\noindent\textbf{What the GIF must show.} One frame per phase, with each person drawn in a colour fixed by identifier, so that a person keeps the same colour across frames. A correct run shows each person keeping its colour and each bite landing on exactly one of them.

\subsection*{E09 --- Infective female, several persons}
\begin{tabular}{@{}ll@{}}
\textit{Mode} & gate \\
\textit{Seed / days} & 20261117 / 3 \\
\textit{Mosquitoes seeded} & 1 female, 0 male \\
\textit{People} & 3; distance from the first mosquito: 40 m, 60 m, 64 m \\
\textit{Closes} & Step 2 \\
\textit{Runs a behaviour test today} & seeding-only \\
\end{tabular}

\par\smallskip\noindent\textbf{Purpose.} Checks the one-way rule of transmission at the bite: an infective female gives several infectious bites across the run when several persons are around, the infection can reach more than one person across the set of runs, and each infectious bite infects at most the person it bites.
\par\smallskip\noindent\textbf{Setup.} One infective host-seeking female with three susceptible persons placed around her at different short distances, run for three days. The infective flag is set on the mosquito and all persons start susceptible.
\par\smallskip\noindent\textbf{What passing or failing implies.} If one bite infects anyone other than the person bitten, the model creates transmission that no contact produced and inflates incidence. If the infection never reaches a second person across the runs, the spread between persons is missing and the model cannot show onward transmission.
\par\smallskip\noindent\textbf{Expected behaviour.} Gate at Step 2. The run contains several infectious bites; the probability that more than one person is infected is above zero; each infectious bite infects at most the person bitten. The check is structural, not a fixed band.
\par\smallskip\noindent\textbf{What the current chain does.} Not meaningful yet: depends on Step 2 (individual humans and the per-bite ledger).
\par\smallskip\noindent\textbf{Checked today by the test suite.}
\begin{itemize}
\item Seeding fidelity: the adult population obeys the conservation identity on every day (adults today = adults yesterday - deaths + emerged, starting from the seeded count), so no seeded individual is dropped or invented; the cohort log and the raster output agree on the number of adults
\item The run completes through the manifest route ('validate\_completeness', resolve, binary) with exit code 0
\end{itemize}
\par\noindent\textbf{What the GIF must show.} One frame per phase. The infective female is a red marker, the persons are grey until infected, and each infectious bite turns its person red. A correct run shows red bite flashes that change only the bitten person, never a neighbour.

\subsection*{E10 --- Sensitivity to the three perception radii}
\begin{tabular}{@{}ll@{}}
\textit{Mode} & report \\
\textit{Seed / days} & 20261117 / 3 \\
\textit{Mosquitoes seeded} & 1 female, 0 male \\
\textit{People} & 1; distance from the first mosquito: 150 m \\
\textit{Variants} & activacion\_co2\_50\_100\_200\_m, olfativa\_10\_30\_70\_m, visual\_0.25\_0.72\_m \\
\textit{Overrides} & co2\_activation\_radius\_m, olfactory\_radius\_m, visual\_radius\_m \\
\textit{Closes} & none \\
\textit{Runs a behaviour test today} & seeding-only \\
\end{tabular}

\par\smallskip\noindent\textbf{Purpose.} Quantifies how much each stage of the sensory cascade (carbon dioxide activation, olfactory tracking, visual approach) moves the outcome, by sweeping one radius at a time. It is a report: it measures sensitivity and does not gate.
\par\smallskip\noindent\textbf{Setup.} One host-seeking female and one susceptible person a short distance away, run three days, once per variant. The three variants live in one file, as entries in its variants list, and each varies one radius with the others at their defaults. Variant one sweeps the carbon dioxide activation radius over 50, 100 and 200 m. Variant two sweeps the olfactory radius over 10, 30 and 70 m. Variant three sweeps the visual radius over 0.25 and 0.72 m. The defaults are the verified values.
\par\smallskip\noindent\textbf{What passing or failing implies.} The result shows which radius the outcome depends on. If the outcome barely moves across a sweep, that stage is doing little work and the cascade is poorly identified there; if it moves a lot, the default needs evidence, which the plan marks as [EP] where there is none. This tells the team which radius to calibrate first.
\par\smallskip\noindent\textbf{Expected behaviour.} Report only, with no pass or fail. For each variant the observed contacts and outcome are written next to the default run. The plan gives no band here, because the sweep is itself the result.
\par\smallskip\noindent\textbf{What the current chain does.} Not meaningful yet: depends on Step 5 (the cascade). Today the overrides are not applied: the engine rejects non-empty overrides until Step 5, so the harness runs each variant without them, with the same seed, and records the overrides as not applied. The three variants therefore reduce to the same run and the report cannot show sensitivity.
\par\smallskip\noindent\textbf{Checked today by the test suite.}
\begin{itemize}
\item Seeding fidelity: the adult population obeys the conservation identity on every day (adults today = adults yesterday - deaths + emerged, starting from the seeded count), so no seeded individual is dropped or invented; the cohort log and the raster output agree on the number of adults
\item The run completes through the manifest route ('validate\_completeness', resolve, binary) with exit code 0
\end{itemize}
\par\noindent\textbf{What the GIF must show.} One GIF per variant. Each shows the female track, the person, and the radius circle drawn at the swept value. Once the overrides apply, the three GIFs should differ in circle size and in when contact happens; if they look the same, the sweep is not reaching the engine.

\subsection*{E11 --- Postprandial window after a full meal}
\begin{tabular}{@{}ll@{}}
\textit{Mode} & gate \\
\textit{Seed / days} & 20261117 / 3 \\
\textit{Mosquitoes seeded} & 1 female, 0 male \\
\textit{People} & 1; distance from the first mosquito: 100 m \\
\textit{Closes} & Step 3 \\
\textit{Runs a behaviour test today} & seeding-only \\
\end{tabular}

\par\smallskip\noindent\textbf{Purpose.} Checks that a female that has completed a meal cannot try again before the postprandial window closes, and records how the return to host-seeking rises afterwards. The window is a gate; the reopening curve is reported only.
\par\smallskip\noindent\textbf{Setup.} One host-seeking female, fed on day zero, in the same cell as a susceptible person, so that she would try to feed again at once if no window existed. The JSON does not yet seed a fed state; the fed state comes from the Step 3 physiology.
\par\smallskip\noindent\textbf{What passing or failing implies.} If the female bites again immediately, the blood-feeding rate is too high and transmission is inflated. The reopening curve describes how fast a fed female resumes host-seeking. It is an estimate without a direct source, so it is reported with its band, not gated.
\par\smallskip\noindent\textbf{Expected behaviour.} Gate. The first attempt after the target is reached is not before 12 hours, given 'blood\_meal\_refractory\_days' = 0.5 day [EP], and falls in a band of 12 hours plus one phase. Reported only, with [EP] tier: progressive reopening of 7 percent at 40 hours, 27 percent at 48 hours and 68 percent at 72 hours. The reopening is not a failure criterion.
\par\smallskip\noindent\textbf{What the current chain does.} Not meaningful yet: depends on Step 3 (post-meal refractory window).
\par\smallskip\noindent\textbf{Checked today by the test suite.}
\begin{itemize}
\item Seeding fidelity: the adult population obeys the conservation identity on every day (adults today = adults yesterday - deaths + emerged, starting from the seeded count), so no seeded individual is dropped or invented; the cohort log and the raster output agree on the number of adults
\item The run completes through the manifest route ('validate\_completeness', resolve, binary) with exit code 0
\end{itemize}
\par\noindent\textbf{What the GIF must show.} One frame per phase, with the time since the meal shown on the female marker and a shaded window of the first 12 hours. A correct run shows no bite or attempt inside the shaded window and the first attempt just after it.

\subsection*{E12 --- Ovarian arrest in the dry season}
\begin{tabular}{@{}ll@{}}
\textit{Mode} & gate \\
\textit{Seed / days} & 20261117 / 30 \\
\textit{Mosquitoes seeded} & 20 female, 0 male \\
\textit{People} & 5; distance from the first mosquito: 141 m, 424 m, 707 m, 728 m, 728 m \\
\textit{Closes} & Step 6 \\
\textit{Runs a behaviour test today} & seeding-only \\
\end{tabular}

\par\smallskip\noindent\textbf{Purpose.} Checks that the fraction of fed females that stop ovarian development during a dry spell (arrest, not death) falls inside the published band for coluzzii. The band is the published data, not a parameter of the scenario.
\par\smallskip\noindent\textbf{Setup.} A group of host-seeking females in one cell, with a few persons nearby, run for thirty days. The environment has a burst of rain followed by a long dry spell, set in the climate input rather than in the scenario file. The cell has no eligible breeding site during the dry spell, so fed females cannot lay.
\par\smallskip\noindent\textbf{What passing or failing implies.} If too few females arrest, dry-season egg production and population are too high and the seasonal dip the data show is missing. If too many arrest, the recovery after the dry season is overstated. Because the band is published data, a fail points to the model, not to a tuned parameter.
\par\smallskip\noindent\textbf{Expected behaviour.} Gate. The fraction of fed females that do not pass stage 2 lies within 19.23 to 34.85 percent, the published range for coluzzii [P]. The scenario measures the fraction and compares it with the band; the band is not adjusted.
\par\smallskip\noindent\textbf{What the current chain does.} Not meaningful yet: depends on Step 6 (arrested state and the dry-spell trigger).
\par\smallskip\noindent\textbf{Checked today by the test suite.}
\begin{itemize}
\item Seeding fidelity: the adult population obeys the conservation identity on every day (adults today = adults yesterday - deaths + emerged, starting from the seeded count), so no seeded individual is dropped or invented; the cohort log and the raster output agree on the number of adults
\item The run completes through the manifest route ('validate\_completeness', resolve, binary) with exit code 0
\end{itemize}
\par\noindent\textbf{What the GIF must show.} One frame per day, with the females coloured by ovarian state (pre-gravid, gravid or arrested) and the rain burst and dry spell drawn as a background band. A correct run shows the arrested share rising during the dry spell and the arrested females still alive when the rain returns.

\subsection*{E13 --- Males and mating supply at one swarm site}
\begin{tabular}{@{}ll@{}}
\textit{Mode} & report \\
\textit{Seed / days} & 20261117 / 3 \\
\textit{Mosquitoes seeded} & 100 female, 100 male \\
\textit{Closes} & none \\
\textit{Runs a behaviour test today} & seeding-only \\
\end{tabular}

\par\smallskip\noindent\textbf{Purpose.} Reports the mating supply one swarm site can offer and how many females remain uninseminated after three nights. It is a report, not a gate, because the plan derives its band by hand from the supply formula before the code is written.
\par\smallskip\noindent\textbf{Setup.} Equal numbers of freshly emerged females and males in a single 1 square km cell, all of whom can reach one swarm site, so every male and female is assigned to it. Some males may stay loose, and the scenario measures how many. The swarm marker density is meant to be fixed in the scenario; the JSON does not yet carry it, and the value enters with Step 8.
\par\smallskip\noindent\textbf{What passing or failing implies.} If the supply is far from the formula, the mating model is inconsistent with the published rates and the uninseminated share is wrong, which matters for population growth after a dry spell. The report lets the team check the number before accepting the swarm model. Because the male count sets the number of pairs, a mismatch appears here first.
\par\smallskip\noindent\textbf{Expected behaviour.} Report band from the formula in section 8.1. Supply per site per night equals the male count times 'p\_mate\_male\_fraction' = 0.072 [P]; for 100 males that is 7.2 pairs. Females present equal the female count times 'female\_attendance\_rate' = 0.1033 [P]; for 100 females that is 10.33. The probability that a female is inseminated is 0.1033 x (1 - exp(-7.2/10.33)) = 0.0518, so about 85.2 percent of females remain uninseminated after three nights [EP]. That figure is a lower bound, because available males decrease as they mate. The earlier published band of 22.3 to 52.8 percent is withdrawn because it cannot be derived from the formula; the implementer computes the report tolerance from this arithmetic.
\par\smallskip\noindent\textbf{What the current chain does.} Not meaningful yet: depends on Step 8 (male site assignment and swarm supply).
\par\smallskip\noindent\textbf{Checked today by the test suite.}
\begin{itemize}
\item Seeding fidelity: the adult population obeys the conservation identity on every day (adults today = adults yesterday - deaths + emerged, starting from the seeded count), so no seeded individual is dropped or invented; the cohort log and the raster output agree on the number of adults
\item The run completes through the manifest route ('validate\_completeness', resolve, binary) with exit code 0
\item Seed sensitivity: two master seeds give different mortality among the 200 seeded adults, so the sub-streams really derive from the single seed
\end{itemize}
\par\noindent\textbf{What the GIF must show.} One frame per night, with the swarm site marked, the males and females near it coloured by assignment and by insemination state, and a running count of inseminated females at the top. A correct run shows the inseminated count rising and levelling off, with a visible share still uninseminated after three nights.

\subsection*{E14 --- Males present before the swarm window}
\begin{tabular}{@{}ll@{}}
\textit{Mode} & gate \\
\textit{Seed / days} & 20261117 / 1 \\
\textit{Mosquitoes seeded} & 1 female, 20 male \\
\textit{Closes} & Step 8 \\
\textit{Runs a behaviour test today} & seeding-only \\
\end{tabular}

\par\smallskip\noindent\textbf{Purpose.} Checks that freshly emerged males assigned to a site do not enter the swarm before their first swarm window, so they cannot inseminate a female early, and that the alarm for uninseminated females is raised.
\par\smallskip\noindent\textbf{Setup.} One host-seeking female with a group of freshly emerged males from the same cell, clustered close to her, run for one day. The day ends before the males' first swarm window. The males carry no state beyond their sex and age. The female is placed close to them, so she is within reach of the site if the site is near.
\par\smallskip\noindent\textbf{What passing or failing implies.} If males enter the swarm before their window, mating happens before the males are physiologically ready and early insemination is overstated. If the alarm is missing, a broken swarm field goes unnoticed, and the alarm is the signal the plan uses to detect that.
\par\smallskip\noindent\textbf{Expected behaviour.} Gate. No female is inseminated by males before their first swarm window. The alarm 'hembras\_sin\_inseminar' is emitted for the cell and night. The female attends the site only if the site lies within her reach.
\par\smallskip\noindent\textbf{What the current chain does.} Not meaningful yet: depends on Step 8 (male site assignment, swarm windows and the alarm).
\par\smallskip\noindent\textbf{Checked today by the test suite.}
\begin{itemize}
\item Seeding fidelity: the adult population obeys the conservation identity on every day (adults today = adults yesterday - deaths + emerged, starting from the seeded count), so no seeded individual is dropped or invented; the cohort log and the raster output agree on the number of adults
\item The run completes through the manifest route ('validate\_completeness', resolve, binary) with exit code 0
\end{itemize}
\par\noindent\textbf{What the GIF must show.} One frame per phase, with the swarm window shaded in the evening. The males are grey before the window and turn to swarm colour only inside it. A correct run shows no male in swarm colour before the shaded window and no insemination during the run.

\subsection*{E15 --- Swarm killing empties the site}
\begin{tabular}{@{}ll@{}}
\textit{Mode} & gate \\
\textit{Seed / days} & 20261117 / 3 \\
\textit{Mosquitoes seeded} & 50 female, 50 male \\
\textit{Closes} & Step 8 \\
\textit{Runs a behaviour test today} & seeding-only \\
\end{tabular}

\par\smallskip\noindent\textbf{Purpose.} Checks that the swarm field is a static, discrete layer: killing the males assigned to one marker removes that site's supply, so the share of uninseminated females rises. It is the test that forces markers to be discrete and removable.
\par\smallskip\noindent\textbf{Setup.} A field of swarm markers with one swarm in the village, equal numbers of males and females, run for three days. The kill of the males assigned to one marker is an intervention the scenario applies. The JSON does not express the kill yet; the flag enters with Step 8.
\par\smallskip\noindent\textbf{What passing or failing implies.} If killing the males at a marker does not reduce supply, the swarm is a number rather than a place, and the model could not represent a targeted intervention. A pass supports the static site layer; a fail means site assignment does not reach the supply.
\par\smallskip\noindent\textbf{Expected behaviour.} Gate, directional. After the kill the supply of the killed site falls to zero, and the share of uninseminated females is higher than in the same-seed run without the kill. The plan gives no numeric threshold for the rise, so the test checks the direction of the change and that the killed site's supply is removed.
\par\smallskip\noindent\textbf{What the current chain does.} Not meaningful yet: depends on Step 8 (site supply and the swarm-killing flag).
\par\smallskip\noindent\textbf{Checked today by the test suite.}
\begin{itemize}
\item Seeding fidelity: the adult population obeys the conservation identity on every day (adults today = adults yesterday - deaths + emerged, starting from the seeded count), so no seeded individual is dropped or invented; the cohort log and the raster output agree on the number of adults
\item The run completes through the manifest route ('validate\_completeness', resolve, binary) with exit code 0
\item Seed sensitivity: two master seeds give different mortality among the seeded adults
\end{itemize}
\par\noindent\textbf{What the GIF must show.} One frame per night, with the markers as fixed points, the males and females coloured by assignment, and the killed marker crossed out after the kill. A correct run shows the killed site's males disappearing and its females staying uninseminated while other sites keep mating.

\subsection*{E16 --- Nets at 60 percent do not offset swarm supply}
\begin{tabular}{@{}ll@{}}
\textit{Mode} & report \\
\textit{Seed / days} & 20261117 / 3 \\
\textit{Mosquitoes seeded} & 50 female, 50 male \\
\textit{Closes} & none \\
\textit{Runs a behaviour test today} & seeding-only \\
\end{tabular}

\par\smallskip\noindent\textbf{Purpose.} Counter-scenario that records a model limit: individual protection from insecticide-treated nets at 60 percent coverage does not restore the swarm supply lost at a site. It is a report because the plan treats the result as a documented limit, not as a model output to tune.
\par\smallskip\noindent\textbf{Setup.} The same swarm-site field as the swarm-killing scenario, with equal numbers of males and females, run for three days, with net coverage set to 60 percent. The net coverage belongs to the scenario definition; the JSON does not carry a net field yet.
\par\smallskip\noindent\textbf{What passing or failing implies.} If nets appeared to restore supply, the model would be suggesting that individual protection can compensate for lost swarm sites, which the plan does not support. Keeping the limit in the report stops a reader from taking net coverage as a mating-supply intervention.
\par\smallskip\noindent\textbf{Expected behaviour.} Report only, with no band in the plan. The report records the uninseminated share and the supply per site with nets and without, and states that the nets do not compensate for the supply loss. The plan records this as a documented limit of the model, not as a result.
\par\smallskip\noindent\textbf{What the current chain does.} Not meaningful yet: depends on Step 8 (the swarm supply the nets are compared against).
\par\smallskip\noindent\textbf{Checked today by the test suite.}
\begin{itemize}
\item Seeding fidelity: the adult population obeys the conservation identity on every day (adults today = adults yesterday - deaths + emerged, starting from the seeded count), so no seeded individual is dropped or invented; the cohort log and the raster output agree on the number of adults
\item The run completes through the manifest route ('validate\_completeness', resolve, binary) with exit code 0
\end{itemize}
\par\noindent\textbf{What the GIF must show.} Same layout as the swarm-killing GIF, with a net coverage indicator at 60 percent drawn on the frame. A correct run shows the sites that lose supply staying weak with nets in place.

\subsection*{E17 --- Sterile males compete for site supply}
\begin{tabular}{@{}ll@{}}
\textit{Mode} & gate \\
\textit{Seed / days} & 20261117 / 3 \\
\textit{Mosquitoes seeded} & 50 female, 50 male \\
\textit{Closes} & Step 10 \\
\textit{Runs a behaviour test today} & seeding-only \\
\end{tabular}

\par\smallskip\noindent\textbf{Purpose.} Checks that sterile males assigned to the same swarm sites as the fertile males take a proportional share of the supply, so that supply per site falls with their number and the uninseminated share rises. It is the sterile-insect test on site supply.
\par\smallskip\noindent\textbf{Setup.} Equal numbers of females and males assigned to the swarm sites, with sterile males added and assigned to the same sites, so they compete for the same supply. Run for three days. The sterility flag is not in the JSON yet; it enters with Step 10.
\par\smallskip\noindent\textbf{What passing or failing implies.} If supply does not fall with the number of sterile males assigned, the model cannot represent sterile-insect release, an intervention a programme may run. Competition among individual sterile males inside a swarm is not simulated (N11, section 8.6), so the test covers supply per site, not competition within the swarm.
\par\smallskip\noindent\textbf{Expected behaviour.} Gate. Supply per site falls in proportion to the number of sterile males assigned to it, and the uninseminated female share rises. The Step 10 criterion also requires that two identical runs match byte for byte (Level 5).
\par\smallskip\noindent\textbf{What the current chain does.} Not meaningful yet: depends on Step 10 (sterile males on the site supply).
\par\smallskip\noindent\textbf{Checked today by the test suite.}
\begin{itemize}
\item Seeding fidelity: the adult population obeys the conservation identity on every day (adults today = adults yesterday - deaths + emerged, starting from the seeded count), so no seeded individual is dropped or invented; the cohort log and the raster output agree on the number of adults
\item The run completes through the manifest route ('validate\_completeness', resolve, binary) with exit code 0
\end{itemize}
\par\noindent\textbf{What the GIF must show.} One frame per night, with fertile and sterile males drawn in different colours at each site and the supply of each site shown as a bar. A correct run shows each site's supply shrinking as its sterile share grows, and the uninseminated count rising with it.

\subsection*{E18 --- Wind deflection and night blocking}
\begin{tabular}{@{}ll@{}}
\textit{Mode} & gate \\
\textit{Seed / days} & 20261117 / 3 \\
\textit{Mosquitoes seeded} & 1 female, 0 male \\
\textit{Closes} & Step 5 \\
\textit{Runs a behaviour test today} & seeding-only \\
\end{tabular}

\par\smallskip\noindent\textbf{Purpose.} Checks the wind branch of the movement law: a female flying against the wind is deflected upwind, and forward progress is blocked when the tail-wind component reaches her flight speed. It is the gate for the wind branch.
\par\smallskip\noindent\textbf{Setup.} One host-seeking female at the centre of the synthetic study area, run for three days. The synthetic area carries a wind field of the 100 m wind components, loaded at run time. The wind is not written in the JSON; it comes from the synthetic area inputs.
\par\smallskip\noindent\textbf{What passing or failing implies.} If the female ignores the wind, the model cannot show how wind affects dispersal and therefore the reach of transmission. If the block never occurs, females would fly against strong winds at full speed. The wind is taken at 100 m, not at flight height, so the branch can only be validated against the synthetic area; this is a declared limit.
\par\smallskip\noindent\textbf{Expected behaviour.} Gate. The flight is deflected upwind by up to 30 degrees ('anemotaxis\_rule', [EP]). Forward progress is blocked when the tail-wind component, projected on the heading, reaches the flight speed of 0.225 to 0.45 m/s [P]. The threshold is computed at run time from the configured flight speed and is never hard-coded. A permanent night block under the measured wind is accepted as a documented result of the 100 m level, not as an error.
\par\smallskip\noindent\textbf{What the current chain does.} Not meaningful yet: depends on Step 5 (wind branch of the cascade). Without a loaded wind field the wind rule is inactive.
\par\smallskip\noindent\textbf{Checked today by the test suite.}
\begin{itemize}
\item Seeding fidelity: the adult population obeys the conservation identity on every day (adults today = adults yesterday - deaths + emerged, starting from the seeded count), so no seeded individual is dropped or invented; the cohort log and the raster output agree on the number of adults
\item The run completes through the manifest route ('validate\_completeness', resolve, binary) with exit code 0
\end{itemize}
\par\noindent\textbf{What the GIF must show.} One frame per phase, with the female track and wind arrows drawn at each frame, the arrow length showing wind speed. A correct run shows the track bent upwind relative to the wind arrows, and forward progress stopping when the tail-wind arrow reaches the flight speed.


% =====================================================================
\section{Generated parameter table}

The complete table of the model's constants, generated from
\code{mal-core/config/model_constants.yaml} by
\code{tools/gen_concepts_values.py}. The per-submodel tables of
Element~7 select the entries that matter for each submodel and add their
provenance.

% The generated table declares four unwrapped columns; the preamble is
% widened here to four wrapping columns instead of editing the generated
% file.
\begingroup
% The generated table has four unwrapped columns whose keys and file
% paths have no natural break point. The preamble is widened here, the
% path macro is made breakable and the table is set one size smaller,
% instead of editing the generated file.
\AtBeginEnvironment{longtable}{\scriptsize}
\setlength{\tabcolsep}{3pt}
\ExplSyntaxOn
\NewDocumentCommand{\brkpath}{m}
  {
    \group_begin:
    \tl_set:Nn \l_tmpa_tl {#1}
    \tl_replace_all:Nnn \l_tmpa_tl { / } {/ \allowbreak}
    \tl_replace_all:Nnn \l_tmpa_tl { \_ } {\_ \allowbreak}
    \tl_replace_all:Nnn \l_tmpa_tl { - } {- \allowbreak}
    \group_end:
  }
\ExplSyntaxOff
\renewcommand{\code}[1]{\texttt{\brkpath{#1}}}
\makeatletter
\let\origlongtableat\LT@array
\def\LT@array[#1]#2{\origlongtableat[#1]{@{}L{3.9cm}L{1.6cm}L{2.0cm}L{8.4cm}@{}}}
\makeatother
% Generated by tools/gen_concepts_values.py -- do not edit.
\paragraph{abm: model parameters}
\begin{longtable}{@{}llll@{}}
\toprule
Key & Value & Unit & Read by \\
\midrule
\endhead
\code{adults.daily\_mort\_basal} & 0.87 & prob/día & mal-core/src/mal\_core/abm/include/mal\_abm\_fast/config.hpp, mal-core/src/mal\_core/abm/include/mal\_abm\_fast/wire.hpp \\
\code{adults.mort\_cap} & 0.93 & fracción & mal-core/src/mal\_core/abm/include/mal\_abm\_fast/config.hpp, mal-core/src/mal\_core/abm/include/mal\_abm\_fast/wire.hpp \\
\code{adults.mort\_floor} & 0.6 & fracción & mal-core/src/mal\_core/abm/include/mal\_abm\_fast/config.hpp, mal-core/src/mal\_core/abm/include/mal\_abm\_fast/wire.hpp \\
\code{adults.opt\_c} & 25.0 & °C & mal-core/src/mal\_core/abm/include/mal\_abm\_fast/config.hpp, mal-core/src/mal\_core/abm/include/mal\_abm\_fast/wire.hpp, mal-core/src/mal\_core/abm/src/mosquito\_submodel.cpp \\
\code{adults.sigma\_c} & 15.0 & °C & mal-core/src/mal\_core/abm/include/mal\_abm\_fast/config.hpp, mal-core/src/mal\_core/abm/include/mal\_abm\_fast/wire.hpp \\
\code{aquatic\_stages.egg\_daily\_mort} & 0.08 & 1/día & mal-core/src/mal\_core/abm/include/mal\_abm\_fast/config.hpp, mal-core/src/mal\_core/abm/src/aquatic\_cohort\_bank.cpp \\
\code{aquatic\_stages.larva\_bh\_alpha} & 0.05 & adimensional & mal-core/src/mal\_core/abm/include/mal\_abm\_fast/config.hpp, mal-core/src/mal\_core/abm/include/mal\_abm\_fast/wire.hpp … (+2 more) \\
\code{aquatic\_stages.larva\_bh\_s0} & 0.95 & fracción & mal-core/src/mal\_core/abm/include/mal\_abm\_fast/config.hpp, mal-core/src/mal\_core/abm/include/mal\_abm\_fast/wire.hpp \\
\code{aquatic\_stages.last\_larval\_instar} & 4 & instar & mal-core/src/mal\_core/abm/include/mal\_abm\_fast/config.hpp, mal-core/src/mal\_core/abm/src/aquatic\_cohort\_bank.cpp \\
\code{aquatic\_stages.pupa\_daily\_mort} & 0.05 & 1/día & mal-core/src/mal\_core/abm/include/mal\_abm\_fast/config.hpp, mal-core/src/mal\_core/abm/src/aquatic\_cohort\_bank.cpp \\
\code{aquatic\_stages.salinity\_missing\_psu} & 0.0 & psu & mal-core/src/mal\_core/abm/include/mal\_abm\_fast/config.hpp, mal-core/src/mal\_core/abm/src/aquatic\_cohort\_bank.cpp, mal-core/src/mal\_core/abm/src/climate.cpp \\
\code{aquatic\_stress.desicc\_base\_daily} & 0.1 & 1/día & mal-core/src/mal\_core/abm/include/mal\_abm\_fast/aquatic\_stress.hpp, mal-core/src/mal\_core/abm/include/mal\_abm\_fast/config.hpp \\
\code{aquatic\_stress.desicc\_grace\_days} & 5.0 & días & mal-core/src/mal\_core/abm/include/mal\_abm\_fast/aquatic\_stress.hpp, mal-core/src/mal\_core/abm/include/mal\_abm\_fast/config.hpp \\
\code{aquatic\_stress.desicc\_ramp\_decay} & 0.3 & 1/día & mal-core/src/mal\_core/abm/include/mal\_abm\_fast/aquatic\_stress.hpp, mal-core/src/mal\_core/abm/include/mal\_abm\_fast/config.hpp \\
\code{aquatic\_stress.vulnerable\_instar\_max} & 2 & instar & mal-core/src/mal\_core/abm/include/mal\_abm\_fast/config.hpp, mal-core/src/mal\_core/abm/src/aquatic\_cohort\_bank.cpp \\
\code{aquatic\_stress.wash\_fraction\_max} & 0.6 & fracción & mal-core/src/mal\_core/abm/include/mal\_abm\_fast/aquatic\_stress.hpp, mal-core/src/mal\_core/abm/include/mal\_abm\_fast/config.hpp \\
\code{aquatic\_stress.wash\_threshold\_mm} & 40.0 & mm & mal-core/src/mal\_core/abm/include/mal\_abm\_fast/aquatic\_cohort\_bank.hpp, mal-core/src/mal\_core/abm/include/mal\_abm\_fast/aquatic\_stress.hpp … (+2 more) \\
\code{capacity.persistent\_availability} & 1.0 & adimensional & mal-core/src/mal\_core/abm/include/mal\_abm\_fast/config.hpp, mal-core/src/mal\_core/abm/include/mal\_abm\_fast/habitat\_availability.hpp \\
\code{contracts.cogs\_state} & "4.0" & string & mal-core/src/mal\_core/abm/include/mal\_abm\_fast/config.hpp, mal-core/src/mal\_core/abm/include/mal\_abm\_fast/output\_contract.hpp, mal-core/src/mal\_core/abm/include/mal\_abm\_fast/wire.hpp \\
\code{contracts.env\_nc} & "3.0" & string & (loaded, never used) \\
\code{contracts.habitat} & "4.0" & string & mal-core/src/mal\_core/abm/include/mal\_abm\_fast/aquatic\_cohort\_bank.hpp, mal-core/src/mal\_core/abm/include/mal\_abm\_fast/aquatic\_stress.hpp … (+28 more) \\
\code{dispersal.flight\_hour} & 20 & h & mal-core/src/mal\_core/abm/include/mal\_abm\_fast/config.hpp, mal-core/src/mal\_core/abm/src/engine.cpp \\
\code{dispersal.night\_end\_hour} & 6 & h & mal-core/src/mal\_core/abm/include/mal\_abm\_fast/config.hpp, mal-core/src/mal\_core/abm/include/mal\_abm\_fast/multirate\_scheduler.hpp \\
\code{dispersal.night\_start\_hour} & 18 & h & mal-core/src/mal\_core/abm/include/mal\_abm\_fast/config.hpp, mal-core/src/mal\_core/abm/include/mal\_abm\_fast/multirate\_scheduler.hpp \\
\code{dispersal.snap\_radius\_m} & 500.0 & m & mal-core/src/mal\_core/abm/include/mal\_abm\_fast/config.hpp, mal-core/src/mal\_core/abm/src/mosquito\_submodel.cpp \\
\code{dispersal.wind\_flight\_hours} & 4.0 & h & mal-core/src/mal\_core/abm/include/mal\_abm\_fast/config.hpp, mal-core/src/mal\_core/abm/include/mal\_abm\_fast/wire.hpp \\
\code{dispersal.wind\_migration\_prob} & 0.05 & 1/día & mal-core/src/mal\_core/abm/include/mal\_abm\_fast/config.hpp, mal-core/src/mal\_core/abm/include/mal\_abm\_fast/wire.hpp \\
\code{dispersal.wind\_self\_speed\_ms} & 1.0 & m/s & mal-core/src/mal\_core/abm/include/mal\_abm\_fast/config.hpp, mal-core/src/mal\_core/abm/include/mal\_abm\_fast/wire.hpp \\
\code{dispersal.wind\_survival} & 0.85 & fracción & mal-core/src/mal\_core/abm/include/mal\_abm\_fast/config.hpp, mal-core/src/mal\_core/abm/include/mal\_abm\_fast/wire.hpp \\
\code{eip.base\_c} & 16.0 & °C & mal-core/src/mal\_core/abm/include/mal\_abm\_fast/config.hpp, mal-core/src/mal\_core/abm/include/mal\_abm\_fast/eip.hpp, mal-core/src/mal\_core/abm/include/mal\_abm\_fast/wire.hpp \\
\code{eip.threshold\_gd} & 110.0 & grados-día & mal-core/src/mal\_core/abm/include/mal\_abm\_fast/config.hpp, mal-core/src/mal\_core/abm/include/mal\_abm\_fast/eip.hpp, mal-core/src/mal\_core/abm/include/mal\_abm\_fast/wire.hpp \\
\code{geodesy.earth\_radius\_m} & 6371008.8 & m & mal-core/src/mal\_core/abm/include/mal\_abm\_fast/aoi.hpp, mal-core/src/mal\_core/abm/include/mal\_abm\_fast/config.hpp, mal-core/src/mal\_core/abm/include/mal\_abm\_fast/wire.hpp \\
\code{geodesy.metres\_per\_deg\_lat} & 111195.08023353292 & m/° & mal-core/src/mal\_core/abm/include/mal\_abm\_fast/aoi.hpp, mal-core/src/mal\_core/abm/include/mal\_abm\_fast/config.hpp … (+2 more) \\
\code{gonotrophic.cycle\_duration\_days} & 2.0 & días & mal-core/src/mal\_core/abm/include/mal\_abm\_fast/config.hpp, mal-core/src/mal\_core/abm/include/mal\_abm\_fast/gonotrophic\_cycle.hpp … (+2 more) \\
\code{gonotrophic.egg\_batch\_max} & 170 & huevos & mal-core/src/mal\_core/abm/include/mal\_abm\_fast/config.hpp, mal-core/src/mal\_core/abm/include/mal\_abm\_fast/gonotrophic\_cycle.hpp, mal-core/src/mal\_core/abm/src/mosquito\_submodel.cpp \\
\code{gonotrophic.egg\_batch\_mean} & 120 & huevos & mal-core/src/mal\_core/abm/include/mal\_abm\_fast/config.hpp, mal-core/src/mal\_core/abm/include/mal\_abm\_fast/gonotrophic\_cycle.hpp … (+2 more) \\
\code{gonotrophic.egg\_batch\_min} & 30 & huevos & mal-core/src/mal\_core/abm/include/mal\_abm\_fast/config.hpp, mal-core/src/mal\_core/abm/include/mal\_abm\_fast/gonotrophic\_cycle.hpp, mal-core/src/mal\_core/abm/src/mosquito\_submodel.cpp \\
\code{gonotrophic.egg\_batch\_n} & 240 & ensayos & mal-core/src/mal\_core/abm/include/mal\_abm\_fast/config.hpp, mal-core/src/mal\_core/abm/include/mal\_abm\_fast/gonotrophic\_cycle.hpp … (+2 more) \\
\code{gonotrophic.feeding\_success\_rate} & 0.825 & fracción & mal-core/src/mal\_core/abm/include/mal\_abm\_fast/config.hpp, mal-core/src/mal\_core/abm/include/mal\_abm\_fast/gonotrophic\_cycle.hpp \\
\code{gonotrophic.first\_cycle\_days} & 4.0 & días & mal-core/src/mal\_core/abm/include/mal\_abm\_fast/config.hpp, mal-core/src/mal\_core/abm/include/mal\_abm\_fast/gonotrophic\_cycle.hpp \\
\code{gonotrophic.resting\_duration\_days} & 0.5 & días & mal-core/src/mal\_core/abm/include/mal\_abm\_fast/config.hpp, mal-core/src/mal\_core/abm/include/mal\_abm\_fast/gonotrophic\_cycle.hpp \\
\code{habitat.activation\_min\_engine} & 0.01 & índice & mal-core/src/mal\_core/abm/include/mal\_abm\_fast/aquatic\_stress.hpp, mal-core/src/mal\_core/abm/include/mal\_abm\_fast/config.hpp … (+3 more) \\
\code{habitat.candidate\_min\_engine} & 0.01 & índice & mal-core/src/mal\_core/abm/include/mal\_abm\_fast/config.hpp, mal-core/src/mal\_core/abm/include/mal\_abm\_fast/coordinator.hpp, mal-core/src/mal\_core/abm/include/mal\_abm\_fast/habitat\_engine.hpp \\
\code{habitat\_types.dynamics.crop\_pluvial} & "pluvial" & categórico & mal-core/src/mal\_core/abm/src/config.cc \\
\code{habitat\_types.dynamics.open\_pluvial} & "pluvial" & categórico & mal-core/src/mal\_core/abm/src/config.cc \\
\code{habitat\_types.dynamics.surface\_persistent} & "persistent" & categórico & mal-core/src/mal\_core/abm/src/config.cc \\
\code{habitat\_types.dynamics.urban\_persistent} & "persistent" & categórico & mal-core/src/mal\_core/abm/src/config.cc \\
\code{habitat\_types.dynamics.wetland\_persistent} & "persistent" & categórico & mal-core/src/mal\_core/abm/src/config.cc \\
\code{habitat\_types.dynamics.wetland\_pluvial} & "pluvial" & categórico & mal-core/src/mal\_core/abm/include/mal\_abm\_fast/config.hpp, mal-core/src/mal\_core/abm/src/habitat\_engine.cpp \\
\code{habitat\_types.types} & "surface\_persistent, wetland\_persistent, urban\_persistent, wetland\_pluvial, open\_pluvial, crop\_pluvial" & --- & mal-core/src/mal\_core/abm/include/mal\_abm\_fast/bite\_ledger.hpp, mal-core/src/mal\_core/abm/include/mal\_abm\_fast/config.hpp … (+5 more) \\
\code{host\_seeking.indoor\_attraction\_add} & 0.72 & adimensional & mal-core/src/mal\_core/abm/include/mal\_abm\_fast/config.hpp, mal-core/src/mal\_core/abm/src/coordinator.cpp \\
\code{host\_seeking.indoor\_fraction\_rural} & 0.3 & fracción & mal-core/src/mal\_core/abm/include/mal\_abm\_fast/config.hpp, mal-core/src/mal\_core/abm/src/host\_landscape.cpp \\
\code{host\_seeking.indoor\_fraction\_urban} & 0.72 & fracción & mal-core/src/mal\_core/abm/include/mal\_abm\_fast/config.hpp, mal-core/src/mal\_core/abm/src/host\_landscape.cpp \\
\code{host\_seeking.starvation\_days} & 7 & días & mal-core/src/mal\_core/abm/include/mal\_abm\_fast/config.hpp, mal-core/src/mal\_core/abm/include/mal\_abm\_fast/wire.hpp \\
\code{host\_seeking.urbanicity\_add} & 0.2 & adimensional & mal-core/src/mal\_core/abm/include/mal\_abm\_fast/config.hpp, mal-core/src/mal\_core/abm/src/coordinator.cpp \\
\code{oviposition.detect\_range\_m} & 2000.0 & m & mal-core/src/mal\_core/abm/include/mal\_abm\_fast/config.hpp, mal-core/src/mal\_core/abm/include/mal\_abm\_fast/oviposition\_seeking.hpp, mal-core/src/mal\_core/abm/src/mosquito\_submodel.cpp \\
\code{oviposition.min\_attraction} & 0.3 & adimensional & mal-core/src/mal\_core/abm/include/mal\_abm\_fast/config.hpp, mal-core/src/mal\_core/abm/include/mal\_abm\_fast/oviposition\_seeking.hpp, mal-core/src/mal\_core/abm/src/oviposition\_seeking.cpp \\
\code{oviposition.scale\_m} & 50.0 & m & mal-core/src/mal\_core/abm/include/mal\_abm\_fast/config.hpp, mal-core/src/mal\_core/abm/include/mal\_abm\_fast/oviposition\_seeking.hpp … (+3 more) \\
\code{oviposition.step\_m} & 50.0 & m & mal-core/src/mal\_core/abm/include/mal\_abm\_fast/config.hpp, mal-core/src/mal\_core/abm/include/mal\_abm\_fast/oviposition\_seeking.hpp … (+2 more) \\
\code{oviposition.timeout\_days} & 7 & días & mal-core/src/mal\_core/abm/include/mal\_abm\_fast/config.hpp, mal-core/src/mal\_core/abm/include/mal\_abm\_fast/wire.hpp \\
\code{pluvial.cold\_start\_z} & 0.0 & adimensional & mal-core/src/mal\_core/abm/include/mal\_abm\_fast/config.hpp, mal-core/src/mal\_core/abm/src/coordinator.cpp \\
\code{pluvial.delta\_days} & 1.0 & días & mal-core/src/mal\_core/abm/include/mal\_abm\_fast/config.hpp, mal-core/src/mal\_core/abm/include/mal\_abm\_fast/habitat\_availability.hpp, mal-core/src/mal\_core/abm/src/coordinator.cpp \\
\code{pluvial.dry\_memory\_days} & 5.0 & días & mal-core/src/mal\_core/abm/include/mal\_abm\_fast/config.hpp, mal-core/src/mal\_core/abm/include/mal\_abm\_fast/coordinator.hpp, mal-core/src/mal\_core/abm/include/mal\_abm\_fast/habitat\_availability.hpp \\
\code{pluvial.rain\_scale\_mm} & 2.5 & mm & mal-core/src/mal\_core/abm/include/mal\_abm\_fast/config.hpp, mal-core/src/mal\_core/abm/include/mal\_abm\_fast/habitat\_availability.hpp \\
\code{rng.aquatic\_seed\_larva\_mortality} & 123 & int & (loaded, never used) \\
\code{rng.aquatic\_seed\_stage\_mortality} & 42 & int & (loaded, never used) \\
\code{rng.aquatic\_seed\_washout} & 789 & int & (loaded, never used) \\
\code{species.arabiensis.activity\_weights} & "0.02, 0.1, 0.8, 0.08" & fracción & mal-core/src/mal\_core/abm/include/mal\_abm\_fast/config.hpp, mal-core/src/mal\_core/abm/include/mal\_abm\_fast/host\_cascade.hpp … (+3 more) \\
\code{species.arabiensis.aestivation\_supported} & true & booleano & (loaded, never used) \\
\code{species.arabiensis.anemotaxis\_deviation\_deg} & 30.0 & ° & mal-core/src/mal\_core/abm/include/mal\_abm\_fast/config.hpp, mal-core/src/mal\_core/abm/include/mal\_abm\_fast/host\_cascade.hpp, mal-core/src/mal\_core/abm/src/host\_cascade.cpp \\
\code{species.arabiensis.anemotaxis\_rule} & "upwind" & categórico & (loaded, never used) \\
\code{species.arabiensis.arrested\_daily\_survival} & 0.986 & 1/día & (loaded, never used) \\
\code{species.arabiensis.blood\_meal\_refractory\_days} & 0.5 & días & (loaded, never used) \\
\code{species.arabiensis.blood\_meal\_size\_ul} & 0.3 & µL & (loaded, never used) \\
\code{species.arabiensis.blood\_meal\_target\_ul} & 0.5 & µL & (loaded, never used) \\
\code{species.arabiensis.casting\_angle\_deg} & 90.0 & ° & mal-core/src/mal\_core/abm/include/mal\_abm\_fast/config.hpp, mal-core/src/mal\_core/abm/include/mal\_abm\_fast/host\_cascade.hpp, mal-core/src/mal\_core/abm/src/host\_cascade.cpp \\
\code{species.arabiensis.casting\_probability} & 0.1 & 1/paso de rastreo & mal-core/src/mal\_core/abm/include/mal\_abm\_fast/config.hpp, mal-core/src/mal\_core/abm/include/mal\_abm\_fast/host\_cascade.hpp, mal-core/src/mal\_core/abm/src/host\_cascade.cpp \\
\code{species.arabiensis.co2\_activation\_radius\_m} & 100.0 & m & mal-core/src/mal\_core/abm/include/mal\_abm\_fast/config.hpp, mal-core/src/mal\_core/abm/include/mal\_abm\_fast/host\_cascade.hpp … (+2 more) \\
\code{species.arabiensis.daily\_move\_capacity\_m} & 590.0 & m & mal-core/src/mal\_core/abm/include/mal\_abm\_fast/config.hpp, mal-core/src/mal\_core/abm/include/mal\_abm\_fast/host\_cascade.hpp … (+3 more) \\
\code{species.arabiensis.directed\_flight\_budget\_h} & 0.4856 & h & mal-core/src/mal\_core/abm/include/mal\_abm\_fast/config.hpp, mal-core/src/mal\_core/abm/include/mal\_abm\_fast/host\_cascade.hpp, mal-core/src/mal\_core/abm/src/host\_cascade.cpp \\
\code{species.arabiensis.eggs\_per\_ul} & 223.0 & huevos/µL & (loaded, never used) \\
\code{species.arabiensis.exploration\_range\_m} & 100.0 & m & mal-core/src/mal\_core/abm/include/mal\_abm\_fast/config.hpp, mal-core/src/mal\_core/abm/include/mal\_abm\_fast/host\_cascade.hpp, mal-core/src/mal\_core/abm/src/host\_cascade.cpp \\
\code{species.arabiensis.exploration\_step\_correlation} & 0.5 & adimensional & mal-core/src/mal\_core/abm/include/mal\_abm\_fast/config.hpp, mal-core/src/mal\_core/abm/include/mal\_abm\_fast/host\_cascade.hpp, mal-core/src/mal\_core/abm/src/host\_cascade.cpp \\
\code{species.arabiensis.exploration\_step\_m} & 50.0 & m & mal-core/src/mal\_core/abm/include/mal\_abm\_fast/config.hpp, mal-core/src/mal\_core/abm/include/mal\_abm\_fast/host\_cascade.hpp, mal-core/src/mal\_core/abm/src/host\_cascade.cpp \\
\code{species.arabiensis.first\_cycle\_target\_ul} & 0.8 & µL & (loaded, never used) \\
\code{species.arabiensis.flight\_speed\_ms} & 0.3375 & m/s & mal-core/src/mal\_core/abm/include/mal\_abm\_fast/config.hpp, mal-core/src/mal\_core/abm/include/mal\_abm\_fast/host\_cascade.hpp, mal-core/src/mal\_core/abm/src/host\_cascade.cpp \\
\code{species.arabiensis.gonotrophic\_arrest\_prob\_dry} & 0.4454 & fracción & (loaded, never used) \\
\code{species.arabiensis.male\_daily\_survival} & 0.773 & 1/día & (loaded, never used) \\
\code{species.arabiensis.male\_ratio} & 0.5 & fracción & (loaded, never used) \\
\code{species.arabiensis.meal\_size\_min\_ul} & 0.1 & µL & (loaded, never used) \\
\code{species.arabiensis.olfactory\_radius\_m} & 50.0 & m & mal-core/src/mal\_core/abm/include/mal\_abm\_fast/config.hpp, mal-core/src/mal\_core/abm/include/mal\_abm\_fast/host\_cascade.hpp, mal-core/src/mal\_core/abm/src/host\_cascade.cpp \\
\code{species.arabiensis.plume\_memory\_s} & 0.7 & s & mal-core/src/mal\_core/abm/include/mal\_abm\_fast/config.hpp, mal-core/src/mal\_core/abm/include/mal\_abm\_fast/host\_cascade.hpp, mal-core/src/mal\_core/abm/src/host\_cascade.cpp \\
\code{species.arabiensis.pref.cattle} & 0.2 & --- & mal-core/src/mal\_core/abm/include/mal\_abm\_fast/config.hpp, mal-core/src/mal\_core/abm/include/mal\_abm\_fast/host\_landscape.hpp … (+6 more) \\
\code{species.arabiensis.pref.goat} & 0.15 & --- & mal-core/src/mal\_core/abm/include/mal\_abm\_fast/config.hpp, mal-core/src/mal\_core/abm/include/mal\_abm\_fast/seeding.hpp … (+3 more) \\
\code{species.arabiensis.pref.human} & 0.5 & --- & mal-core/src/mal\_core/abm/include/mal\_abm\_fast/config.hpp, mal-core/src/mal\_core/abm/include/mal\_abm\_fast/engine.hpp … (+18 more) \\
\code{species.arabiensis.pref.sheep} & 0.1 & --- & mal-core/src/mal\_core/abm/include/mal\_abm\_fast/config.hpp, mal-core/src/mal\_core/abm/include/mal\_abm\_fast/host\_landscape.hpp … (+6 more) \\
\code{species.arabiensis.pref.wildlife} & 0.05 & --- & mal-core/src/mal\_core/abm/include/mal\_abm\_fast/config.hpp, mal-core/src/mal\_core/abm/include/mal\_abm\_fast/seeding.hpp … (+3 more) \\
\code{species.arabiensis.probing\_success\_rate} & 0.7275 & fracción & (loaded, never used) \\
\code{species.arabiensis.salinity\_hi\_tol\_ppt} & 15.0 & psu & mal-core/src/mal\_core/abm/include/mal\_abm\_fast/config.hpp, mal-core/src/mal\_core/abm/include/mal\_abm\_fast/coordinator.hpp … (+3 more) \\
\code{species.arabiensis.salinity\_opt\_ppt} & 0.0 & psu & mal-core/src/mal\_core/abm/include/mal\_abm\_fast/config.hpp, mal-core/src/mal\_core/abm/include/mal\_abm\_fast/species\_params.hpp, mal-core/src/mal\_core/abm/src/species\_params.cpp \\
\code{species.arabiensis.salinity\_slope} & 0.15 & psu⁻¹ & mal-core/src/mal\_core/abm/include/mal\_abm\_fast/config.hpp, mal-core/src/mal\_core/abm/include/mal\_abm\_fast/species\_params.hpp, mal-core/src/mal\_core/abm/src/species\_params.cpp \\
\code{species.arabiensis.terminal\_radius\_m} & 0.15 & m & mal-core/src/mal\_core/abm/include/mal\_abm\_fast/config.hpp, mal-core/src/mal\_core/abm/include/mal\_abm\_fast/host\_cascade.hpp, mal-core/src/mal\_core/abm/src/host\_cascade.cpp \\
\code{species.arabiensis.visual\_radius\_m} & 0.485 & m & mal-core/src/mal\_core/abm/include/mal\_abm\_fast/config.hpp, mal-core/src/mal\_core/abm/include/mal\_abm\_fast/host\_cascade.hpp, mal-core/src/mal\_core/abm/src/host\_cascade.cpp \\
\code{species.arabiensis.weights.crop\_pluvial} & 0.1 & --- & mal-core/src/mal\_core/abm/src/config.cc \\
\code{species.arabiensis.weights.open\_pluvial} & 0.55 & --- & mal-core/src/mal\_core/abm/src/config.cc \\
\code{species.arabiensis.weights.surface\_persistent} & 0.1 & --- & mal-core/src/mal\_core/abm/src/config.cc \\
\code{species.arabiensis.weights.urban\_persistent} & 0.15 & --- & mal-core/src/mal\_core/abm/src/config.cc \\
\code{species.arabiensis.weights.wetland\_persistent} & 0.1 & --- & mal-core/src/mal\_core/abm/src/config.cc \\
\code{species.arabiensis.weights.wetland\_pluvial} & 0.0 & fracción & mal-core/src/mal\_core/abm/include/mal\_abm\_fast/config.hpp, mal-core/src/mal\_core/abm/src/habitat\_engine.cpp \\
\code{species.arabiensis.wind\_block\_threshold\_ms} & 0.3375 & m/s & mal-core/src/mal\_core/abm/include/mal\_abm\_fast/config.hpp, mal-core/src/mal\_core/abm/include/mal\_abm\_fast/host\_cascade.hpp, mal-core/src/mal\_core/abm/src/host\_cascade.cpp \\
\code{species.arabiensis.wind\_reference} & "u100/v100" & string & (loaded, never used) \\
\code{species.coluzzii.activity\_weights} & "0.0, 0.255, 0.48, 0.265" & fracción & mal-core/src/mal\_core/abm/include/mal\_abm\_fast/config.hpp, mal-core/src/mal\_core/abm/include/mal\_abm\_fast/host\_cascade.hpp … (+3 more) \\
\code{species.coluzzii.aestivation\_supported} & true & booleano & (loaded, never used) \\
\code{species.coluzzii.anemotaxis\_deviation\_deg} & 30.0 & ° & mal-core/src/mal\_core/abm/include/mal\_abm\_fast/config.hpp, mal-core/src/mal\_core/abm/include/mal\_abm\_fast/host\_cascade.hpp, mal-core/src/mal\_core/abm/src/host\_cascade.cpp \\
\code{species.coluzzii.anemotaxis\_rule} & "upwind" & categórico & (loaded, never used) \\
\code{species.coluzzii.arrested\_daily\_survival} & 0.986 & 1/día & (loaded, never used) \\
\code{species.coluzzii.blood\_meal\_refractory\_days} & 0.5 & días & (loaded, never used) \\
\code{species.coluzzii.blood\_meal\_size\_ul} & 0.3 & µL & (loaded, never used) \\
\code{species.coluzzii.blood\_meal\_target\_ul} & 0.5 & µL & (loaded, never used) \\
\code{species.coluzzii.casting\_angle\_deg} & 90.0 & ° & mal-core/src/mal\_core/abm/include/mal\_abm\_fast/config.hpp, mal-core/src/mal\_core/abm/include/mal\_abm\_fast/host\_cascade.hpp, mal-core/src/mal\_core/abm/src/host\_cascade.cpp \\
\code{species.coluzzii.casting\_probability} & 0.1 & 1/paso de rastreo & mal-core/src/mal\_core/abm/include/mal\_abm\_fast/config.hpp, mal-core/src/mal\_core/abm/include/mal\_abm\_fast/host\_cascade.hpp, mal-core/src/mal\_core/abm/src/host\_cascade.cpp \\
\code{species.coluzzii.co2\_activation\_radius\_m} & 100.0 & m & mal-core/src/mal\_core/abm/include/mal\_abm\_fast/config.hpp, mal-core/src/mal\_core/abm/include/mal\_abm\_fast/host\_cascade.hpp … (+2 more) \\
\code{species.coluzzii.daily\_move\_capacity\_m} & 590.0 & m & mal-core/src/mal\_core/abm/include/mal\_abm\_fast/config.hpp, mal-core/src/mal\_core/abm/include/mal\_abm\_fast/host\_cascade.hpp … (+3 more) \\
\code{species.coluzzii.directed\_flight\_budget\_h} & 0.4856 & h & mal-core/src/mal\_core/abm/include/mal\_abm\_fast/config.hpp, mal-core/src/mal\_core/abm/include/mal\_abm\_fast/host\_cascade.hpp, mal-core/src/mal\_core/abm/src/host\_cascade.cpp \\
\code{species.coluzzii.eggs\_per\_ul} & 223.0 & huevos/µL & (loaded, never used) \\
\code{species.coluzzii.exploration\_range\_m} & 100.0 & m & mal-core/src/mal\_core/abm/include/mal\_abm\_fast/config.hpp, mal-core/src/mal\_core/abm/include/mal\_abm\_fast/host\_cascade.hpp, mal-core/src/mal\_core/abm/src/host\_cascade.cpp \\
\code{species.coluzzii.exploration\_step\_correlation} & 0.5 & adimensional & mal-core/src/mal\_core/abm/include/mal\_abm\_fast/config.hpp, mal-core/src/mal\_core/abm/include/mal\_abm\_fast/host\_cascade.hpp, mal-core/src/mal\_core/abm/src/host\_cascade.cpp \\
\code{species.coluzzii.exploration\_step\_m} & 50.0 & m & mal-core/src/mal\_core/abm/include/mal\_abm\_fast/config.hpp, mal-core/src/mal\_core/abm/include/mal\_abm\_fast/host\_cascade.hpp, mal-core/src/mal\_core/abm/src/host\_cascade.cpp \\
\code{species.coluzzii.first\_cycle\_target\_ul} & 0.8 & µL & (loaded, never used) \\
\code{species.coluzzii.flight\_speed\_ms} & 0.3375 & m/s & mal-core/src/mal\_core/abm/include/mal\_abm\_fast/config.hpp, mal-core/src/mal\_core/abm/include/mal\_abm\_fast/host\_cascade.hpp, mal-core/src/mal\_core/abm/src/host\_cascade.cpp \\
\code{species.coluzzii.gonotrophic\_arrest\_prob\_dry} & 0.2704 & fracción & (loaded, never used) \\
\code{species.coluzzii.male\_daily\_survival} & 0.773 & 1/día & (loaded, never used) \\
\code{species.coluzzii.male\_ratio} & 0.5 & fracción & (loaded, never used) \\
\code{species.coluzzii.meal\_size\_min\_ul} & 0.1 & µL & (loaded, never used) \\
\code{species.coluzzii.olfactory\_radius\_m} & 50.0 & m & mal-core/src/mal\_core/abm/include/mal\_abm\_fast/config.hpp, mal-core/src/mal\_core/abm/include/mal\_abm\_fast/host\_cascade.hpp, mal-core/src/mal\_core/abm/src/host\_cascade.cpp \\
\code{species.coluzzii.plume\_memory\_s} & 0.7 & s & mal-core/src/mal\_core/abm/include/mal\_abm\_fast/config.hpp, mal-core/src/mal\_core/abm/include/mal\_abm\_fast/host\_cascade.hpp, mal-core/src/mal\_core/abm/src/host\_cascade.cpp \\
\code{species.coluzzii.pref.cattle} & 0.005 & --- & mal-core/src/mal\_core/abm/include/mal\_abm\_fast/config.hpp, mal-core/src/mal\_core/abm/include/mal\_abm\_fast/host\_landscape.hpp … (+6 more) \\
\code{species.coluzzii.pref.goat} & 0.003 & --- & mal-core/src/mal\_core/abm/include/mal\_abm\_fast/config.hpp, mal-core/src/mal\_core/abm/include/mal\_abm\_fast/seeding.hpp … (+3 more) \\
\code{species.coluzzii.pref.human} & 0.99 & --- & mal-core/src/mal\_core/abm/include/mal\_abm\_fast/config.hpp, mal-core/src/mal\_core/abm/include/mal\_abm\_fast/engine.hpp … (+18 more) \\
\code{species.coluzzii.pref.sheep} & 0.001 & --- & mal-core/src/mal\_core/abm/include/mal\_abm\_fast/config.hpp, mal-core/src/mal\_core/abm/include/mal\_abm\_fast/host\_landscape.hpp … (+6 more) \\
\code{species.coluzzii.pref.wildlife} & 0.001 & --- & mal-core/src/mal\_core/abm/include/mal\_abm\_fast/config.hpp, mal-core/src/mal\_core/abm/include/mal\_abm\_fast/seeding.hpp … (+3 more) \\
\code{species.coluzzii.probing\_success\_rate} & 0.7275 & fracción & (loaded, never used) \\
\code{species.coluzzii.salinity\_hi\_tol\_ppt} & 30.0 & psu & mal-core/src/mal\_core/abm/include/mal\_abm\_fast/config.hpp, mal-core/src/mal\_core/abm/include/mal\_abm\_fast/coordinator.hpp … (+3 more) \\
\code{species.coluzzii.salinity\_opt\_ppt} & 0.0 & psu & mal-core/src/mal\_core/abm/include/mal\_abm\_fast/config.hpp, mal-core/src/mal\_core/abm/include/mal\_abm\_fast/species\_params.hpp, mal-core/src/mal\_core/abm/src/species\_params.cpp \\
\code{species.coluzzii.salinity\_slope} & 0.08 & psu⁻¹ & mal-core/src/mal\_core/abm/include/mal\_abm\_fast/config.hpp, mal-core/src/mal\_core/abm/include/mal\_abm\_fast/species\_params.hpp, mal-core/src/mal\_core/abm/src/species\_params.cpp \\
\code{species.coluzzii.terminal\_radius\_m} & 0.15 & m & mal-core/src/mal\_core/abm/include/mal\_abm\_fast/config.hpp, mal-core/src/mal\_core/abm/include/mal\_abm\_fast/host\_cascade.hpp, mal-core/src/mal\_core/abm/src/host\_cascade.cpp \\
\code{species.coluzzii.visual\_radius\_m} & 0.485 & m & mal-core/src/mal\_core/abm/include/mal\_abm\_fast/config.hpp, mal-core/src/mal\_core/abm/include/mal\_abm\_fast/host\_cascade.hpp, mal-core/src/mal\_core/abm/src/host\_cascade.cpp \\
\code{species.coluzzii.weights.crop\_pluvial} & 0.05 & --- & mal-core/src/mal\_core/abm/src/config.cc \\
\code{species.coluzzii.weights.open\_pluvial} & 0.25 & --- & mal-core/src/mal\_core/abm/src/config.cc \\
\code{species.coluzzii.weights.surface\_persistent} & 0.4 & --- & mal-core/src/mal\_core/abm/src/config.cc \\
\code{species.coluzzii.weights.urban\_persistent} & 0.2 & --- & mal-core/src/mal\_core/abm/src/config.cc \\
\code{species.coluzzii.weights.wetland\_persistent} & 0.1 & --- & mal-core/src/mal\_core/abm/src/config.cc \\
\code{species.coluzzii.weights.wetland\_pluvial} & 0.0 & fracción & mal-core/src/mal\_core/abm/include/mal\_abm\_fast/config.hpp, mal-core/src/mal\_core/abm/src/habitat\_engine.cpp \\
\code{species.coluzzii.wind\_block\_threshold\_ms} & 0.3375 & m/s & mal-core/src/mal\_core/abm/include/mal\_abm\_fast/config.hpp, mal-core/src/mal\_core/abm/include/mal\_abm\_fast/host\_cascade.hpp, mal-core/src/mal\_core/abm/src/host\_cascade.cpp \\
\code{species.coluzzii.wind\_reference} & "u100/v100" & string & (loaded, never used) \\
\code{species.default\_species} & "coluzzii" & string & mal-core/src/mal\_core/abm/include/mal\_abm\_fast/config.hpp, mal-core/src/mal\_core/abm/include/mal\_abm\_fast/coordinator.hpp, mal-core/src/mal\_core/abm/include/mal\_abm\_fast/species\_params.hpp \\
\code{species.funestus.activity\_weights} & "0.02, 0.1, 0.8, 0.08" & fracción & mal-core/src/mal\_core/abm/include/mal\_abm\_fast/config.hpp, mal-core/src/mal\_core/abm/include/mal\_abm\_fast/host\_cascade.hpp … (+3 more) \\
\code{species.funestus.aestivation\_supported} & false & booleano & (loaded, never used) \\
\code{species.funestus.anemotaxis\_deviation\_deg} & 30.0 & ° & mal-core/src/mal\_core/abm/include/mal\_abm\_fast/config.hpp, mal-core/src/mal\_core/abm/include/mal\_abm\_fast/host\_cascade.hpp, mal-core/src/mal\_core/abm/src/host\_cascade.cpp \\
\code{species.funestus.anemotaxis\_rule} & "upwind" & categórico & (loaded, never used) \\
\code{species.funestus.arrested\_daily\_survival} & 0.986 & 1/día & (loaded, never used) \\
\code{species.funestus.blood\_meal\_refractory\_days} & 0.5 & días & (loaded, never used) \\
\code{species.funestus.blood\_meal\_size\_ul} & 0.3 & µL & (loaded, never used) \\
\code{species.funestus.blood\_meal\_target\_ul} & 0.5 & µL & (loaded, never used) \\
\code{species.funestus.casting\_angle\_deg} & 90.0 & ° & mal-core/src/mal\_core/abm/include/mal\_abm\_fast/config.hpp, mal-core/src/mal\_core/abm/include/mal\_abm\_fast/host\_cascade.hpp, mal-core/src/mal\_core/abm/src/host\_cascade.cpp \\
\code{species.funestus.casting\_probability} & 0.1 & 1/paso de rastreo & mal-core/src/mal\_core/abm/include/mal\_abm\_fast/config.hpp, mal-core/src/mal\_core/abm/include/mal\_abm\_fast/host\_cascade.hpp, mal-core/src/mal\_core/abm/src/host\_cascade.cpp \\
\code{species.funestus.co2\_activation\_radius\_m} & 100.0 & m & mal-core/src/mal\_core/abm/include/mal\_abm\_fast/config.hpp, mal-core/src/mal\_core/abm/include/mal\_abm\_fast/host\_cascade.hpp … (+2 more) \\
\code{species.funestus.daily\_move\_capacity\_m} & 590.0 & m & mal-core/src/mal\_core/abm/include/mal\_abm\_fast/config.hpp, mal-core/src/mal\_core/abm/include/mal\_abm\_fast/host\_cascade.hpp … (+3 more) \\
\code{species.funestus.directed\_flight\_budget\_h} & 0.4856 & h & mal-core/src/mal\_core/abm/include/mal\_abm\_fast/config.hpp, mal-core/src/mal\_core/abm/include/mal\_abm\_fast/host\_cascade.hpp, mal-core/src/mal\_core/abm/src/host\_cascade.cpp \\
\code{species.funestus.eggs\_per\_ul} & 223.0 & huevos/µL & (loaded, never used) \\
\code{species.funestus.exploration\_range\_m} & 100.0 & m & mal-core/src/mal\_core/abm/include/mal\_abm\_fast/config.hpp, mal-core/src/mal\_core/abm/include/mal\_abm\_fast/host\_cascade.hpp, mal-core/src/mal\_core/abm/src/host\_cascade.cpp \\
\code{species.funestus.exploration\_step\_correlation} & 0.5 & adimensional & mal-core/src/mal\_core/abm/include/mal\_abm\_fast/config.hpp, mal-core/src/mal\_core/abm/include/mal\_abm\_fast/host\_cascade.hpp, mal-core/src/mal\_core/abm/src/host\_cascade.cpp \\
\code{species.funestus.exploration\_step\_m} & 50.0 & m & mal-core/src/mal\_core/abm/include/mal\_abm\_fast/config.hpp, mal-core/src/mal\_core/abm/include/mal\_abm\_fast/host\_cascade.hpp, mal-core/src/mal\_core/abm/src/host\_cascade.cpp \\
\code{species.funestus.first\_cycle\_target\_ul} & 0.8 & µL & (loaded, never used) \\
\code{species.funestus.flight\_speed\_ms} & 0.3375 & m/s & mal-core/src/mal\_core/abm/include/mal\_abm\_fast/config.hpp, mal-core/src/mal\_core/abm/include/mal\_abm\_fast/host\_cascade.hpp, mal-core/src/mal\_core/abm/src/host\_cascade.cpp \\
\code{species.funestus.gonotrophic\_arrest\_prob\_dry} & 0.0 & fracción & (loaded, never used) \\
\code{species.funestus.male\_daily\_survival} & 0.773 & 1/día & (loaded, never used) \\
\code{species.funestus.male\_ratio} & 0.5 & fracción & (loaded, never used) \\
\code{species.funestus.meal\_size\_min\_ul} & 0.1 & µL & (loaded, never used) \\
\code{species.funestus.olfactory\_radius\_m} & 50.0 & m & mal-core/src/mal\_core/abm/include/mal\_abm\_fast/config.hpp, mal-core/src/mal\_core/abm/include/mal\_abm\_fast/host\_cascade.hpp, mal-core/src/mal\_core/abm/src/host\_cascade.cpp \\
\code{species.funestus.plume\_memory\_s} & 0.7 & s & mal-core/src/mal\_core/abm/include/mal\_abm\_fast/config.hpp, mal-core/src/mal\_core/abm/include/mal\_abm\_fast/host\_cascade.hpp, mal-core/src/mal\_core/abm/src/host\_cascade.cpp \\
\code{species.funestus.pref.cattle} & 0.005 & --- & mal-core/src/mal\_core/abm/include/mal\_abm\_fast/config.hpp, mal-core/src/mal\_core/abm/include/mal\_abm\_fast/host\_landscape.hpp … (+6 more) \\
\code{species.funestus.pref.goat} & 0.003 & --- & mal-core/src/mal\_core/abm/include/mal\_abm\_fast/config.hpp, mal-core/src/mal\_core/abm/include/mal\_abm\_fast/seeding.hpp … (+3 more) \\
\code{species.funestus.pref.human} & 0.99 & --- & mal-core/src/mal\_core/abm/include/mal\_abm\_fast/config.hpp, mal-core/src/mal\_core/abm/include/mal\_abm\_fast/engine.hpp … (+18 more) \\
\code{species.funestus.pref.sheep} & 0.001 & --- & mal-core/src/mal\_core/abm/include/mal\_abm\_fast/config.hpp, mal-core/src/mal\_core/abm/include/mal\_abm\_fast/host\_landscape.hpp … (+6 more) \\
\code{species.funestus.pref.wildlife} & 0.001 & --- & mal-core/src/mal\_core/abm/include/mal\_abm\_fast/config.hpp, mal-core/src/mal\_core/abm/include/mal\_abm\_fast/seeding.hpp … (+3 more) \\
\code{species.funestus.probing\_success\_rate} & 0.7275 & fracción & (loaded, never used) \\
\code{species.funestus.salinity\_hi\_tol\_ppt} & 15.0 & psu & mal-core/src/mal\_core/abm/include/mal\_abm\_fast/config.hpp, mal-core/src/mal\_core/abm/include/mal\_abm\_fast/coordinator.hpp … (+3 more) \\
\code{species.funestus.salinity\_opt\_ppt} & 0.0 & psu & mal-core/src/mal\_core/abm/include/mal\_abm\_fast/config.hpp, mal-core/src/mal\_core/abm/include/mal\_abm\_fast/species\_params.hpp, mal-core/src/mal\_core/abm/src/species\_params.cpp \\
\code{species.funestus.salinity\_slope} & 0.15 & psu⁻¹ & mal-core/src/mal\_core/abm/include/mal\_abm\_fast/config.hpp, mal-core/src/mal\_core/abm/include/mal\_abm\_fast/species\_params.hpp, mal-core/src/mal\_core/abm/src/species\_params.cpp \\
\code{species.funestus.terminal\_radius\_m} & 0.15 & m & mal-core/src/mal\_core/abm/include/mal\_abm\_fast/config.hpp, mal-core/src/mal\_core/abm/include/mal\_abm\_fast/host\_cascade.hpp, mal-core/src/mal\_core/abm/src/host\_cascade.cpp \\
\code{species.funestus.visual\_radius\_m} & 0.485 & m & mal-core/src/mal\_core/abm/include/mal\_abm\_fast/config.hpp, mal-core/src/mal\_core/abm/include/mal\_abm\_fast/host\_cascade.hpp, mal-core/src/mal\_core/abm/src/host\_cascade.cpp \\
\code{species.funestus.weights.crop\_pluvial} & 0.05 & --- & mal-core/src/mal\_core/abm/src/config.cc \\
\code{species.funestus.weights.open\_pluvial} & 0.1 & --- & mal-core/src/mal\_core/abm/src/config.cc \\
\code{species.funestus.weights.surface\_persistent} & 0.45 & --- & mal-core/src/mal\_core/abm/src/config.cc \\
\code{species.funestus.weights.urban\_persistent} & 0.1 & --- & mal-core/src/mal\_core/abm/src/config.cc \\
\code{species.funestus.weights.wetland\_persistent} & 0.3 & --- & mal-core/src/mal\_core/abm/src/config.cc \\
\code{species.funestus.weights.wetland\_pluvial} & 0.0 & fracción & mal-core/src/mal\_core/abm/include/mal\_abm\_fast/config.hpp, mal-core/src/mal\_core/abm/src/habitat\_engine.cpp \\
\code{species.funestus.wind\_block\_threshold\_ms} & 0.3375 & m/s & mal-core/src/mal\_core/abm/include/mal\_abm\_fast/config.hpp, mal-core/src/mal\_core/abm/include/mal\_abm\_fast/host\_cascade.hpp, mal-core/src/mal\_core/abm/src/host\_cascade.cpp \\
\code{species.funestus.wind\_reference} & "u100/v100" & string & (loaded, never used) \\
\code{species.gambiae\_ss.activity\_weights} & "0.02, 0.1, 0.8, 0.08" & fracción & mal-core/src/mal\_core/abm/include/mal\_abm\_fast/config.hpp, mal-core/src/mal\_core/abm/include/mal\_abm\_fast/host\_cascade.hpp … (+3 more) \\
\code{species.gambiae\_ss.aestivation\_supported} & true & booleano & (loaded, never used) \\
\code{species.gambiae\_ss.anemotaxis\_deviation\_deg} & 30.0 & ° & mal-core/src/mal\_core/abm/include/mal\_abm\_fast/config.hpp, mal-core/src/mal\_core/abm/include/mal\_abm\_fast/host\_cascade.hpp, mal-core/src/mal\_core/abm/src/host\_cascade.cpp \\
\code{species.gambiae\_ss.anemotaxis\_rule} & "upwind" & categórico & (loaded, never used) \\
\code{species.gambiae\_ss.arrested\_daily\_survival} & 0.986 & 1/día & (loaded, never used) \\
\code{species.gambiae\_ss.blood\_meal\_refractory\_days} & 0.5 & días & (loaded, never used) \\
\code{species.gambiae\_ss.blood\_meal\_size\_ul} & 0.3 & µL & (loaded, never used) \\
\code{species.gambiae\_ss.blood\_meal\_target\_ul} & 0.5 & µL & (loaded, never used) \\
\code{species.gambiae\_ss.casting\_angle\_deg} & 90.0 & ° & mal-core/src/mal\_core/abm/include/mal\_abm\_fast/config.hpp, mal-core/src/mal\_core/abm/include/mal\_abm\_fast/host\_cascade.hpp, mal-core/src/mal\_core/abm/src/host\_cascade.cpp \\
\code{species.gambiae\_ss.casting\_probability} & 0.1 & 1/paso de rastreo & mal-core/src/mal\_core/abm/include/mal\_abm\_fast/config.hpp, mal-core/src/mal\_core/abm/include/mal\_abm\_fast/host\_cascade.hpp, mal-core/src/mal\_core/abm/src/host\_cascade.cpp \\
\code{species.gambiae\_ss.co2\_activation\_radius\_m} & 100.0 & m & mal-core/src/mal\_core/abm/include/mal\_abm\_fast/config.hpp, mal-core/src/mal\_core/abm/include/mal\_abm\_fast/host\_cascade.hpp … (+2 more) \\
\code{species.gambiae\_ss.daily\_move\_capacity\_m} & 590.0 & m & mal-core/src/mal\_core/abm/include/mal\_abm\_fast/config.hpp, mal-core/src/mal\_core/abm/include/mal\_abm\_fast/host\_cascade.hpp … (+3 more) \\
\code{species.gambiae\_ss.directed\_flight\_budget\_h} & 0.4856 & h & mal-core/src/mal\_core/abm/include/mal\_abm\_fast/config.hpp, mal-core/src/mal\_core/abm/include/mal\_abm\_fast/host\_cascade.hpp, mal-core/src/mal\_core/abm/src/host\_cascade.cpp \\
\code{species.gambiae\_ss.eggs\_per\_ul} & 223.0 & huevos/µL & (loaded, never used) \\
\code{species.gambiae\_ss.exploration\_range\_m} & 100.0 & m & mal-core/src/mal\_core/abm/include/mal\_abm\_fast/config.hpp, mal-core/src/mal\_core/abm/include/mal\_abm\_fast/host\_cascade.hpp, mal-core/src/mal\_core/abm/src/host\_cascade.cpp \\
\code{species.gambiae\_ss.exploration\_step\_correlation} & 0.5 & adimensional & mal-core/src/mal\_core/abm/include/mal\_abm\_fast/config.hpp, mal-core/src/mal\_core/abm/include/mal\_abm\_fast/host\_cascade.hpp, mal-core/src/mal\_core/abm/src/host\_cascade.cpp \\
\code{species.gambiae\_ss.exploration\_step\_m} & 50.0 & m & mal-core/src/mal\_core/abm/include/mal\_abm\_fast/config.hpp, mal-core/src/mal\_core/abm/include/mal\_abm\_fast/host\_cascade.hpp, mal-core/src/mal\_core/abm/src/host\_cascade.cpp \\
\code{species.gambiae\_ss.first\_cycle\_target\_ul} & 0.8 & µL & (loaded, never used) \\
\code{species.gambiae\_ss.flight\_speed\_ms} & 0.3375 & m/s & mal-core/src/mal\_core/abm/include/mal\_abm\_fast/config.hpp, mal-core/src/mal\_core/abm/include/mal\_abm\_fast/host\_cascade.hpp, mal-core/src/mal\_core/abm/src/host\_cascade.cpp \\
\code{species.gambiae\_ss.gonotrophic\_arrest\_prob\_dry} & 0.6838 & fracción & (loaded, never used) \\
\code{species.gambiae\_ss.male\_daily\_survival} & 0.773 & 1/día & (loaded, never used) \\
\code{species.gambiae\_ss.male\_ratio} & 0.5 & fracción & (loaded, never used) \\
\code{species.gambiae\_ss.meal\_size\_min\_ul} & 0.1 & µL & (loaded, never used) \\
\code{species.gambiae\_ss.olfactory\_radius\_m} & 50.0 & m & mal-core/src/mal\_core/abm/include/mal\_abm\_fast/config.hpp, mal-core/src/mal\_core/abm/include/mal\_abm\_fast/host\_cascade.hpp, mal-core/src/mal\_core/abm/src/host\_cascade.cpp \\
\code{species.gambiae\_ss.plume\_memory\_s} & 0.7 & s & mal-core/src/mal\_core/abm/include/mal\_abm\_fast/config.hpp, mal-core/src/mal\_core/abm/include/mal\_abm\_fast/host\_cascade.hpp, mal-core/src/mal\_core/abm/src/host\_cascade.cpp \\
\code{species.gambiae\_ss.pref.cattle} & 0.005 & --- & mal-core/src/mal\_core/abm/include/mal\_abm\_fast/config.hpp, mal-core/src/mal\_core/abm/include/mal\_abm\_fast/host\_landscape.hpp … (+6 more) \\
\code{species.gambiae\_ss.pref.goat} & 0.003 & --- & mal-core/src/mal\_core/abm/include/mal\_abm\_fast/config.hpp, mal-core/src/mal\_core/abm/include/mal\_abm\_fast/seeding.hpp … (+3 more) \\
\code{species.gambiae\_ss.pref.human} & 0.99 & --- & mal-core/src/mal\_core/abm/include/mal\_abm\_fast/config.hpp, mal-core/src/mal\_core/abm/include/mal\_abm\_fast/engine.hpp … (+18 more) \\
\code{species.gambiae\_ss.pref.sheep} & 0.001 & --- & mal-core/src/mal\_core/abm/include/mal\_abm\_fast/config.hpp, mal-core/src/mal\_core/abm/include/mal\_abm\_fast/host\_landscape.hpp … (+6 more) \\
\code{species.gambiae\_ss.pref.wildlife} & 0.001 & --- & mal-core/src/mal\_core/abm/include/mal\_abm\_fast/config.hpp, mal-core/src/mal\_core/abm/include/mal\_abm\_fast/seeding.hpp … (+3 more) \\
\code{species.gambiae\_ss.probing\_success\_rate} & 0.7275 & fracción & (loaded, never used) \\
\code{species.gambiae\_ss.salinity\_hi\_tol\_ppt} & 20.0 & psu & mal-core/src/mal\_core/abm/include/mal\_abm\_fast/config.hpp, mal-core/src/mal\_core/abm/include/mal\_abm\_fast/coordinator.hpp … (+3 more) \\
\code{species.gambiae\_ss.salinity\_opt\_ppt} & 0.0 & psu & mal-core/src/mal\_core/abm/include/mal\_abm\_fast/config.hpp, mal-core/src/mal\_core/abm/include/mal\_abm\_fast/species\_params.hpp, mal-core/src/mal\_core/abm/src/species\_params.cpp \\
\code{species.gambiae\_ss.salinity\_slope} & 0.12 & psu⁻¹ & mal-core/src/mal\_core/abm/include/mal\_abm\_fast/config.hpp, mal-core/src/mal\_core/abm/include/mal\_abm\_fast/species\_params.hpp, mal-core/src/mal\_core/abm/src/species\_params.cpp \\
\code{species.gambiae\_ss.terminal\_radius\_m} & 0.15 & m & mal-core/src/mal\_core/abm/include/mal\_abm\_fast/config.hpp, mal-core/src/mal\_core/abm/include/mal\_abm\_fast/host\_cascade.hpp, mal-core/src/mal\_core/abm/src/host\_cascade.cpp \\
\code{species.gambiae\_ss.visual\_radius\_m} & 0.485 & m & mal-core/src/mal\_core/abm/include/mal\_abm\_fast/config.hpp, mal-core/src/mal\_core/abm/include/mal\_abm\_fast/host\_cascade.hpp, mal-core/src/mal\_core/abm/src/host\_cascade.cpp \\
\code{species.gambiae\_ss.weights.crop\_pluvial} & 0.1 & --- & mal-core/src/mal\_core/abm/src/config.cc \\
\code{species.gambiae\_ss.weights.open\_pluvial} & 0.65 & --- & mal-core/src/mal\_core/abm/src/config.cc \\
\code{species.gambiae\_ss.weights.surface\_persistent} & 0.05 & --- & mal-core/src/mal\_core/abm/src/config.cc \\
\code{species.gambiae\_ss.weights.urban\_persistent} & 0.15 & --- & mal-core/src/mal\_core/abm/src/config.cc \\
\code{species.gambiae\_ss.weights.wetland\_persistent} & 0.05 & --- & mal-core/src/mal\_core/abm/src/config.cc \\
\code{species.gambiae\_ss.weights.wetland\_pluvial} & 0.0 & fracción & mal-core/src/mal\_core/abm/include/mal\_abm\_fast/config.hpp, mal-core/src/mal\_core/abm/src/habitat\_engine.cpp \\
\code{species.gambiae\_ss.wind\_block\_threshold\_ms} & 0.3375 & m/s & mal-core/src/mal\_core/abm/include/mal\_abm\_fast/config.hpp, mal-core/src/mal\_core/abm/include/mal\_abm\_fast/host\_cascade.hpp, mal-core/src/mal\_core/abm/src/host\_cascade.cpp \\
\code{species.gambiae\_ss.wind\_reference} & "u100/v100" & string & (loaded, never used) \\
\code{species.k\_ref} & 50.0 & ind·km⁻²·d⁻¹ & mal-core/src/mal\_core/abm/include/mal\_abm\_fast/config.hpp, mal-core/src/mal\_core/abm/include/mal\_abm\_fast/habitat\_availability.hpp \\
\code{species.melas.activity\_weights} & "0.02, 0.1, 0.8, 0.08" & fracción & mal-core/src/mal\_core/abm/include/mal\_abm\_fast/config.hpp, mal-core/src/mal\_core/abm/include/mal\_abm\_fast/host\_cascade.hpp … (+3 more) \\
\code{species.melas.aestivation\_supported} & false & booleano & (loaded, never used) \\
\code{species.melas.anemotaxis\_deviation\_deg} & 30.0 & ° & mal-core/src/mal\_core/abm/include/mal\_abm\_fast/config.hpp, mal-core/src/mal\_core/abm/include/mal\_abm\_fast/host\_cascade.hpp, mal-core/src/mal\_core/abm/src/host\_cascade.cpp \\
\code{species.melas.anemotaxis\_rule} & "upwind" & categórico & (loaded, never used) \\
\code{species.melas.arrested\_daily\_survival} & 0.986 & 1/día & (loaded, never used) \\
\code{species.melas.blood\_meal\_refractory\_days} & 0.5 & días & (loaded, never used) \\
\code{species.melas.blood\_meal\_size\_ul} & 0.3 & µL & (loaded, never used) \\
\code{species.melas.blood\_meal\_target\_ul} & 0.5 & µL & (loaded, never used) \\
\code{species.melas.casting\_angle\_deg} & 90.0 & ° & mal-core/src/mal\_core/abm/include/mal\_abm\_fast/config.hpp, mal-core/src/mal\_core/abm/include/mal\_abm\_fast/host\_cascade.hpp, mal-core/src/mal\_core/abm/src/host\_cascade.cpp \\
\code{species.melas.casting\_probability} & 0.1 & 1/paso de rastreo & mal-core/src/mal\_core/abm/include/mal\_abm\_fast/config.hpp, mal-core/src/mal\_core/abm/include/mal\_abm\_fast/host\_cascade.hpp, mal-core/src/mal\_core/abm/src/host\_cascade.cpp \\
\code{species.melas.co2\_activation\_radius\_m} & 100.0 & m & mal-core/src/mal\_core/abm/include/mal\_abm\_fast/config.hpp, mal-core/src/mal\_core/abm/include/mal\_abm\_fast/host\_cascade.hpp … (+2 more) \\
\code{species.melas.daily\_move\_capacity\_m} & 590.0 & m & mal-core/src/mal\_core/abm/include/mal\_abm\_fast/config.hpp, mal-core/src/mal\_core/abm/include/mal\_abm\_fast/host\_cascade.hpp … (+3 more) \\
\code{species.melas.directed\_flight\_budget\_h} & 0.4856 & h & mal-core/src/mal\_core/abm/include/mal\_abm\_fast/config.hpp, mal-core/src/mal\_core/abm/include/mal\_abm\_fast/host\_cascade.hpp, mal-core/src/mal\_core/abm/src/host\_cascade.cpp \\
\code{species.melas.eggs\_per\_ul} & 223.0 & huevos/µL & (loaded, never used) \\
\code{species.melas.exploration\_range\_m} & 100.0 & m & mal-core/src/mal\_core/abm/include/mal\_abm\_fast/config.hpp, mal-core/src/mal\_core/abm/include/mal\_abm\_fast/host\_cascade.hpp, mal-core/src/mal\_core/abm/src/host\_cascade.cpp \\
\code{species.melas.exploration\_step\_correlation} & 0.5 & adimensional & mal-core/src/mal\_core/abm/include/mal\_abm\_fast/config.hpp, mal-core/src/mal\_core/abm/include/mal\_abm\_fast/host\_cascade.hpp, mal-core/src/mal\_core/abm/src/host\_cascade.cpp \\
\code{species.melas.exploration\_step\_m} & 50.0 & m & mal-core/src/mal\_core/abm/include/mal\_abm\_fast/config.hpp, mal-core/src/mal\_core/abm/include/mal\_abm\_fast/host\_cascade.hpp, mal-core/src/mal\_core/abm/src/host\_cascade.cpp \\
\code{species.melas.first\_cycle\_target\_ul} & 0.8 & µL & (loaded, never used) \\
\code{species.melas.flight\_speed\_ms} & 0.3375 & m/s & mal-core/src/mal\_core/abm/include/mal\_abm\_fast/config.hpp, mal-core/src/mal\_core/abm/include/mal\_abm\_fast/host\_cascade.hpp, mal-core/src/mal\_core/abm/src/host\_cascade.cpp \\
\code{species.melas.gonotrophic\_arrest\_prob\_dry} & 0.0 & fracción & (loaded, never used) \\
\code{species.melas.male\_daily\_survival} & 0.773 & 1/día & (loaded, never used) \\
\code{species.melas.male\_ratio} & 0.5 & fracción & (loaded, never used) \\
\code{species.melas.meal\_size\_min\_ul} & 0.1 & µL & (loaded, never used) \\
\code{species.melas.olfactory\_radius\_m} & 50.0 & m & mal-core/src/mal\_core/abm/include/mal\_abm\_fast/config.hpp, mal-core/src/mal\_core/abm/include/mal\_abm\_fast/host\_cascade.hpp, mal-core/src/mal\_core/abm/src/host\_cascade.cpp \\
\code{species.melas.plume\_memory\_s} & 0.7 & s & mal-core/src/mal\_core/abm/include/mal\_abm\_fast/config.hpp, mal-core/src/mal\_core/abm/include/mal\_abm\_fast/host\_cascade.hpp, mal-core/src/mal\_core/abm/src/host\_cascade.cpp \\
\code{species.melas.pref.cattle} & 0.005 & --- & mal-core/src/mal\_core/abm/include/mal\_abm\_fast/config.hpp, mal-core/src/mal\_core/abm/include/mal\_abm\_fast/host\_landscape.hpp … (+6 more) \\
\code{species.melas.pref.goat} & 0.003 & --- & mal-core/src/mal\_core/abm/include/mal\_abm\_fast/config.hpp, mal-core/src/mal\_core/abm/include/mal\_abm\_fast/seeding.hpp … (+3 more) \\
\code{species.melas.pref.human} & 0.99 & --- & mal-core/src/mal\_core/abm/include/mal\_abm\_fast/config.hpp, mal-core/src/mal\_core/abm/include/mal\_abm\_fast/engine.hpp … (+18 more) \\
\code{species.melas.pref.sheep} & 0.001 & --- & mal-core/src/mal\_core/abm/include/mal\_abm\_fast/config.hpp, mal-core/src/mal\_core/abm/include/mal\_abm\_fast/host\_landscape.hpp … (+6 more) \\
\code{species.melas.pref.wildlife} & 0.001 & --- & mal-core/src/mal\_core/abm/include/mal\_abm\_fast/config.hpp, mal-core/src/mal\_core/abm/include/mal\_abm\_fast/seeding.hpp … (+3 more) \\
\code{species.melas.probing\_success\_rate} & 0.7275 & fracción & (loaded, never used) \\
\code{species.melas.salinity\_hi\_tol\_ppt} & 38.0 & psu & mal-core/src/mal\_core/abm/include/mal\_abm\_fast/config.hpp, mal-core/src/mal\_core/abm/include/mal\_abm\_fast/coordinator.hpp … (+3 more) \\
\code{species.melas.salinity\_opt\_ppt} & 0.0 & psu & mal-core/src/mal\_core/abm/include/mal\_abm\_fast/config.hpp, mal-core/src/mal\_core/abm/include/mal\_abm\_fast/species\_params.hpp, mal-core/src/mal\_core/abm/src/species\_params.cpp \\
\code{species.melas.salinity\_slope} & 0.02 & psu⁻¹ & mal-core/src/mal\_core/abm/include/mal\_abm\_fast/config.hpp, mal-core/src/mal\_core/abm/include/mal\_abm\_fast/species\_params.hpp, mal-core/src/mal\_core/abm/src/species\_params.cpp \\
\code{species.melas.terminal\_radius\_m} & 0.15 & m & mal-core/src/mal\_core/abm/include/mal\_abm\_fast/config.hpp, mal-core/src/mal\_core/abm/include/mal\_abm\_fast/host\_cascade.hpp, mal-core/src/mal\_core/abm/src/host\_cascade.cpp \\
\code{species.melas.visual\_radius\_m} & 0.485 & m & mal-core/src/mal\_core/abm/include/mal\_abm\_fast/config.hpp, mal-core/src/mal\_core/abm/include/mal\_abm\_fast/host\_cascade.hpp, mal-core/src/mal\_core/abm/src/host\_cascade.cpp \\
\code{species.melas.weights.crop\_pluvial} & 0.05 & --- & mal-core/src/mal\_core/abm/src/config.cc \\
\code{species.melas.weights.open\_pluvial} & 0.4 & --- & mal-core/src/mal\_core/abm/src/config.cc \\
\code{species.melas.weights.surface\_persistent} & 0.25 & --- & mal-core/src/mal\_core/abm/src/config.cc \\
\code{species.melas.weights.urban\_persistent} & 0.15 & --- & mal-core/src/mal\_core/abm/src/config.cc \\
\code{species.melas.weights.wetland\_persistent} & 0.15 & --- & mal-core/src/mal\_core/abm/src/config.cc \\
\code{species.melas.weights.wetland\_pluvial} & 0.0 & fracción & mal-core/src/mal\_core/abm/include/mal\_abm\_fast/config.hpp, mal-core/src/mal\_core/abm/src/habitat\_engine.cpp \\
\code{species.melas.wind\_block\_threshold\_ms} & 0.3375 & m/s & mal-core/src/mal\_core/abm/include/mal\_abm\_fast/config.hpp, mal-core/src/mal\_core/abm/include/mal\_abm\_fast/host\_cascade.hpp, mal-core/src/mal\_core/abm/src/host\_cascade.cpp \\
\code{species.melas.wind\_reference} & "u100/v100" & string & (loaded, never used) \\
\code{species.neutral.weights.crop\_pluvial} & 0.05 & --- & mal-core/src/mal\_core/abm/src/config.cc \\
\code{species.neutral.weights.open\_pluvial} & 0.4 & --- & mal-core/src/mal\_core/abm/src/config.cc \\
\code{species.neutral.weights.surface\_persistent} & 0.25 & --- & mal-core/src/mal\_core/abm/src/config.cc \\
\code{species.neutral.weights.urban\_persistent} & 0.15 & --- & mal-core/src/mal\_core/abm/src/config.cc \\
\code{species.neutral.weights.wetland\_persistent} & 0.15 & --- & mal-core/src/mal\_core/abm/src/config.cc \\
\code{species.neutral.weights.wetland\_pluvial} & 0.0 & fracción & mal-core/src/mal\_core/abm/include/mal\_abm\_fast/config.hpp, mal-core/src/mal\_core/abm/src/habitat\_engine.cpp \\
\code{species.stephensi.activity\_weights} & "0.02, 0.1, 0.8, 0.08" & fracción & mal-core/src/mal\_core/abm/include/mal\_abm\_fast/config.hpp, mal-core/src/mal\_core/abm/include/mal\_abm\_fast/host\_cascade.hpp … (+3 more) \\
\code{species.stephensi.aestivation\_supported} & false & booleano & (loaded, never used) \\
\code{species.stephensi.anemotaxis\_deviation\_deg} & 30.0 & ° & mal-core/src/mal\_core/abm/include/mal\_abm\_fast/config.hpp, mal-core/src/mal\_core/abm/include/mal\_abm\_fast/host\_cascade.hpp, mal-core/src/mal\_core/abm/src/host\_cascade.cpp \\
\code{species.stephensi.anemotaxis\_rule} & "upwind" & categórico & (loaded, never used) \\
\code{species.stephensi.arrested\_daily\_survival} & 0.986 & 1/día & (loaded, never used) \\
\code{species.stephensi.blood\_meal\_refractory\_days} & 0.5 & días & (loaded, never used) \\
\code{species.stephensi.blood\_meal\_size\_ul} & 0.3 & µL & (loaded, never used) \\
\code{species.stephensi.blood\_meal\_target\_ul} & 0.5 & µL & (loaded, never used) \\
\code{species.stephensi.casting\_angle\_deg} & 90.0 & ° & mal-core/src/mal\_core/abm/include/mal\_abm\_fast/config.hpp, mal-core/src/mal\_core/abm/include/mal\_abm\_fast/host\_cascade.hpp, mal-core/src/mal\_core/abm/src/host\_cascade.cpp \\
\code{species.stephensi.casting\_probability} & 0.1 & 1/paso de rastreo & mal-core/src/mal\_core/abm/include/mal\_abm\_fast/config.hpp, mal-core/src/mal\_core/abm/include/mal\_abm\_fast/host\_cascade.hpp, mal-core/src/mal\_core/abm/src/host\_cascade.cpp \\
\code{species.stephensi.co2\_activation\_radius\_m} & 100.0 & m & mal-core/src/mal\_core/abm/include/mal\_abm\_fast/config.hpp, mal-core/src/mal\_core/abm/include/mal\_abm\_fast/host\_cascade.hpp … (+2 more) \\
\code{species.stephensi.daily\_move\_capacity\_m} & 590.0 & m & mal-core/src/mal\_core/abm/include/mal\_abm\_fast/config.hpp, mal-core/src/mal\_core/abm/include/mal\_abm\_fast/host\_cascade.hpp … (+3 more) \\
\code{species.stephensi.directed\_flight\_budget\_h} & 0.4856 & h & mal-core/src/mal\_core/abm/include/mal\_abm\_fast/config.hpp, mal-core/src/mal\_core/abm/include/mal\_abm\_fast/host\_cascade.hpp, mal-core/src/mal\_core/abm/src/host\_cascade.cpp \\
\code{species.stephensi.eggs\_per\_ul} & 223.0 & huevos/µL & (loaded, never used) \\
\code{species.stephensi.exploration\_range\_m} & 100.0 & m & mal-core/src/mal\_core/abm/include/mal\_abm\_fast/config.hpp, mal-core/src/mal\_core/abm/include/mal\_abm\_fast/host\_cascade.hpp, mal-core/src/mal\_core/abm/src/host\_cascade.cpp \\
\code{species.stephensi.exploration\_step\_correlation} & 0.5 & adimensional & mal-core/src/mal\_core/abm/include/mal\_abm\_fast/config.hpp, mal-core/src/mal\_core/abm/include/mal\_abm\_fast/host\_cascade.hpp, mal-core/src/mal\_core/abm/src/host\_cascade.cpp \\
\code{species.stephensi.exploration\_step\_m} & 50.0 & m & mal-core/src/mal\_core/abm/include/mal\_abm\_fast/config.hpp, mal-core/src/mal\_core/abm/include/mal\_abm\_fast/host\_cascade.hpp, mal-core/src/mal\_core/abm/src/host\_cascade.cpp \\
\code{species.stephensi.first\_cycle\_target\_ul} & 0.8 & µL & (loaded, never used) \\
\code{species.stephensi.flight\_speed\_ms} & 0.3375 & m/s & mal-core/src/mal\_core/abm/include/mal\_abm\_fast/config.hpp, mal-core/src/mal\_core/abm/include/mal\_abm\_fast/host\_cascade.hpp, mal-core/src/mal\_core/abm/src/host\_cascade.cpp \\
\code{species.stephensi.gonotrophic\_arrest\_prob\_dry} & 0.0 & fracción & (loaded, never used) \\
\code{species.stephensi.male\_daily\_survival} & 0.773 & 1/día & (loaded, never used) \\
\code{species.stephensi.male\_ratio} & 0.5 & fracción & (loaded, never used) \\
\code{species.stephensi.meal\_size\_min\_ul} & 0.1 & µL & (loaded, never used) \\
\code{species.stephensi.olfactory\_radius\_m} & 50.0 & m & mal-core/src/mal\_core/abm/include/mal\_abm\_fast/config.hpp, mal-core/src/mal\_core/abm/include/mal\_abm\_fast/host\_cascade.hpp, mal-core/src/mal\_core/abm/src/host\_cascade.cpp \\
\code{species.stephensi.plume\_memory\_s} & 0.7 & s & mal-core/src/mal\_core/abm/include/mal\_abm\_fast/config.hpp, mal-core/src/mal\_core/abm/include/mal\_abm\_fast/host\_cascade.hpp, mal-core/src/mal\_core/abm/src/host\_cascade.cpp \\
\code{species.stephensi.pref.cattle} & 0.15 & --- & mal-core/src/mal\_core/abm/include/mal\_abm\_fast/config.hpp, mal-core/src/mal\_core/abm/include/mal\_abm\_fast/host\_landscape.hpp … (+6 more) \\
\code{species.stephensi.pref.goat} & 0.05 & --- & mal-core/src/mal\_core/abm/include/mal\_abm\_fast/config.hpp, mal-core/src/mal\_core/abm/include/mal\_abm\_fast/seeding.hpp … (+3 more) \\
\code{species.stephensi.pref.human} & 0.7 & --- & mal-core/src/mal\_core/abm/include/mal\_abm\_fast/config.hpp, mal-core/src/mal\_core/abm/include/mal\_abm\_fast/engine.hpp … (+18 more) \\
\code{species.stephensi.pref.sheep} & 0.05 & --- & mal-core/src/mal\_core/abm/include/mal\_abm\_fast/config.hpp, mal-core/src/mal\_core/abm/include/mal\_abm\_fast/host\_landscape.hpp … (+6 more) \\
\code{species.stephensi.pref.wildlife} & 0.05 & --- & mal-core/src/mal\_core/abm/include/mal\_abm\_fast/config.hpp, mal-core/src/mal\_core/abm/include/mal\_abm\_fast/seeding.hpp … (+3 more) \\
\code{species.stephensi.probing\_success\_rate} & 0.7275 & fracción & (loaded, never used) \\
\code{species.stephensi.salinity\_hi\_tol\_ppt} & 20.0 & psu & mal-core/src/mal\_core/abm/include/mal\_abm\_fast/config.hpp, mal-core/src/mal\_core/abm/include/mal\_abm\_fast/coordinator.hpp … (+3 more) \\
\code{species.stephensi.salinity\_opt\_ppt} & 0.0 & psu & mal-core/src/mal\_core/abm/include/mal\_abm\_fast/config.hpp, mal-core/src/mal\_core/abm/include/mal\_abm\_fast/species\_params.hpp, mal-core/src/mal\_core/abm/src/species\_params.cpp \\
\code{species.stephensi.salinity\_slope} & 0.15 & psu⁻¹ & mal-core/src/mal\_core/abm/include/mal\_abm\_fast/config.hpp, mal-core/src/mal\_core/abm/include/mal\_abm\_fast/species\_params.hpp, mal-core/src/mal\_core/abm/src/species\_params.cpp \\
\code{species.stephensi.terminal\_radius\_m} & 0.15 & m & mal-core/src/mal\_core/abm/include/mal\_abm\_fast/config.hpp, mal-core/src/mal\_core/abm/include/mal\_abm\_fast/host\_cascade.hpp, mal-core/src/mal\_core/abm/src/host\_cascade.cpp \\
\code{species.stephensi.visual\_radius\_m} & 0.485 & m & mal-core/src/mal\_core/abm/include/mal\_abm\_fast/config.hpp, mal-core/src/mal\_core/abm/include/mal\_abm\_fast/host\_cascade.hpp, mal-core/src/mal\_core/abm/src/host\_cascade.cpp \\
\code{species.stephensi.weights.crop\_pluvial} & 0.05 & --- & mal-core/src/mal\_core/abm/src/config.cc \\
\code{species.stephensi.weights.open\_pluvial} & 0.4 & --- & mal-core/src/mal\_core/abm/src/config.cc \\
\code{species.stephensi.weights.surface\_persistent} & 0.25 & --- & mal-core/src/mal\_core/abm/src/config.cc \\
\code{species.stephensi.weights.urban\_persistent} & 0.15 & --- & mal-core/src/mal\_core/abm/src/config.cc \\
\code{species.stephensi.weights.wetland\_persistent} & 0.15 & --- & mal-core/src/mal\_core/abm/src/config.cc \\
\code{species.stephensi.weights.wetland\_pluvial} & 0.0 & fracción & mal-core/src/mal\_core/abm/include/mal\_abm\_fast/config.hpp, mal-core/src/mal\_core/abm/src/habitat\_engine.cpp \\
\code{species.stephensi.wind\_block\_threshold\_ms} & 0.3375 & m/s & mal-core/src/mal\_core/abm/include/mal\_abm\_fast/config.hpp, mal-core/src/mal\_core/abm/include/mal\_abm\_fast/host\_cascade.hpp, mal-core/src/mal\_core/abm/src/host\_cascade.cpp \\
\code{species.stephensi.wind\_reference} & "u100/v100" & string & (loaded, never used) \\
\code{thermal.efd\_max\_eggs\_per\_day} & 0.5 & huevos·hembra⁻¹·día⁻¹ & mal-core/src/mal\_core/abm/include/mal\_abm\_fast/birth\_rate.hpp, mal-core/src/mal\_core/abm/include/mal\_abm\_fast/config.hpp \\
\code{thermal.efd\_tmax\_c} & 34.0 & °C & mal-core/src/mal\_core/abm/include/mal\_abm\_fast/birth\_rate.hpp, mal-core/src/mal\_core/abm/include/mal\_abm\_fast/config.hpp \\
\code{thermal.efd\_tmin\_c} & 16.0 & °C & mal-core/src/mal\_core/abm/include/mal\_abm\_fast/birth\_rate.hpp, mal-core/src/mal\_core/abm/include/mal\_abm\_fast/config.hpp \\
\code{thermal.efd\_topt\_c} & 25.0 & °C & mal-core/src/mal\_core/abm/include/mal\_abm\_fast/birth\_rate.hpp, mal-core/src/mal\_core/abm/include/mal\_abm\_fast/config.hpp \\
\code{thermal.egg\_development\_scale} & 0.00062 & escala & mal-core/src/mal\_core/abm/include/mal\_abm\_fast/config.hpp, mal-core/src/mal\_core/abm/include/mal\_abm\_fast/thermal\_responses.hpp \\
\code{thermal.egg\_development\_tmax\_c} & 40.0 & °C & mal-core/src/mal\_core/abm/include/mal\_abm\_fast/config.hpp, mal-core/src/mal\_core/abm/include/mal\_abm\_fast/thermal\_responses.hpp \\
\code{thermal.egg\_development\_tmin\_c} & 10.0 & °C & mal-core/src/mal\_core/abm/include/mal\_abm\_fast/config.hpp, mal-core/src/mal\_core/abm/include/mal\_abm\_fast/thermal\_responses.hpp \\
\code{thermal.larva\_development\_scale} & 0.001 & escala & mal-core/src/mal\_core/abm/include/mal\_abm\_fast/config.hpp, mal-core/src/mal\_core/abm/include/mal\_abm\_fast/thermal\_responses.hpp \\
\code{thermal.larva\_development\_tmax\_c} & 40.0 & °C & mal-core/src/mal\_core/abm/include/mal\_abm\_fast/config.hpp, mal-core/src/mal\_core/abm/include/mal\_abm\_fast/thermal\_responses.hpp \\
\code{thermal.larva\_development\_tmin\_c} & 10.0 & °C & mal-core/src/mal\_core/abm/include/mal\_abm\_fast/config.hpp, mal-core/src/mal\_core/abm/include/mal\_abm\_fast/thermal\_responses.hpp \\
\code{thermal.pupa\_development\_scale} & 0.00066 & escala & mal-core/src/mal\_core/abm/include/mal\_abm\_fast/config.hpp, mal-core/src/mal\_core/abm/include/mal\_abm\_fast/thermal\_responses.hpp \\
\code{thermal.pupa\_development\_tmax\_c} & 40.0 & °C & mal-core/src/mal\_core/abm/include/mal\_abm\_fast/config.hpp, mal-core/src/mal\_core/abm/include/mal\_abm\_fast/thermal\_responses.hpp \\
\code{thermal.pupa\_development\_tmin\_c} & 12.0 & °C & mal-core/src/mal\_core/abm/include/mal\_abm\_fast/config.hpp, mal-core/src/mal\_core/abm/include/mal\_abm\_fast/thermal\_responses.hpp \\
\code{transmission.beta\_hv} & 0.4 & prob/mordida & mal-core/src/mal\_core/abm/include/mal\_abm\_fast/config.hpp, mal-core/src/mal\_core/abm/include/mal\_abm\_fast/transmission.hpp … (+3 more) \\
\code{transmission.beta\_vh} & 0.5 & prob/mordida & mal-core/src/mal\_core/abm/include/mal\_abm\_fast/config.hpp, mal-core/src/mal\_core/abm/include/mal\_abm\_fast/transmission.hpp … (+3 more) \\
\code{transmission.human\_immunity\_days} & 180 & días & mal-core/src/mal\_core/abm/include/mal\_abm\_fast/config.hpp, mal-core/src/mal\_core/abm/include/mal\_abm\_fast/transmission.hpp, mal-core/src/mal\_core/abm/src/main.cpp \\
\code{transmission.human\_incubation\_days} & 12 & días & mal-core/src/mal\_core/abm/include/mal\_abm\_fast/config.hpp, mal-core/src/mal\_core/abm/include/mal\_abm\_fast/transmission.hpp … (+3 more) \\
\code{transmission.human\_infectious\_days} & 20 & días & mal-core/src/mal\_core/abm/include/mal\_abm\_fast/config.hpp, mal-core/src/mal\_core/abm/include/mal\_abm\_fast/transmission.hpp … (+3 more) \\
\code{transmission.initial\_human\_prevalence} & 0.05 & fracción & mal-core/src/mal\_core/abm/include/mal\_abm\_fast/config.hpp, mal-core/src/mal\_core/abm/include/mal\_abm\_fast/transmission.hpp … (+3 more) \\
\code{transmission.initial\_vector\_infected\_frac} & 0.0 & fracción & mal-core/src/mal\_core/abm/include/mal\_abm\_fast/config.hpp, mal-core/src/mal\_core/abm/include/mal\_abm\_fast/transmission.hpp … (+3 more) \\
\code{urban.beta\_urban} & 0.75 & adimensional & mal-core/src/mal\_core/abm/include/mal\_abm\_fast/config.hpp, mal-core/src/mal\_core/abm/include/mal\_abm\_fast/habitat\_availability.hpp \\
\code{urban.rain\_pools\_built\_damage} & 0.3 & adimensional & mal-core/src/mal\_core/abm/include/mal\_abm\_fast/config.hpp, mal-core/src/mal\_core/abm/include/mal\_abm\_fast/habitat\_availability.hpp \\
\bottomrule
\end{longtable}

\endgroup

% =====================================================================
\begin{thebibliography}{99}

\bibitem{Bayoh2003}
Bayoh, M.N. and Lindsay, S.W. (2003)
Effect of temperature on the development of the aquatic stages of
\emph{Anopheles gambiae} sensu stricto.
\emph{Bulletin of Entomological Research}, 93, 375--381.

\bibitem{Brown2022}
Brown, C.F., Brumby, S.P., Guzder-Williams, B., et al. (2022)
Dynamic World, near real-time global 10\,m land use land cover mapping.
\emph{Scientific Data}, 9, 251.
DOI: 10.1038/s41597-022-01307-4.

\bibitem{Funk2014}
Funk, C.C., Peterson, P.J., Landsfeld, M.F., et al. (2014)
A quasi-global precipitation time series for drought monitoring.
\emph{U.S. Geological Survey Data Series} 832, 4\,p.
Data: UCSB Climate Hazards Center, CHIRPS-2.0.

\bibitem{Grimm2006}
Grimm, V., Berger, U., Bastiansen, F., et al. (2006)
A standard protocol for describing individual-based and agent-based models.
\emph{Ecological Modelling}, 198, 115--126.
DOI: 10.1016/j.ecolmodel.2006.04.023.

\bibitem{Grimm2010}
Grimm, V., Berger, U., DeAngelis, D.L., Polhill, J.G., Giske, J. and
Railsback, S.F. (2010)
The ODD protocol: a review and first update.
\emph{Ecological Modelling}, 221(23), 2760--2768.
DOI: 10.1016/j.ecolmodel.2010.08.019.

\bibitem{GrimmRailsback2005}
Grimm, V. and Railsback, S.F. (2005)
\emph{Individual-Based Modeling and Ecology}.
Princeton University Press, Princeton, NJ.

\bibitem{Huestis2019}
Huestis, D.L., Dao, A., Diallo, M., et al. (2019)
Windborne long-distance migration of malaria mosquitoes in the Sahel.
\emph{Nature}, 574, 404--408.

\bibitem{Lloyd2019}
Lloyd, C.T., Chamberlain, H., Kerr, D., et al. (2019)
Global spatio-temporally harmonised datasets for producing high-resolution
gridded population distribution datasets.
\emph{Big Earth Data}, 3, 108--139.

\bibitem{Macdonald1957}
Macdonald, G. (1957)
\emph{The Epidemiology and Control of Malaria}.
Oxford University Press, London.

\bibitem{Mordecai2013}
Mordecai, E.A., Paaijmans, K.P., Johnson, L.R., et al. (2013)
Optimal temperature for malaria transmission is dramatically lower than
previously predicted.
\emph{Ecology Letters}, 16, 22--30.

\bibitem{Munoz2021}
Mu\~noz-Sabater, J., Dutra, E., Agust\'i-Panareda, A., et al. (2021)
ERA5-Land: a state-of-the-art global reanalysis dataset for land
applications.
\emph{Earth System Science Data}, 13, 4349--4383.

\bibitem{Pekel2016}
Pekel, J.-F., Cottam, A., Gorelick, N. and Belward, A.S. (2016)
High-resolution mapping of global surface water and its long-term changes.
\emph{Nature}, 540, 418--422.

\bibitem{Railsback2001}
Railsback, S.F. (2001)
Conceptual design and simulation analysis of individual-based models.
In: \emph{Proceedings of the Third International Symposium on Tools and
Methods of Competitive Engineering}, Delft University Press, pp.~1099--1112.

\bibitem{Schiavina2023}
Schiavina, M., Pesaresi, M. and Melchiorri, M. (2023)
GHS-SMOD R2023A: GHSL settlement model grid.
European Commission, Joint Research Centre (JRC), Ispra.

\bibitem{Smith2004}
Smith, D.L. and McKenzie, F.E. (2004)
Statics and dynamics of malaria infection in \emph{Anopheles} mosquitoes.
\emph{Malaria Journal}, 3, 13.

\bibitem{Wessel1996}
Wessel, P. and Smith, W.H.F. (1996)
A global, self-consistent, hierarchical, high-resolution shoreline database.
\emph{Journal of Geophysical Research}, 101(B4), 8741--8743.

\bibitem{Zurell2010}
Zurell, D., Jeltsch, F. and Grimm, V. (2010)
Ecological modelling and data availability: do we really need all of these?
\emph{Journal of Applied Ecology}, 47, 623--626.

\end{thebibliography}

\vspace{1em}
\noindent\footnotesize
\textbf{Provenance note.} Several comments in
\code{model_constants.yaml} cite author--year anchors without a full
bibliographic record in the repository (for example Costantini 1996 and
Thomas 2013 for the dispersal priors, Faiman 2020 and Sanogo 2021 for wind
flight, Okal 2013 for the oviposition plume, Giraldo 2023 for the host
plume, and Lardeux for the adult survival curve). They are reported in
the tables as provenance labels, not as verified references. The numerical
values remain engineering priors or literature anchors as marked; none
has been re-fitted to Ghanaian data for this release.

\end{document}
